% Fonction exponentielle népérienne — Mathématiques Tle E — Cours signé OnBuch+
% Programme officiel MINESEC. Compiler avec : tectonic M07-fonction-exponentielle-neperienne.tex
\newcommand{\DOCMATIERE}{Mathématiques}
\newcommand{\DOCNIVEAU}{Tle E}
\newcommand{\DOCMODULE}{Relations et opérations fondamentales dans l'ensemble des nombres réels}
\newcommand{\DOCLECON}{07}
\newcommand{\DOCTITRE}{Fonction exponentielle népérienne}
% OnBuch+ — préambule « cours signés » ENRICHI (compilé avec tectonic / XeLaTeX)
% Système visuel « Pop » d'OnBuch+ : fond crème (jamais blanc), bordures encre
% épaisses, ombres « dures » décalées, titres Archivo Black en capitales.
% Version MATHÉMATIQUES : graphes (pgfplots/tikz), boîtes pédagogiques
% (définition, propriété, méthode, attention, le savais-tu, exercice résolu,
% exercice, TP, bilan), page de garde et signature OnBuch+ ancrée en bas.
%
% Macros attendues dans main.tex AVANT \input{preamble} :
%   \DOCMATIERE  \DOCNIVEAU  \DOCMODULE  \DOCLECON (numéro)  \DOCTITRE
\documentclass[11pt]{article}

\usepackage{fontspec}
\usepackage[french]{babel}
\usepackage[a4paper,margin=2cm,top=3.4cm,bottom=2.8cm,headheight=1.4cm,footskip=1.3cm]{geometry}
\usepackage{amsmath,amssymb,amsfonts}
\usepackage{mathtools} % \xmapsto, \coloneqq... souvent utilisés par les modèles
\usepackage{cancel} % simplifications barrées
% \abs{z} (module, valeur absolue) et \norm{u} (norme), utilisés par les modèles
\providecommand{\abs}{}\let\abs\relax\DeclarePairedDelimiter{\abs}{\lvert}{\rvert}
\providecommand{\norm}{}\let\norm\relax\DeclarePairedDelimiter{\norm}{\lVert}{\rVert}
\usepackage[table]{xcolor} % [table] : fournit aussi \rowcolors (alternance de lignes)
\usepackage{pagecolor}
\usepackage{tikz}
\usetikzlibrary{arrows.meta,positioning,calc,shapes.geometric,shapes.misc,shapes.arrows,decorations.pathmorphing,decorations.markings,patterns,shadows,mindmap,trees,fit,backgrounds,matrix,intersections}
\usepackage{pgfplots}
\usepackage{pgf-pie} % diagrammes circulaires \pie (statistiques)
\pgfplotsset{compat=1.18}
\usepgfplotslibrary{fillbetween,groupplots} % aires entre courbes (calcul intégral)
\usepackage{enumitem}
\usepackage{fancyhdr}
% [most] charge toutes les bibliothèques tcolorbox (listings, minted, raster,
% documentation...) et gonfle inutilement la mémoire TeX sur un document long
% (~300 boîtes) : on ne charge que ce qui est réellement utilisé.
\usepackage[skins,breakable]{tcolorbox}
\usepackage{graphicx}
\usepackage{colortbl}
\usepackage{booktabs}
\usepackage{tabularx}
\usepackage{diagbox} % case d'en-tête diagonale des tableaux à double entrée
% Colonnes extensibles centrée / gauche / droite (C, L, R), souvent utilisées par les modèles
\newcolumntype{C}{>{\centering\arraybackslash}X}
\newcolumntype{L}{>{\raggedright\arraybackslash}X}
\newcolumntype{R}{>{\raggedleft\arraybackslash}X}
\usepackage{array}
\usepackage{multicol}
\usepackage{siunitx}
% parse-numbers=false : accepte aussi \SI{2,5 \times 10^{-3}}{...}
\sisetup{locale=FR,output-decimal-marker={,},group-minimum-digits=4,parse-numbers=false}
\DeclareSIUnit\ohm{\ensuremath{\symup{\Omega}}} % U+2126 absent de la police
\usepackage{qrcode}
% \degree : commande fréquemment utilisée par les modèles pour un angle ou une
% température en dehors de \SI{}{} (non fournie par les paquets chargés).
\providecommand{\degree}{^{\circ}}
% \qed (fin de démonstration), utilisé par les modèles sans amsthm
\providecommand{\qed}{\hfill$\square$}
% \barre : double barre des valeurs interdites dans un tableau de variations (tkz-tab)
\providecommand{\barre}{\ensuremath{\|}}
% environnement proof (amsthm non chargé) utilisé par les modèles
\ifdefined\proof\else\newenvironment{proof}[1][Démonstration]{\par\noindent\textit{#1.}\ }{\qed\par}\fi
\usepackage{lastpage}
\usepackage{hyperref}
\hypersetup{hidelinks}

% ============================================================
% Palette « Pop » OnBuch+ — mêmes hex que la classe P (plus_ui.dart)
% ============================================================
\definecolor{poporange}{HTML}{F59321}
\definecolor{popdark}{HTML}{E07A0C}
\definecolor{popcream}{HTML}{FFF6E8}
\definecolor{popink}{HTML}{1C1714}
\definecolor{poppurple}{HTML}{7A5AE0}
\definecolor{popgreen}{HTML}{2E9E6A}
\definecolor{poppink}{HTML}{E0457B}
\definecolor{popblue}{HTML}{2F7DE1}
\definecolor{popgold}{HTML}{C98A12}
\definecolor{popmuted}{HTML}{8A7F76}
\definecolor{popline}{HTML}{EADFD2}
\definecolor{popgray}{HTML}{8A7F76}
\colorlet{popgrey}{popgray}
% Alias tolérants (le modèle nomme parfois les couleurs différemment)
\colorlet{popwhite}{white}
\colorlet{popblack}{popink}
\colorlet{popred}{poppink}
\colorlet{popyellow}{popgold}
% Teintes claires pour fonds de tableaux / zones
\colorlet{popblueL}{popblue!10!white}
\colorlet{poporangeL}{poporange!14!white}
\colorlet{popgreenL}{popgreen!12!white}
\colorlet{poppurpleL}{poppurple!12!white}
\colorlet{poppinkL}{poppink!12!white}
\colorlet{popgoldL}{popgold!14!white}

\pagecolor{popcream}

% ============================================================
% Typographie
% ============================================================
\defaultfontfeatures{Ligatures=TeX}
\setmainfont{PlusJakartaSans-Medium.ttf}[Path=../../../fonts/,BoldFont=PlusJakartaSans-Bold.ttf]
% Maths en sans-serif (Fira Math) avec chiffres et lettres droites de Plus
% Jakarta, pour que les formules se fondent dans le texte.
\usepackage[math-style=ISO,bold-style=ISO]{unicode-math}
\setmathfont{FiraMath-Regular.otf}[Path=../../../fonts/]
\setmathfont{PlusJakartaSans-Medium.ttf}[Path=../../../fonts/,range={up/num,up/{Latin,latin},\mathup}]
\usepackage{icomma} % virgule décimale sans espace en mode math
% Caractères mathématiques Unicode absents des polices de texte (Plus Jakarta,
% Archivo Black) : rendus par leur équivalent mathématique, en texte comme en
% formule — sinon ils s'affichent en carré vide (ex. « Arithmétique dans ℤ »).
\usepackage{newunicodechar}
\newunicodechar{ℝ}{\ensuremath{\mathbb{R}}}
\newunicodechar{ℤ}{\ensuremath{\mathbb{Z}}}
\newunicodechar{ℂ}{\ensuremath{\mathbb{C}}}
\newunicodechar{ℕ}{\ensuremath{\mathbb{N}}}
\newunicodechar{ℚ}{\ensuremath{\mathbb{Q}}}
\newunicodechar{𝔻}{\ensuremath{\mathbb{D}}}
\newunicodechar{α}{\ensuremath{\alpha}}
\newunicodechar{β}{\ensuremath{\beta}}
\newunicodechar{γ}{\ensuremath{\gamma}}
\newunicodechar{δ}{\ensuremath{\delta}}
\newunicodechar{ε}{\ensuremath{\varepsilon}}
\newunicodechar{θ}{\ensuremath{\theta}}
\newunicodechar{λ}{\ensuremath{\lambda}}
\newunicodechar{μ}{\ensuremath{\mu}}
\newunicodechar{σ}{\ensuremath{\sigma}}
\newunicodechar{τ}{\ensuremath{\tau}}
\newunicodechar{φ}{\ensuremath{\varphi}}
\newunicodechar{ψ}{\ensuremath{\psi}}
\newunicodechar{ω}{\ensuremath{\omega}}
\newunicodechar{Σ}{\ensuremath{\Sigma}}
% majuscules produites par \MakeUppercase dans les titres (α-aminés -> Α)
\newunicodechar{Α}{\ensuremath{\alpha}}
\newunicodechar{Β}{\ensuremath{\beta}}
\newunicodechar{Γ}{\ensuremath{\Gamma}}
\newunicodechar{Δ}{\ensuremath{\Delta}}
\newunicodechar{Ε}{\ensuremath{\varepsilon}}
\newunicodechar{Θ}{\ensuremath{\Theta}}
\newunicodechar{Λ}{\ensuremath{\Lambda}}
\newunicodechar{Μ}{\ensuremath{\mu}}
\newunicodechar{Τ}{\ensuremath{\tau}}
\newunicodechar{Φ}{\ensuremath{\Phi}}
\newunicodechar{Ψ}{\ensuremath{\Psi}}
\newunicodechar{Ω}{\ensuremath{\Omega}}
\newunicodechar{↦}{\ensuremath{\mapsto}}
\newunicodechar{∈}{\ensuremath{\in}}
\newunicodechar{∉}{\ensuremath{\notin}}
\newunicodechar{∧}{\ensuremath{\wedge}}
\newunicodechar{⇔}{\ensuremath{\Leftrightarrow}}
\newunicodechar{⟺}{\ensuremath{\Longleftrightarrow}}
\newunicodechar{⇒}{\ensuremath{\Rightarrow}}
\newunicodechar{⟹}{\ensuremath{\Longrightarrow}}
\newunicodechar{∘}{\ensuremath{\circ}}
\newunicodechar{∪}{\ensuremath{\cup}}
\newunicodechar{∩}{\ensuremath{\cap}}
\newunicodechar{≡}{\ensuremath{\equiv}}
\newunicodechar{⊂}{\ensuremath{\subset}}
\newunicodechar{⊥}{\ensuremath{\perp}}
\newunicodechar{∃}{\ensuremath{\exists}}
\newunicodechar{∀}{\ensuremath{\forall}}
\newunicodechar{∛}{\ensuremath{\sqrt[3]{\,}}}
\newunicodechar{∜}{\ensuremath{\sqrt[4]{\,}}}
\newunicodechar{⃗}{\ensuremath{{}^{\to}}}
\newunicodechar{ⁿ}{\ifmmode^{n}\else\textsuperscript{\ensuremath{n}}\fi}
\newunicodechar{ˣ}{\ifmmode^{x}\else\textsuperscript{\ensuremath{x}}\fi}
\newunicodechar{ᵇ}{\ifmmode^{b}\else\textsuperscript{\ensuremath{b}}\fi}
\newunicodechar{ʳ}{\ifmmode^{r}\else\textsuperscript{\ensuremath{r}}\fi}
\newunicodechar{ᵘ}{\ifmmode^{u}\else\textsuperscript{\ensuremath{u}}\fi}
\newunicodechar{ʸ}{\ifmmode^{y}\else\textsuperscript{\ensuremath{y}}\fi}
\newunicodechar{ᵃ}{\ifmmode^{a}\else\textsuperscript{\ensuremath{a}}\fi}
\newunicodechar{ᵉ}{\ifmmode^{e}\else\textsuperscript{\ensuremath{e}}\fi}
\newunicodechar{ᵏ}{\ifmmode^{k}\else\textsuperscript{\ensuremath{k}}\fi}
\newunicodechar{ᵖ}{\ifmmode^{p}\else\textsuperscript{\ensuremath{p}}\fi}
\newunicodechar{ᴺ}{\ifmmode^{N}\else\textsuperscript{\ensuremath{N}}\fi}
\newunicodechar{ᶜ}{\ifmmode^{c}\else\textsuperscript{\ensuremath{c}}\fi}
\newunicodechar{ʰ}{\ifmmode^{h}\else\textsuperscript{\ensuremath{h}}\fi}
\newunicodechar{ᵠ}{\ifmmode^{q}\else\textsuperscript{\ensuremath{q}}\fi}
\newunicodechar{ₙ}{\ifmmode_{n}\else\textsubscript{\ensuremath{n}}\fi}
\newunicodechar{ₐ}{\ifmmode_{a}\else\textsubscript{\ensuremath{a}}\fi}
\newunicodechar{ᵢ}{\ifmmode_{i}\else\textsubscript{\ensuremath{i}}\fi}
\newunicodechar{ᵣ}{\ifmmode_{r}\else\textsubscript{\ensuremath{r}}\fi}
\newunicodechar{ₖ}{\ifmmode_{k}\else\textsubscript{\ensuremath{k}}\fi}
\newunicodechar{ᵦ}{\ifmmode_{\beta}\else\textsubscript{\ensuremath{\beta}}\fi}
\newunicodechar{ₓ}{\ifmmode_{x}\else\textsubscript{\ensuremath{x}}\fi}
\newfontfamily\popdisplay{ArchivoBlack-Regular.ttf}[Path=../../../fonts/]
\newfontfamily\poplogo{SpaceGrotesk-Bold.ttf}[Path=../../../fonts/]
\newfontfamily\popsemi{PlusJakartaSans-SemiBold.ttf}[Path=../../../fonts/]
\newfontfamily\popxbold{PlusJakartaSans-ExtraBold.ttf}[Path=../../../fonts/]
\newfontfamily\popxboldls{PlusJakartaSans-ExtraBold.ttf}[Path=../../../fonts/,LetterSpace=3.0]

\setlist[enumerate]{leftmargin=1.7em,itemsep=5pt,topsep=4pt}
\setlist[itemize]{leftmargin=1.7em,itemsep=4pt,topsep=4pt,label={\color{poporange}\textbullet}}
\setlength{\parindent}{0pt}
\setlength{\parskip}{5pt}
\renewcommand{\arraystretch}{1.25}

% pgfplots : style Pop par défaut
\pgfplotsset{
  every axis/.append style={axis line style={popink,line width=1pt},tick style={popink},
    grid style={popline,line width=0.5pt},label style={font=\small},tick label style={font=\footnotesize}},
  popcurve/.style={poporange,line width=1.8pt,no markers,smooth},
}
% Même style disponible dans un \draw TikZ simple (hors pgfplots)
\tikzset{popcurve/.style={poporange,line width=1.8pt}}

% ============================================================
% Pill / squiggle
% ============================================================
\newcommand{\poppill}[3]{%
  \tikz[baseline=-0.55ex]\node[draw=popink,line width=1.2pt,rounded corners=2.6pt,%
    fill=#2,inner xsep=6.5pt,inner ysep=3pt]{%
    \popxboldls\fontsize{7}{7}\selectfont\color{#3}\MakeUppercase{#1}};%
}
\newcommand{\popsquiggle}[1]{%
  \tikz[baseline=-0.4ex]{
    \draw[#1,line width=2.6pt,line cap=round]
      plot [smooth] coordinates {(0,0.09) (0.34,0.01) (0.68,0.21) (1.02,0.01) (1.36,0.10)};
  }%
}
% Mot-clé surligné (marqueur orange sous le texte)
\newcommand{\cle}[1]{{\popxbold\color{popdark}#1}}

% ============================================================
% En-tête / pied
% ============================================================
\pagestyle{fancy}
\fancyhf{}
\renewcommand{\headrulewidth}{0pt}
\renewcommand{\footrulewidth}{0pt}
\lhead{\raisebox{-0.15ex}{\poplogo\fontsize{13}{13}\selectfont\color{popink}OnBuch\color{poporange}+}}
\rhead{\poppill{\DOCMATIERE\ \textbullet\ \DOCNIVEAU\ \textbullet\ Leçon \DOCLECON}{white}{popink}}
\lfoot{\popxbold\fontsize{7.6}{7.6}\selectfont\color{popmuted}\MakeUppercase{Cours signé OnBuch+}\ \textbullet\ {\popsemi\DOCTITRE}}
\rfoot{\tikz[baseline=-0.6ex]\node[rounded corners=8.5pt,fill=popink,minimum height=17pt,minimum width=17pt,inner xsep=5pt,inner sep=0]{\popxbold\fontsize{8}{8}\selectfont\color{popcream}\thepage\,/\,\pageref*{LastPage}};}

% ============================================================
% Titres
% ============================================================
\newcommand{\chaptitre}[2]{%
  \begin{tcolorbox}[enhanced,colback=poporange,colframe=popink,boxrule=2.6pt,arc=11pt,
      drop shadow=popink,left=15pt,right=15pt,top=12pt,bottom=11pt,before skip=3pt,after skip=18pt]
    \poppill{#2}{popink}{popcream}\par\vspace{9pt}
    {\raggedright\hyphenpenalty=10000\exhyphenpenalty=10000\popdisplay\fontsize{22}{24}\selectfont\color{popcream}\MakeUppercase{#1}\par}
  \end{tcolorbox}
}
% Section numérotée : pastille orange avec le numéro + titre Archivo
\newcounter{popsec}
\newcommand{\coursec}[1]{%
  \stepcounter{popsec}\setcounter{popsub}{0}%
  \par\vspace{16pt}\noindent
  \begin{minipage}{\linewidth}
  \tikz[baseline=-0.7ex]\node[draw=popink,line width=1.6pt,fill=poporange,rounded corners=4pt,minimum size=22pt,inner sep=0]{\popdisplay\fontsize{12}{12}\selectfont\color{popcream}\thepopsec};%
  \hspace{8pt}{\popdisplay\fontsize{15}{17}\selectfont\color{popink}\MakeUppercase{#1}}\par
  \vspace{4pt}\noindent{\color{poporange}\rule{36pt}{3.6pt}}
  \end{minipage}\par\nopagebreak\vspace{8pt}
}
\newcounter{popsub}
\newcommand{\courssub}[1]{%
  \stepcounter{popsub}%
  \par\vspace{9pt}\noindent{\popxbold\fontsize{11.5}{13}\selectfont\color{popink}{\color{poporange}\thepopsec.\thepopsub}\hspace{6pt}#1}\par\nopagebreak\vspace{4pt}
}

% ============================================================
% Boîtes pédagogiques — même recette : fond blanc, bordure épaisse colorée,
% ombre dure encre, badge plein. Titre optionnel [#1].
% ============================================================
\tcbset{popbase/.style={breakable,enhanced,colback=white,boxrule=2.3pt,arc=9pt,
  shadow={3pt}{-3pt}{0pt}{popink},left=14pt,right=14pt,top=11pt,bottom=11pt,before skip=11pt,after skip=15pt,
  fontupper=\popsemi\fontsize{10.6}{14.5}\selectfont\color{popink},
  fontlower=\popsemi\fontsize{10.6}{14.5}\selectfont\color{popink}}}
% Coche dessinée (le glyphe ✓ n'existe pas dans les polices du document)
\newcommand{\popcheck}[1]{\tikz[baseline=-0.5ex]\draw[#1,line width=1.6pt,line cap=round,line join=round](-0.13,0.0)--(-0.04,-0.09)--(0.13,0.1);}
\newcommand{\popboxhead}[3]{\poppill{#1}{#2}{white}\ifx\relax#3\relax\else\hspace{6pt}{\popxbold\fontsize{10.6}{12}\selectfont\color{popink}#3}\fi\par\vspace{7pt}}

\newtcolorbox{aretenir}[1][]{popbase,colframe=popgreen,before upper={\popboxhead{À retenir}{popgreen}{#1}}}
\newtcolorbox{exemplebox}[1][]{popbase,colframe=poporange,before upper={\popboxhead{Exemple}{poporange}{#1}}}
\newtcolorbox{definition}[1][]{popbase,colframe=popblue,before upper={\popboxhead{Définition}{popblue}{#1}}}
\newtcolorbox{propriete}[1][]{popbase,colframe=poppurple,before upper={\popboxhead{Propriété}{poppurple}{#1}}}
\newtcolorbox{methode}[1][]{popbase,colframe=popgold,before upper={\popboxhead{Méthode}{popgold}{#1}}}
\newtcolorbox{attention}[1][]{popbase,colframe=poppink,before upper={\popboxhead{Attention}{poppink}{#1}}}
\newtcolorbox{experience}[1][]{popbase,colframe=popblue,colback=popblueL,before upper={\popboxhead{Expérience}{popblue}{#1}}}
% « Le savais-tu ? » : carte violette pleine (contexte, culture, Cameroun)
\newtcolorbox{savaistu}[1][]{popbase,colback=poppurple,colframe=popink,
  fontupper=\popsemi\fontsize{10.4}{14.2}\selectfont\color{white},
  before upper={\poppill{Le savais-tu\,?}{white}{poppurple}\ifx\relax#1\relax\else\hspace{6pt}{\popxbold\color{white}#1}\fi\par\vspace{7pt}}}
% Exercice résolu : énoncé en haut, \tcblower puis la solution sur fond crème
\newcounter{popexo}\newcounter{popexr}
\newtcolorbox{exoresolu}[1][]{popbase,colframe=poporange,colbacklower=popcream,
  segmentation style={popink,line width=1.2pt,dashed},
  before upper={\stepcounter{popexr}\popboxhead{Exercice résolu \thepopexr}{poporange}{#1}},
  before lower={\poppill{Solution}{popink}{popcream}\par\vspace{6pt}}}
% Exercice (non résolu en place) — la correction est dans la partie Corrigés
\newtcolorbox{exercice}[1][]{popbase,colframe=popink,
  before upper={\stepcounter{popexo}\popboxhead{Exercice \thepopexo}{popink}{#1}}}
% Corrigé
\newtcolorbox{corrige}[1][]{popbase,colframe=popgreen,colback=popgreenL,
  before upper={\popboxhead{Corrigé}{popgreen}{#1}}}
% Objectifs / prérequis / bilan
\newtcolorbox{objectifs}{popbase,colframe=popink,colback=poporangeL,
  before upper={\popboxhead{Objectifs de la leçon}{popink}{}}}
\newtcolorbox{prerequis}{popbase,colframe=popink,colback=popblueL,
  before upper={\popboxhead{Prérequis}{popblue}{}}}
\newtcolorbox{bilan}{popbase,colframe=popink,colback=white,boxrule=2.8pt,
  before upper={\popboxhead{Fiche bilan}{poporange}{L'essentiel à connaître pour l'examen}}}
% Figure encadrée avec légende
\newtcolorbox{popfigure}[1][]{enhanced,breakable=false,colback=white,colframe=popline,boxrule=1.4pt,arc=8pt,
  left=10pt,right=10pt,top=10pt,bottom=8pt,before skip=10pt,after skip=12pt,halign=center,
  lower separated=false,halign lower=center,
  fontlower=\popsemi\fontsize{9}{11.5}\selectfont\color{popmuted},
  #1}
\newcommand{\legende}[1]{\par\vspace{6pt}{\popsemi\fontsize{9}{11.5}\selectfont\color{popmuted}#1}}

% ============================================================
% Signature OnBuch+
% ============================================================
\newcommand{\poplogoinline}[1]{{\poplogo\fontsize{#1}{#1}\selectfont\color{popink}OnBuch\color{poporange}+}}
\newcommand{\signatureonbuch}{%
  \begin{tcolorbox}[enhanced,colback=popink,colframe=popink,boxrule=2.2pt,arc=10pt,
      drop shadow=poporange,left=14pt,right=14pt,top=10pt,bottom=10pt,before skip=0pt,after skip=0pt]
    \tikz[baseline=-0.7ex]\node[circle,fill=popgreen,draw=popcream,line width=1.3pt,minimum size=22pt,inner sep=0]{\popcheck{white}};%
    \hspace{9pt}\begin{minipage}[c]{0.62\linewidth}
      {\popxbold\fontsize{12}{13}\selectfont\color{popcream}Signé\ \textbullet\ {\poplogo OnBuch\color{poporange}+}}\par\vspace{2pt}
      {\raggedright\popsemi\fontsize{8}{10.5}\selectfont\color{popline}\MakeUppercase{Équipe pédagogique OnBuch+}\\\MakeUppercase{Conforme au programme officiel MINESEC}\par}
    \end{minipage}\hfill
    {\poplogo\fontsize{22}{22}\selectfont\color{popcream}OnBuch\color{poporange}+}
  \end{tcolorbox}%
}

% ============================================================
% Page de garde
% #1 titre · #2 matière · #3 niveau · #4 module · #5 n° leçon · #6 durée ·
% #7 description / objectifs (texte ou itemize)
% ============================================================
\newcommand{\pagegarde}[7]{%
  \thispagestyle{empty}
  \newgeometry{a4paper,margin=1.7cm,top=1.6cm,bottom=1.4cm}%
  \begin{tikzpicture}[remember picture,overlay]
    \fill[poporange!18] ([xshift=-3.2cm,yshift=-3.6cm]current page.north east) circle (4.6cm);
    \fill[poppurple!14] ([xshift=2.4cm,yshift=7.6cm]current page.south west) circle (3.8cm);
  \end{tikzpicture}%
  {\poplogo\fontsize{30}{30}\selectfont\color{popink}OnBuch\color{poporange}+}\hfill
  \raisebox{0.6ex}{\poppill{Cours signé \textbullet\ Premium}{popink}{popcream}}\par
  \vspace{1pt}\popsquiggle{poppurple}\par
  \vspace{8pt}
  \begin{tcolorbox}[enhanced,colback=poporange,colframe=popink,boxrule=2.8pt,arc=13pt,
      drop shadow=popink,left=18pt,right=18pt,top=12pt,bottom=13pt,after skip=10pt]
    \poppill{#2 \textbullet\ #3}{popink}{popcream}\hspace{5pt}\poppill{#4}{white}{popink}\par\vspace{8pt}
    {\popdisplay\fontsize{14}{14}\selectfont\color{popink}LEÇON #5}\par\vspace{4pt}
    {\raggedright\hyphenpenalty=10000\exhyphenpenalty=10000\popdisplay\fontsize{25}{27}\selectfont\color{popcream}\MakeUppercase{#1}\par}
    \vspace{8pt}
    {\popxbold\fontsize{9.5}{9.5}\selectfont\color{popink}\MakeUppercase{Durée conseillée : #6}}
  \end{tcolorbox}
  \begin{tcolorbox}[enhanced,colback=white,colframe=popink,boxrule=2.2pt,arc=10pt,
      drop shadow=popink,left=16pt,right=16pt,top=10pt,bottom=10pt,after skip=8pt]
    \poppill{Dans cette leçon}{poppurple}{white}\par\vspace{5pt}
    \popsemi\fontsize{9.3}{12.6}\selectfont\color{popink}#7
  \end{tcolorbox}
  \noindent\begin{minipage}[t]{0.487\linewidth}
    \begin{tcolorbox}[enhanced,colback=popgreen,colframe=popink,boxrule=2.2pt,arc=10pt,
        drop shadow=popink,left=11pt,right=11pt,top=10pt,bottom=10pt,height=3.1cm,valign=center]
      \tcbox[colback=white,colframe=popink,boxrule=1.3pt,arc=5pt,boxsep=0pt,nobeforeafter,
        left=3pt,right=3pt,top=3pt,bottom=3pt,valign=center]{\qrcode[height=1.9cm]{https://play.google.com/store/apps/details?id=com.luvvix.onbuch&hl=fr}}%
      \hspace{8pt}\begin{minipage}[c]{0.5\linewidth}
        {\popdisplay\fontsize{11}{12.5}\selectfont\color{white}\MakeUppercase{Télécharge OnBuch}}\par\vspace{4pt}
        {\raggedright\popsemi\fontsize{8.4}{11}\selectfont\color{white}Gratuit \textbullet\ Google Play\\Tuteur IA, annales, concours, résultats\par}
      \end{minipage}
    \end{tcolorbox}
  \end{minipage}\hfill
  \begin{minipage}[t]{0.487\linewidth}
    \begin{tcolorbox}[enhanced,colback=poppurple,colframe=popink,boxrule=2.2pt,arc=10pt,
        drop shadow=popink,left=11pt,right=11pt,top=10pt,bottom=10pt,height=3.1cm,valign=center]
      \tikz[baseline=-0.7ex]\node[circle,fill=white,draw=popink,line width=1.3pt,minimum size=1.6cm,inner sep=0]{\popdisplay\fontsize{24}{24}\selectfont\color{poppurple}+};%
      \hspace{8pt}\begin{minipage}[c]{0.6\linewidth}
        {\popdisplay\fontsize{11}{12.5}\selectfont\color{white}\MakeUppercase{Passe à OnBuch+}}\par\vspace{4pt}
        {\raggedright\popsemi\fontsize{8.4}{11}\selectfont\color{white}Tous les cours signés, Tuteur IA illimité, fiches PDF — onglet OnBuch+ de l'app\par}
      \end{minipage}
    \end{tcolorbox}
  \end{minipage}
  \vspace{8pt}
  \popcontenu
  \vfill
  \signatureonbuch
  \restoregeometry
  \newpage
}

% Rangée « Ce cours contient » (page de garde)
\newcommand{\poptile}[3]{% #1 couleur #2 gros texte #3 légende
  \begin{tcolorbox}[enhanced,colback=#1,colframe=popink,boxrule=2pt,arc=9pt,shadow={2.5pt}{-2.5pt}{0pt}{popink},
      left=6pt,right=6pt,top=9pt,bottom=9pt,halign=center,valign=center,height=1.75cm,width=\linewidth,nobeforeafter]
    {\popdisplay\fontsize{12.5}{13}\selectfont\color{white}\MakeUppercase{#2}}\par\vspace{3pt}
    {\centering\popsemi\fontsize{7.8}{9.5}\selectfont\color{white}#3\par}
  \end{tcolorbox}}
\newcommand{\popcontenu}{%
  \par\noindent{\popxboldls\fontsize{7.5}{7.5}\selectfont\color{popmuted}CE COURS SIGNÉ CONTIENT}\par\vspace{7pt}
  \noindent\begin{minipage}[t]{0.235\linewidth}\poptile{popblue}{Cours}{complet, illustré, pas à pas}\end{minipage}\hfill
  \begin{minipage}[t]{0.235\linewidth}\poptile{poporange}{Méthodes}{et exercices résolus}\end{minipage}\hfill
  \begin{minipage}[t]{0.235\linewidth}\poptile{poppink}{Exercices}{type examen + corrigés}\end{minipage}\hfill
  \begin{minipage}[t]{0.235\linewidth}\poptile{popgreen}{Fiche bilan}{l'essentiel à retenir}\end{minipage}\par}

% Sommaire
\newtcolorbox{sommaire}{enhanced,colback=white,colframe=popink,boxrule=2.2pt,arc=10pt,
  drop shadow=popink,left=18pt,right=18pt,top=13pt,bottom=13pt,before skip=6pt,after skip=20pt,
  before upper={\poppill{Sommaire}{poppurple}{white}\par\vspace{9pt}\popsemi\fontsize{10.6}{14.5}\selectfont\color{popink}}}

% Signature de fin de cours — poussée tout en bas de la dernière page
\newcommand{\findecours}{%
  \par\vfill\nopagebreak
  \begin{center}\popsquiggle{poporange}\end{center}\vspace{4pt}
  \signatureonbuch
}
\AtBeginDocument{\def\llbracket{\mathopen{[\mkern-3.5mu[}}\def\rrbracket{\mathclose{]\mkern-3.5mu]}}} % intervalles d'entiers (glyphe ⟦ absent de la police)
% Glyphes absents de Fira Math : redéfinis à partir de glyphes disponibles
\AtBeginDocument{%
  \def\setminus{\mathbin{\backslash}}\let\smallsetminus\setminus
  \def\vdots{\mathord{\vbox{\baselineskip3.2pt\lineskiplimit0pt\kern4pt\hbox{.}\hbox{.}\hbox{.}}}}%
  \def\ddots{\mathinner{\mkern1mu\raise6.5pt\hbox{.}\mkern2mu\raise3.6pt\hbox{.}\mkern2mu\raise0.7pt\hbox{.}\mkern1mu}}%
  \def\triangle{\mathord{\scalebox{0.9}{$\symup{\Delta}$}}}%
  \def\nabla{\mathord{\raisebox{\depth}{\scalebox{1}[-1]{$\symup{\Delta}$}}}}%
  \def\checkmark{\popcheck{popdark}}%
  \def\star{\mathbin{\ast}}\def\bigstar{\mathord{\ast}}%
  \def\diamond{\mathbin{\scalebox{0.6}{$\Diamond$}}}%
}

%% FIGFIX-BEGIN v4 — correctifs globaux des figures (docs/onbuch-generation/tools/figfix)
\usepackage{adjustbox}\usepackage{etoolbox}
% toute figure TikZ/pgfplots plus large que la ligne est réduite à la largeur disponible
\BeforeBeginEnvironment{tikzpicture}{\begin{adjustbox}{max width=\linewidth}}
\AfterEndEnvironment{tikzpicture}{\end{adjustbox}}
% légendes pgfplots lisibles par-dessus les courbes
\pgfplotsset{every axis/.append style={legend style={fill=white,fill opacity=0.92,text opacity=1,draw=black!15,inner sep=3pt}}}
% étiquettes d'arêtes (syntaxe edge["..."]) avec fond blanc
\tikzset{every edge quotes/.append style={fill=white,inner sep=1.2pt}}
%% FIGFIX-END
\begin{document}

\pagegarde{Fonction exponentielle népérienne}{Mathématiques}{Tle E}{Relations et opérations fondamentales dans l'ensemble des nombres réels}{07}{8 h}{Cette leçon introduit la fonction exponentielle népérienne comme bijection réciproque du logarithme népérien. Elle en établit les propriétés algébriques, les limites, la dérivation et l'intégration, puis montre comment elle permet de résoudre des équations et de modéliser des phénomènes de croissance et de décroissance.
\vspace{4pt}
\begin{multicols}{2}\small
\begin{itemize}
\item À la fin de cette leçon, je sais définir la fonction exponentielle népérienne comme réciproque de ln et utiliser les relations e\^{}(ln x) = x et ln(e\^{}x) = x
\item À la fin de cette leçon, je sais utiliser les propriétés algébriques de exp pour transformer et simplifier une expression
\item À la fin de cette leçon, je sais résoudre une équation, une inéquation ou un système dont l'inconnue intervient dans e\^{}x ou ln x
\item À la fin de cette leçon, je sais déterminer les limites de exp aux bornes de son domaine et celles de fonctions composées de la forme e\^{}(ax+b)
\item À la fin de cette leçon, je sais dériver une fonction contenant e\^{}u et en déduire son sens de variation
\item À la fin de cette leçon, je sais déterminer une primitive de la forme u′ e\^{}u
\end{itemize}\end{multicols}}

\begin{sommaire}
\begin{multicols}{2}
\begin{enumerate}
\item Définition et premières propriétés de la fonction exponentielle
\item Propriétés algébriques de l'exponentielle
\item Équations, inéquations et systèmes avec e\^{}x
\item Limites de la fonction exponentielle
\item Dérivation, variations et courbe représentative
\item Primitives et modélisation par l'exponentielle
\item Fonction exponentielle de base a et fonctions puissances
\item Activité d'intégration
\item Méthodes et exercices résolus
\item Exercices
\item Corrigés des exercices
\item Fiche bilan
\end{enumerate}
\end{multicols}
\end{sommaire}

\begin{objectifs}
\begin{itemize}
\item À la fin de cette leçon, je sais définir la fonction exponentielle népérienne comme réciproque de ln et utiliser les relations e\^{}(ln x) = x et ln(e\^{}x) = x
\item À la fin de cette leçon, je sais utiliser les propriétés algébriques de exp pour transformer et simplifier une expression
\item À la fin de cette leçon, je sais résoudre une équation, une inéquation ou un système dont l'inconnue intervient dans e\^{}x ou ln x
\item À la fin de cette leçon, je sais déterminer les limites de exp aux bornes de son domaine et celles de fonctions composées de la forme e\^{}(ax+b)
\item À la fin de cette leçon, je sais dériver une fonction contenant e\^{}u et en déduire son sens de variation
\item À la fin de cette leçon, je sais déterminer une primitive de la forme u′ e\^{}u
\item À la fin de cette leçon, je sais étudier complètement une fonction contenant une exponentielle et tracer sa courbe représentative
\item À la fin de cette leçon, je sais modéliser une situation concrète (croissance, décroissance, refroidissement) par une fonction exponentielle et exploiter le modèle
\end{itemize}
\end{objectifs}

\begin{prerequis}
\begin{itemize}
\item Fonction logarithme népérien : définition, propriétés algébriques, dérivée, variations (leçon M06)
\item Notion de bijection et de fonction réciproque (courbes symétriques par rapport à la droite y = x)
\item Dérivation : dérivées usuelles, dérivée d'une composée u' × v'(u), étude du signe de la dérivée
\item Calcul de limites : limites usuelles, opérations sur les limites, limites de fonctions composées
\item Notion de primitive et primitives usuelles vues en début d'année de Terminale
\end{itemize}
\end{prerequis}

\newpage

\coursec{Définition et premières propriétés de la fonction exponentielle}

Au Cameroun, la population de Douala croît chaque année d'environ $3,5\,\%$ : les démographes disent qu'elle suit une \emph{croissance exponentielle}. De même, un technicien du centre hospitalier de Yaoundé suit la décroissance d'un produit radioactif utilisé en imagerie médicale, et un météorologue décrit le refroidissement d'un corps laissé à l'air libre. Toutes ces situations ont un point commun : elles se modélisent par une fonction remarquable, la \cle{fonction exponentielle népérienne}, qui « défait » exactement ce que fait le logarithme népérien.

\textbf{Problématique.} Comment définir cette fonction, quelles sont ses propriétés fondamentales, et comment l'utiliser pour résoudre des équations et étudier des phénomènes d'évolution ?

\courssub{Rappel : la fonction logarithme népérien est une bijection}

Avant de « défaire » ce que fait $\ln$, rappelons ce que nous savons de cette fonction étudiée dans la leçon précédente.

\begin{propriete}[Bijectivité de la fonction ln]
La fonction logarithme népérien $\ln$ est \textbf{continue} et \textbf{strictement croissante} sur $]0\,;\,+\infty[$, avec
\[ \lim_{x \to 0^+} \ln x = -\infty \qquad \text{et} \qquad \lim_{x \to +\infty} \ln x = +\infty. \]
Par conséquent, d'après le théorème des valeurs intermédiaires (version bijection), $\ln$ réalise une \cle{bijection} de $]0\,;\,+\infty[$ sur $\mathbb{R}$.
\end{propriete}

Concrètement, cela signifie deux choses :
\begin{itemize}
\item \textbf{Existence} : pour tout réel $y$, l'équation $\ln x = y$ admet \emph{au moins} une solution dans $]0\,;\,+\infty[$ ;
\item \textbf{Unicité} : cette solution est \emph{unique}, car $\ln$ est strictement monotone (elle ne prend jamais deux fois la même valeur).
\end{itemize}

\begin{exemplebox}[Une solution unique]
L'équation $\ln x = 0$ admet une unique solution : $x = 1$. L'équation $\ln x = 1$ admet aussi une unique solution, que nous ne savons pas encore écrire autrement que « le nombre dont le logarithme vaut $1$ ». C'est précisément ce nombre que nous allons baptiser $e$.
\end{exemplebox}

Toute bijection admet une \cle{bijection réciproque} : la fonction qui « remonte les flèches ». C'est exactement ainsi que naît l'exponentielle.

\courssub{Définition de la fonction exponentielle népérienne}

\textbf{Intuition.} La fonction $\ln$ transforme un nombre strictement positif $x$ en un réel $y$. La fonction réciproque fait le voyage en sens inverse : à un réel $y$, elle associe l'\emph{unique} nombre strictement positif $x$ tel que $\ln x = y$. Cette « machine à remonter le logarithme » est la fonction exponentielle.

\begin{definition}[Fonction exponentielle népérienne]
La fonction $\ln$ étant une bijection de $]0\,;\,+\infty[$ sur $\mathbb{R}$, elle admet une bijection réciproque, appelée \cle{fonction exponentielle népérienne} et notée $\exp$ :
\[ \exp : \mathbb{R} \longrightarrow \, ]0\,;\,+\infty[. \]
Pour tout réel $x$, on note $\exp(x) = e^x$ (se lire « exponentielle de $x$ » ou « $e$ puissance $x$ »).
\end{definition}

\begin{popfigure}
\begin{center}
\begin{tikzpicture}[scale=1.1, >=Stealth]
% Ensemble de départ R
\draw[thick, popblue] (0,0) ellipse (1.6 and 2.2);
\node[popblue, font=\bfseries] at (0,2.7) {$\mathbb{R}$};
\node at (-0.5,0.6) {$-2$};
\node at (-0.5,0) {$0$};
\node at (-0.5,-0.6) {$3$};
% Ensemble d'arrivée
\draw[thick, poporange] (6,0) ellipse (1.8 and 2.2);
\node[poporange, font=\bfseries] at (6,2.7) {$]0\,;\,+\infty[$};
\node at (6.6,0.6) {$e^{-2}$};
\node at (6.6,0) {$1$};
\node at (6.6,-0.6) {$e^{3}$};
% Flèches exp
\draw[->, thick, popgreen] (-0.1,0.6) to[bend left=15] node[above, sloped, font=\small]{$\exp$} (5.2,0.6);
\draw[->, thick, popgreen] (-0.1,0) to[bend left=15] (5.2,0);
\draw[->, thick, popgreen] (-0.1,-0.6) to[bend left=15] (5.2,-0.6);
% Flèches ln (retour)
\draw[->, thick, poppurple, dashed] (5.2,0.3) to[bend left=15] node[below, sloped, font=\small]{$\ln$} (-0.1,0.3);
\draw[->, thick, poppurple, dashed] (5.2,-0.9) to[bend left=15] (-0.1,-0.9);
\end{tikzpicture}
\end{center}
\legende{Schéma fléché des bijections réciproques : $\exp$ envoie $\mathbb{R}$ sur $]0\,;\,+\infty[$, et $\ln$ fait le chemin inverse.}
\end{popfigure}

\begin{attention}[Domaines à ne pas confondre]
La fonction $\exp$ est définie sur $\mathbb{R}$ \emph{tout entier} : on peut calculer $e^x$ pour n'importe quel réel $x$, positif, négatif ou nul. En revanche, ses \emph{valeurs} sont toujours strictement positives. C'est l'inverse de $\ln$, qui n'accepte que des entrées strictement positives mais produit n'importe quel réel.
\end{attention}

\courssub{Le nombre $e$}

\begin{definition}[Le nombre $e$]
On appelle \cle{nombre $e$} l'image de $1$ par la fonction exponentielle :
\[ e = \exp(1) = e^1. \]
C'est l'unique réel strictement positif dont le logarithme népérien vaut $1$ : $\ln e = 1$. Sa valeur approchée est
\[ e \approx 2,71828\,18284\,59\ldots \]
\end{definition}

\begin{savaistu}[Leonhard Euler et le nombre $e$]
Le nombre $e \approx 2,71828\ldots$ est, comme $\pi$, un nombre \textbf{irrationnel} : son écriture décimale est infinie et non périodique, et il ne peut pas s'écrire comme quotient de deux entiers. Il est même \emph{transcendant}, c'est-à-dire racine d'aucun polynôme à coefficients entiers. Il porte le nom du mathématicien suisse \textbf{Leonhard Euler} (1707--1783), l'un des plus prolifiques de l'histoire, qui en montra l'importance capitale en analyse. Le nombre $e$ apparaît naturellement dès qu'on étudie des phénomènes de croissance continue : intérêts composés, désintégration radioactive, propagation d'une épidémie... ou croissance de la population de Douala !
\end{savaistu}

\courssub{Relations fondamentales entre $\exp$ et $\ln$}

Puisque $\exp$ et $\ln$ sont réciproques l'une de l'autre, chacune « annule » l'effet de l'autre. C'est la propriété la plus importante de cette section.

\begin{propriete}[Relations fondamentales (conséquences directes de la définition de la réciproque)]
\begin{itemize}
\item Pour tout réel $x$ : $\ln\left(e^x\right) = x$.
\item Pour tout réel $x > 0$ : $e^{\ln x} = x$.
\end{itemize}
\end{propriete}

\begin{experience}[Pourquoi ces formules sont-elles vraies ?]
Raisonnons pas à pas sur la première relation.
\begin{enumerate}
\item Par définition de la bijection réciproque, $e^x = \exp(x)$ est \emph{l'unique} nombre strictement positif dont le logarithme vaut $x$.
\item Autrement dit, si l'on pose $y = e^x$, alors par construction $\ln y = x$.
\item En remplaçant $y$ par $e^x$, on obtient exactement $\ln(e^x) = x$.
\end{enumerate}
Le même raisonnement, en échangeant les rôles, donne $e^{\ln x} = x$ pour $x > 0$ : $e^{\ln x}$ est l'unique nombre strictement positif dont le logarithme vaut $\ln x$... et $x$ fait parfaitement l'affaire !
\end{experience}

\begin{attention}[Erreur fréquente : oublier la condition $x > 0$]
L'égalité $e^{\ln x} = x$ n'a de sens que si $\ln x$ \emph{existe}, c'est-à-dire si $x > 0$. Écrire $e^{\ln(-3)} = -3$ est une \textbf{double erreur} : $\ln(-3)$ n'existe pas, et de toute façon une exponentielle ne peut jamais être négative. En revanche, $\ln(e^x) = x$ est valable pour \emph{tout} réel $x$, sans aucune restriction. Retiens : \textbf{l'exponentielle accepte tout, le logarithme n'accepte que le strictement positif}.
\end{attention}

\begin{exemplebox}[Calculs directs]
\begin{itemize}
\item $\ln\left(e^{-2}\right) = -2$ (la relation $\ln(e^x)=x$ s'applique à tout réel, donc en particulier à $x=-2$) ;
\item $e^{\ln 5} = 5$ (car $5 > 0$) ;
\item $\ln\left(e^{0}\right) = 0$, donc $e^0$ est le nombre dont le logarithme vaut $0$ : c'est $1$. D'où $e^0 = 1$ ;
\item $e^{\ln 1} = 1$ et $\ln(e^1) = 1$, ce qui redonne $e^1 = e$ et $\ln e = 1$.
\end{itemize}
\end{exemplebox}

\begin{aretenir}[Valeurs remarquables à connaître par cœur]
\[ e^0 = 1, \qquad e^1 = e \approx 2,718, \qquad \ln e = 1, \qquad \ln 1 = 0. \]
\end{aretenir}

\courssub{Stricte positivité de l'exponentielle}

\begin{propriete}[Pour tout réel $x$, $e^x > 0$]
Pour tout réel $x$, on a $e^x > 0$.
\end{propriete}

\begin{experience}[Démonstration guidée]
Cette démonstration, très courte, est un excellent entraînement à la rédaction.
\begin{enumerate}
\item Par définition, la fonction $\exp$ est à valeurs dans $]0\,;\,+\infty[$ : c'est son \emph{ensemble image} (ensemble d'arrivée de la bijection).
\item Donc pour tout réel $x$, le nombre $e^x = \exp(x)$ appartient à $]0\,;\,+\infty[$.
\item Conclusion : $e^x > 0$ pour tout réel $x$. \hfill$\square$
\end{enumerate}
\end{experience}

\begin{attention}[Piège classique dans les équations et inéquations]
La stricte positivité de $e^x$ a deux conséquences pratiques immédiates :
\begin{itemize}
\item une équation du type $e^x = -4$ ou $e^x = 0$ n'a \textbf{aucune solution} : inutile de chercher ;
\item dans une inéquation, on peut multiplier ou diviser par $e^x$ \textbf{sans changer le sens de l'inégalité}, puisque $e^x$ est strictement positif.
\end{itemize}
\end{attention}

\courssub{L'équivalence fondamentale $\ln x = y \iff x = e^y$}

Voici l'outil qui permet de résoudre toutes les équations faisant intervenir un logarithme.

\begin{propriete}[Équivalence fondamentale]
Pour tout réel $x > 0$ et tout réel $y$ :
\[ \ln x = y \iff x = e^y. \]
\end{propriete}

\begin{experience}[Démonstration guidée de l'équivalence]
Démontrons les deux implications.
\begin{enumerate}
\item \textbf{Sens direct} ($\Rightarrow$). Supposons $\ln x = y$ avec $x > 0$. Appliquons $\exp$ aux deux membres : $e^{\ln x} = e^y$. Or $e^{\ln x} = x$ (relation fondamentale, valable car $x > 0$). Donc $x = e^y$.
\item \textbf{Sens réciproque} ($\Leftarrow$). Supposons $x = e^y$. Alors $x > 0$ (stricte positivité de l'exponentielle), donc $\ln x$ existe. Appliquons $\ln$ aux deux membres : $\ln x = \ln(e^y) = y$ (relation fondamentale, valable pour tout réel $y$).
\end{enumerate}
L'équivalence est donc établie. Remarque au passage que le sens réciproque prouve que la condition $x > 0$ est automatiquement satisfaite quand $x = e^y$. \hfill$\square$
\end{experience}

\begin{exoresolu}[Résoudre $\ln x = 3$]
Résoudre dans $\mathbb{R}$ l'équation $\ln x = 3$.
\tcblower
\textbf{Étape 1 : ensemble de validité.} $\ln x$ n'existe que si $x > 0$. On cherche donc les solutions dans $]0\,;\,+\infty[$.

\textbf{Étape 2 : application de l'équivalence fondamentale.}
\[ \ln x = 3 \iff x = e^3. \]

\textbf{Étape 3 : valeur approchée et conclusion.} À la calculatrice, $e^3 \approx 20,0855$, soit $e^3 \approx 20,09$ au centième. Comme $e^3 > 0$, la solution est bien dans l'ensemble de validité.
\[ \mathcal{S} = \left\{ e^3 \right\} \approx \{20,09\}. \]

\textbf{Vérification.} $\ln(e^3) = 3$ par la relation fondamentale : la solution convient.
\end{exoresolu}

\begin{exoresolu}[Une équation avec condition]
Résoudre dans $\mathbb{R}$ l'équation $\ln(2x - 1) = -2$.
\tcblower
\textbf{Étape 1 : ensemble de validité.} Il faut $2x - 1 > 0$, c'est-à-dire $x > \dfrac{1}{2}$.

\textbf{Étape 2 : équivalence fondamentale.}
\[ \ln(2x-1) = -2 \iff 2x - 1 = e^{-2} \iff x = \frac{1 + e^{-2}}{2}. \]

\textbf{Étape 3 : vérification de la condition.} Comme $e^{-2} > 0$, on a $1 + e^{-2} > 1$, donc $x = \dfrac{1+e^{-2}}{2} > \dfrac{1}{2}$ : la condition est satisfaite. Numériquement, $e^{-2} \approx 0,1353$, donc $x \approx \dfrac{1,1353}{2} \approx 0,568$.

\textbf{Conclusion.} $\mathcal{S} = \left\{ \dfrac{1 + e^{-2}}{2} \right\}$.
\end{exoresolu}

\begin{methode}[Résoudre une équation de la forme $\ln(u(x)) = a$]
\begin{enumerate}
\item Déterminer l'\textbf{ensemble de validité} : résoudre $u(x) > 0$.
\item Appliquer l'équivalence fondamentale : $\ln(u(x)) = a \iff u(x) = e^a$.
\item Résoudre l'équation obtenue (qui ne contient plus de logarithme).
\item \textbf{Confronter} chaque solution trouvée à l'ensemble de validité et conclure.
\end{enumerate}
\end{methode}

\courssub{Lien graphique : symétrie des courbes de $\exp$ et de $\ln$}

Il existe une propriété géométrique générale : les courbes représentatives de deux bijections réciproques sont toujours \cle{symétriques} par rapport à la droite d'équation $y = x$ (la première bissectrice du repère), dans un repère orthonormé.

En effet, dire que le point $M(a\,;\,b)$ appartient à la courbe de $\exp$ signifie $b = e^a$, c'est-à-dire $\ln b = a$ : le point $M'(b\,;\,a)$ appartient alors à la courbe de $\ln$. Or échanger les coordonnées revient exactement à effectuer la symétrie par rapport à la droite $y = x$.

\begin{propriete}[Symétrie des courbes]
Dans un repère orthonormé, les courbes représentatives des fonctions $\exp$ et $\ln$ sont symétriques par rapport à la droite d'équation $y = x$.
\end{propriete}

\begin{popfigure}
\begin{center}
\begin{tikzpicture}
\begin{axis}[
  width=0.85\linewidth,
  axis lines=middle,
  xlabel={$x$}, ylabel={$y$},
  xmin=-3.2, xmax=4.2,
  ymin=-3.2, ymax=4.5,
  xtick={-3,-2,-1,1,2,3,4},
  ytick={-3,-2,-1,1,2,3,4},
  grid=both, grid style={popline!40},
  tick label style={font=\small},
  legend style={at={(0.03,0.97)}, anchor=north west, font=\small},
  samples=200
]
% Droite y = x
\addplot[popmuted, dashed, thick, domain=-3:4] {x};
\addlegendentry{$y = x$}
% Courbe exp
\addplot[popcurve, poporange, very thick, domain=-3.2:1.43] {exp(x)};
\addlegendentry{$y = e^x$}
% Courbe ln
\addplot[popcurve, popblue, very thick, domain=0.05:4.2] {ln(x)};
\addlegendentry{$y = \ln x$}
% Points symétriques
\addplot[only marks, mark=*, popdark] coordinates {(0,1) (1,0) (1,2.71828) (2.71828,1)};
% Flèches de symétrie
\draw[<->, poppurple, thick] (axis cs:0.15,0.95) -- (axis cs:0.95,0.15);
\draw[<->, poppurple, thick] (axis cs:1.15,2.62) -- (axis cs:2.62,1.15);
\node[popdark, font=\small, anchor=west] at (axis cs:0.05,1.15) {$(0\,;\,1)$};
\node[popdark, font=\small, anchor=south] at (axis cs:1.05,0.02) {$(1\,;\,0)$};
\node[popdark, font=\small, anchor=west] at (axis cs:1.05,2.85) {$(1\,;\,e)$};
\node[popdark, font=\small, anchor=south west] at (axis cs:2.75,1.0) {$(e\,;\,1)$};
\end{axis}
\end{tikzpicture}
\end{center}
\legende{Courbes de $\exp$ (orange) et de $\ln$ (bleu) dans un repère orthonormé. Elles sont symétriques par rapport à la droite $y = x$ (pointillés) : les flèches violettes relient des points symétriques, comme $(0\,;\,1)$ et $(1\,;\,0)$, ou $(1\,;\,e)$ et $(e\,;\,1)$.}
\end{popfigure}

Lecture de la figure : la courbe de $\exp$ passe par $(0\,;\,1)$ (car $e^0 = 1$) et par $(1\,;\,e)$ ; par symétrie, celle de $\ln$ passe par $(1\,;\,0)$ (car $\ln 1 = 0$) et par $(e\,;\,1)$. On observe aussi que la courbe de $\exp$ « monte » de plus en plus vite vers $+\infty$ quand $x$ grandit, tandis qu'elle se rapproche de l'axe des abscisses sans jamais le toucher quand $x$ devient très négatif : nous étudierons ces limites en détail dans la section 4.

\begin{aretenir}[L'essentiel de la section]
\begin{itemize}
\item $\exp$ est la \textbf{bijection réciproque} de $\ln$ : elle va de $\mathbb{R}$ vers $]0\,;\,+\infty[$.
\item Relations fondamentales : $\ln(e^x) = x$ pour tout réel $x$ ; $e^{\ln x} = x$ pour tout $x > 0$.
\item Pour tout réel $x$, $e^x > 0$ ; valeurs remarquables : $e^0 = 1$, $e^1 = e \approx 2,718$.
\item Équivalence fondamentale : pour $x > 0$, $\ln x = y \iff x = e^y$.
\item Les courbes de $\exp$ et $\ln$ sont symétriques par rapport à la droite $y = x$.
\end{itemize}
\end{aretenir}

\begin{exercice}[À toi de jouer]
\begin{enumerate}
\item Simplifier : $\ln\left(e^{5}\right)$ ; $e^{\ln 7}$ ; $\ln\left(e^{-1}\right)$ ; $e^{\ln \frac{1}{2}}$.
\item Résoudre dans $\mathbb{R}$ : a) $\ln x = -1$ ; b) $\ln(3x + 2) = 0$ ; c) $\ln(x^2) = 4$ (attention aux conditions !).
\item Un démographe estime que la population $P$ d'une ville (en millions d'habitants) vérifie $\ln P = 0,035\,t + \ln 2$, où $t$ est le temps en années. Exprimer $P$ en fonction de $t$.
\end{enumerate}
\end{exercice}

\begin{corrige}[Exercice : À toi de jouer]
\begin{enumerate}
\item $\ln(e^5) = 5$ ; $e^{\ln 7} = 7$ (car $7 > 0$) ; $\ln(e^{-1}) = -1$ ; $e^{\ln \frac{1}{2}} = \dfrac{1}{2}$ (car $\tfrac12 > 0$).
\item a) $\ln x = -1 \iff x = e^{-1} \approx 0,368$. Solution positive : $\mathcal{S} = \{e^{-1}\}$.

b) Validité : $3x + 2 > 0 \iff x > -\dfrac{2}{3}$. Alors $\ln(3x+2) = 0 \iff 3x + 2 = e^0 = 1 \iff x = -\dfrac{1}{3}$. Or $-\dfrac13 > -\dfrac23$ : $\mathcal{S} = \left\{ -\dfrac{1}{3} \right\}$.

c) Validité : $x^2 > 0 \iff x \neq 0$. Alors $x^2 = e^4$, donc $x = e^2$ ou $x = -e^2$, toutes deux non nulles : $\mathcal{S} = \{-e^2\,;\,e^2\}$.
\item On a $\ln P = \ln 2 + 0,035\,t$. Par l'équivalence fondamentale : $P = e^{\ln 2 + 0,035\,t}$. En anticipant sur les propriétés algébriques de la section suivante, $P = e^{\ln 2} \times e^{0,035\,t} = 2\,e^{0,035\,t}$ : la ville part de $2$ millions d'habitants et croît de façon exponentielle.
\end{enumerate}
\end{corrige}

\coursec{Propriétés algébriques de l'exponentielle}

Dans la section précédente, nous avons défini la fonction exponentielle népérienne comme la \cle{réciproque} du logarithme népérien : pour tout réel $x$, $e^x$ est l'unique réel strictement positif dont le logarithme vaut $x$. Nous avons retenu les deux identités fondamentales
\[
\ln(e^x) = x \quad \text{pour tout } x \in \mathbb{R}, \qquad e^{\ln x} = x \quad \text{pour tout } x > 0.
\]
Mais une fonction ne devient vraiment utile que lorsqu'on sait \emph{calculer} avec elle. Or le logarithme possède un talent remarquable : il transforme les produits en sommes ($\ln(ab) = \ln a + \ln b$). Puisque l'exponentielle fait exactement le chemin inverse, on devine qu'elle va, elle, transformer les \cle{sommes en produits}. C'est précisément le contenu de cette section : la célèbre relation $e^{a+b} = e^a \times e^b$ et toutes ses conséquences. Ces formules sont les outils de base de presque tous les exercices du Baccalauréat portant sur l'exponentielle : apprends-les comme tu connais les tables de multiplication.

\courssub{La relation fonctionnelle fondamentale}

\textbf{Intuition.}\par
Prenons un exemple numérique pour sentir ce qui se passe. On sait que $e^0 = 1$ et $e^1 = e \approx 2{,}718$. Que vaut $e^{1+1} = e^2$ ? Si l'exponentielle « transforme les sommes en produits », on devrait avoir $e^2 = e^1 \times e^1 = e \times e \approx 7{,}389$. Et c'est effectivement le cas : la calculatrice confirme $e^2 \approx 7{,}389$. De même, $e^{0+1} = e^0 \times e^1 = 1 \times e = e$ : la formule reste cohérente. Ce comportement n'est pas un hasard, c'est une loi générale, que nous allons démontrer rigoureusement.

\begin{propriete}[Théorème : relation fonctionnelle de l'exponentielle]
Pour tous réels $a$ et $b$ :
\[
\boxed{e^{a+b} = e^a \times e^b}
\]
Cette démonstration est \textbf{exigible au Baccalauréat} : elle figure parmi les « démonstrations de cours » classiques.
\end{propriete}

\begin{experience}[Démonstration guidée du théorème]
L'idée directrice est simple : pour montrer que deux nombres \emph{strictement positifs} sont égaux, il suffit de montrer qu'ils ont le \emph{même logarithme népérien}, car la fonction $\ln$ est strictement croissante, donc injective (deux nombres distincts ont des images distinctes).

\textbf{Étape 1.} Posons $X = e^a \times e^b$. Comme $e^a > 0$ et $e^b > 0$, le produit $X$ est strictement positif : on peut donc calculer $\ln X$.

\textbf{Étape 2.} Utilisons la propriété algébrique du logarithme : le logarithme d'un produit est la somme des logarithmes.
\[
\ln\!\left(e^a \times e^b\right) = \ln(e^a) + \ln(e^b).
\]

\textbf{Étape 3.} Or, par définition de l'exponentielle comme réciproque de $\ln$, on a $\ln(e^a) = a$ et $\ln(e^b) = b$. Donc :
\[
\ln\!\left(e^a \times e^b\right) = a + b.
\]

\textbf{Étape 4.} D'autre part, toujours par définition de la réciproque :
\[
\ln\!\left(e^{a+b}\right) = a + b.
\]

\textbf{Étape 5.} Les deux nombres strictement positifs $e^a \times e^b$ et $e^{a+b}$ ont le même logarithme népérien. Par injectivité de $\ln$, ils sont égaux :
\[
e^a \times e^b = e^{a+b}. \qquad \blacksquare
\]
\end{experience}

\begin{attention}[L'erreur classique à bannir]
La relation dit $e^{a+b} = e^a \times e^b$, et \textbf{surtout pas} $e^{a+b} = e^a + e^b$ ! Vérifie avec des nombres : $e^{1+1} = e^2 \approx 7{,}39$, alors que $e^1 + e^1 = 2e \approx 5{,}44$. Les deux résultats sont différents. L'exponentielle transforme les \emph{sommes en produits}, pas les sommes en sommes. De même, nous verrons que $(e^a)^n = e^{na}$ et non $e^{(a^n)}$.
\end{attention}

\courssub{Premières conséquences : inverse et quotient}

La relation fonctionnelle est un « théorème mère » : toutes les autres formules algébriques de l'exponentielle en découlent en une ou deux lignes.

\begin{propriete}[Inverse et quotient]
Pour tous réels $a$ et $b$ :
\[
e^{-a} = \frac{1}{e^a} \qquad \text{et} \qquad e^{a-b} = \frac{e^a}{e^b}.
\]
En particulier, $e^a \neq 0$ pour tout réel $a$ : l'exponentielle ne s'annule jamais.
\end{propriete}

\begin{experience}[Démonstrations guidées]
\textbf{Pour l'inverse.} Appliquons la relation fonctionnelle avec $b = -a$ :
\[
e^a \times e^{-a} = e^{a + (-a)} = e^0 = 1.
\]
Le produit $e^a \times e^{-a}$ vaut $1$, ce qui signifie exactement que $e^{-a}$ est l'inverse de $e^a$ : $e^{-a} = \dfrac{1}{e^a}$. $\blacksquare$

\textbf{Pour le quotient.} Écrivons $a - b = a + (-b)$ et appliquons la relation fonctionnelle puis le résultat précédent :
\[
e^{a-b} = e^{a + (-b)} = e^a \times e^{-b} = e^a \times \frac{1}{e^b} = \frac{e^a}{e^b}. \qquad \blacksquare
\]
\end{experience}

\begin{exemplebox}[Application immédiate]
Calculons $\dfrac{e^5}{e^8}$. D'après la formule du quotient :
\[
\frac{e^5}{e^8} = e^{5-8} = e^{-3} = \frac{1}{e^3}.
\]
On retiendra le réflexe : \cle{quotient d'exponentielles} $=$ \cle{soustraction des exposants}.
\end{exemplebox}

\courssub{Puissances d'une exponentielle}

\begin{propriete}[Puissances entières]
Pour tout réel $a$ et tout entier relatif $n \in \mathbb{Z}$ :
\[
\left(e^a\right)^n = e^{na}.
\]
\end{propriete}

\begin{experience}[Démonstration guidée]
\textbf{Cas $n \in \mathbb{N}$, par récurrence.} Pour $n = 0$ : $(e^a)^0 = 1 = e^0 = e^{0 \times a}$ : la propriété est vraie au rang $0$. Supposons-la vraie à un rang $n \in \mathbb{N}$, c'est-à-dire $(e^a)^n = e^{na}$ (hypothèse de récurrence). Alors au rang $n+1$ :
\[
(e^a)^{n+1} = (e^a)^n \times e^a = e^{na} \times e^a = e^{na + a} = e^{(n+1)a},
\]
où l'on a utilisé l'hypothèse de récurrence puis la relation fonctionnelle. La propriété est héréditaire ; étant vraie au rang $0$, elle est vraie pour tout $n \in \mathbb{N}$. $\blacksquare$

\textbf{Cas $n$ négatif.} Si $n < 0$, écrivons $n = -m$ avec $m \in \mathbb{N}^*$. Alors, en utilisant le cas positif déjà établi et la formule de l'inverse :
\[
(e^a)^n = (e^a)^{-m} = \frac{1}{(e^a)^m} = \frac{1}{e^{ma}} = e^{-ma} = e^{na}. \qquad \blacksquare
\]
\end{experience}

\begin{propriete}[Racine carrée d'une exponentielle]
Pour tout réel $a$ :
\[
e^{a/2} = \sqrt{e^a}.
\]
Plus généralement, pour tout entier $n \geq 2$, $e^{a/n}$ est le nombre positif dont la puissance $n$-ième vaut $e^a$.
\end{propriete}

En effet, $\left(e^{a/2}\right)^2 = e^{2 \times \frac{a}{2}} = e^a$ d'après la propriété des puissances, et $e^{a/2} > 0$ : c'est donc bien la racine carrée de $e^a$.

\begin{exemplebox}[Exemples chiffrés]
\begin{itemize}
\item $\left(e^3\right)^2 = e^{3 \times 2} = e^6$ ;
\item $\left(e^x\right)^3 = e^{3x}$ pour tout réel $x$ ;
\item $\sqrt{e} = e^{1/2} \approx 1{,}649$ ; en effet $\left(e^{1/2}\right)^2 = e^1 = e$ ;
\item $\left(e^{-2}\right)^4 = e^{-8} = \dfrac{1}{e^8}$.
\end{itemize}
\end{exemplebox}

La figure suivante résume l'organisation logique de toutes ces formules : tout découle de la relation fonctionnelle.

\begin{popfigure}
\begin{center}
\begin{tikzpicture}[
  node distance=1.6cm and 2.2cm,
  central/.style={rectangle, rounded corners, draw=popblue, fill=popblueL, thick, align=center, font=\bfseries, inner sep=6pt},
  branche/.style={rectangle, rounded corners, draw=poporange, fill=poporangeL, thick, align=center, inner sep=5pt},
  lien/.style={-{Stealth[length=3mm]}, thick, popdark}
]
\node[central, align=center] (rf) {Relation fonctionnelle\\ $e^{a+b} = e^a \times e^b$};
\node[branche, above left=of rf, align=center] (inv) {Inverse\\ $e^{-a} = \dfrac{1}{e^a}$};
\node[branche, below left=of rf, align=center] (quo) {Quotient\\ $e^{a-b} = \dfrac{e^a}{e^b}$};
\node[branche, above right=of rf, align=center] (puis) {Puissances\\ $(e^a)^n = e^{na}$};
\node[branche, below right=of rf, align=center] (rac) {Racine\\ $e^{a/2} = \sqrt{e^a}$};
\draw[lien] (rf) -- (inv) node[midway, above, sloped, font=\small] {$b=-a$};
\draw[lien] (rf) -- (quo) node[midway, below, sloped, font=\small] {$a+(-b)$};
\draw[lien] (rf) -- (puis) node[midway, above, sloped, font=\small] {récurrence};
\draw[lien] (rf) -- (rac) node[midway, below, sloped, font=\small] {$n=\tfrac12$};
\draw[lien, popmuted, dashed] (inv) -- (puis);
\end{tikzpicture}
\end{center}
\legende{Arbre de dérivation logique : toutes les formules algébriques de l'exponentielle découlent de la relation fonctionnelle $e^{a+b} = e^a e^b$.}
\end{popfigure}

\courssub{Égalité et comparaison d'exponentielles}

\begin{propriete}[Égalité et inégalité]
Pour tous réels $a$ et $b$ :
\[
e^a = e^b \iff a = b \qquad \text{et} \qquad e^a < e^b \iff a < b.
\]
Autrement dit, la fonction exponentielle est \cle{strictement croissante} sur $\mathbb{R}$ (ce sera démontré dans la section 5 grâce à sa dérivée ; nous l'admettons ici).
\end{propriete}

\textbf{Justification de l'équivalence d'égalité.} Si $a = b$, alors évidemment $e^a = e^b$. Réciproquement, si $e^a = e^b$, alors en appliquant $\ln$ (qui est définie sur $]0 ; +\infty[$, donc sur les valeurs de l'exponentielle) : $\ln(e^a) = \ln(e^b)$, c'est-à-dire $a = b$. $\blacksquare$

\begin{aretenir}[Le réflexe fondamental]
Face à une égalité ou une inégalité entre exponentielles, on « descend » au niveau des exposants, \textbf{sans changer le sens de l'inégalité} car exp est croissante :
\[
e^{u(x)} = e^{v(x)} \iff u(x) = v(x) \quad ; \quad e^{u(x)} < e^{v(x)} \iff u(x) < v(x).
\]
Ce réflexe sera la clé de la section 3 (équations et inéquations).
\end{aretenir}

\courssub{Simplifier des expressions mêlant exponentielle et logarithme}

\begin{methode}[Simplifier une expression avec des exponentielles]
\begin{enumerate}
\item Repérer les produits, quotients et puissances d'exponentielles.
\item Regrouper les exposants grâce à $e^a e^b = e^{a+b}$, $\dfrac{e^a}{e^b} = e^{a-b}$ et $(e^a)^n = e^{na}$.
\item Si un logarithme apparaît dans un exposant, utiliser $e^{\ln a} = a$ (pour $a > 0$) ou écrire un nombre positif sous forme exponentielle : $a = e^{\ln a}$.
\item Simplifier l'exposant obtenu ; ne pas oublier que $e^0 = 1$.
\end{enumerate}
\end{methode}

\begin{propriete}[Écriture exponentielle d'un nombre positif]
Tout réel $a > 0$ peut s'écrire sous forme exponentielle :
\[
a = e^{\ln a}.
\]
Cette écriture est très utile pour « faire entrer » un nombre dans un calcul d'exponentielles.
\end{propriete}

\begin{exoresolu}[Simplifications types du Baccalauréat]
Simplifier les expressions suivantes :
\[
A = \frac{e^3 \times e^{-5}}{(e^{-1})^2}, \qquad B = e^{2\ln 3}, \qquad C = \ln(e^x) - x, \qquad D = \frac{e^{x+1}}{e^x}, \qquad E = 5\,e^{x} \times e^{-x}.
\]
\tcblower
\textbf{Expression $A$.} On regroupe les exposants : au numérateur $e^3 \times e^{-5} = e^{3-5} = e^{-2}$ ; au dénominateur $(e^{-1})^2 = e^{-2}$. Donc :
\[
A = \frac{e^{-2}}{e^{-2}} = e^{-2-(-2)} = e^0 = 1.
\]
On pouvait aussi tout regrouper d'un coup : $A = e^{3 - 5 + 2} = e^0 = 1$.

\textbf{Expression $B$.} On utilise la propriété $n \ln a = \ln(a^n)$ du logarithme : $2\ln 3 = \ln(3^2) = \ln 9$, donc :
\[
B = e^{2\ln 3} = e^{\ln 9} = 9.
\]

\textbf{Expression $C$.} Par définition de la réciproque, $\ln(e^x) = x$ pour tout réel $x$, donc :
\[
C = x - x = 0.
\]

\textbf{Expression $D$.} Formule du quotient :
\[
D = e^{(x+1) - x} = e^1 = e.
\]
Remarque : le résultat ne dépend pas de $x$ !

\textbf{Expression $E$.} $E = 5 \times e^{x + (-x)} = 5 e^0 = 5 \times 1 = 5$. Là encore, $x$ disparaît.
\end{exoresolu}

\begin{attention}[Pièges de rédaction fréquents]
\begin{itemize}
\item \textbf{Ne jamais écrire} $e^{a+b} = e^a + e^b$ : c'est faux (voir le contre-exemple numérique plus haut).
\item $(e^a)^n = e^{na}$, et non $e^{(a^n)}$. Par exemple $(e^2)^3 = e^6$, alors que $e^{(2^3)} = e^8$.
\item $e^{\ln a} = a$ n'est valable que si $a > 0$ : écrire $e^{\ln(-3)} = -3$ n'a aucun sens, car $\ln(-3)$ n'existe pas.
\item $e^0 = 1$ et non $0$ ; l'exponentielle ne s'annule jamais : l'équation $e^x = 0$ n'a pas de solution.
\item Dans $e^{2\ln 3}$, c'est bien le \emph{coefficient} $2$ qui « monte » en puissance ($2\ln 3 = \ln 9$), grâce à la propriété $n\ln a = \ln(a^n)$ du logarithme.
\end{itemize}
\end{attention}

\courssub{Tableau récapitulatif : le miroir exp / ln}

L'exponentielle et le logarithme sont réciproques l'une de l'autre : chaque formule de l'une est le « reflet » d'une formule de l'autre. Ce parallélisme est une excellente aide-mémoire.

\begin{center}
\begin{tabularx}{\linewidth}{|c|X|X|}
\hline
\rowcolor{popblueL}
\textbf{Opération} & \textbf{Exponentielle (sommes $\to$ produits)} & \textbf{Logarithme (produits $\to$ sommes)} \\
\hline
Produit / somme & $e^{a+b} = e^a \times e^b$ & $\ln(ab) = \ln a + \ln b$ \\
\hline
Inverse & $e^{-a} = \dfrac{1}{e^a}$ & $\ln\!\left(\dfrac{1}{a}\right) = -\ln a$ \\
\hline
Quotient / différence & $e^{a-b} = \dfrac{e^a}{e^b}$ & $\ln\!\left(\dfrac{a}{b}\right) = \ln a - \ln b$ \\
\hline
Puissance & $(e^a)^n = e^{na}$ & $\ln(a^n) = n \ln a$ \\
\hline
Racine carrée & $e^{a/2} = \sqrt{e^a}$ & $\ln(\sqrt{a}) = \dfrac{1}{2}\ln a$ \\
\hline
Élément neutre & $e^0 = 1$ & $\ln 1 = 0$ \\
\hline
Valeur remarquable & $e^1 = e$ & $\ln e = 1$ \\
\hline
Réciprocité & $\ln(e^x) = x$ & $e^{\ln x} = x$ \; ($x > 0$) \\
\hline
\end{tabularx}
\end{center}

\begin{exercice}[À toi de jouer]
\begin{enumerate}
\item Simplifier : $F = \dfrac{(e^2)^3 \times e^{-4}}{e^{3}}$ ; \quad $G = e^{3\ln 2} + e^{\ln 5}$ ; \quad $H = \dfrac{e^{2x+3}}{e^{x-1}}$.
\item Montrer que pour tout réel $x$ : $\left(e^x + e^{-x}\right)^2 - \left(e^x - e^{-x}\right)^2 = 4$.
\item Écrire le nombre $7$ sous forme exponentielle, puis simplifier $7 \times e^{-\ln 7}$.
\end{enumerate}
\end{exercice}

\begin{corrige}[Exercice : À toi de jouer]
\textbf{1.} $F = \dfrac{e^{6} \times e^{-4}}{e^{3}} = \dfrac{e^{2}}{e^{3}} = e^{2-3} = e^{-1} = \dfrac{1}{e}$.

$G = e^{\ln(2^3)} + 5 = e^{\ln 8} + 5 = 8 + 5 = 13$.

$H = e^{(2x+3)-(x-1)} = e^{x+4}$.

\textbf{2.} On développe : $\left(e^x + e^{-x}\right)^2 = e^{2x} + 2 e^x e^{-x} + e^{-2x} = e^{2x} + 2 + e^{-2x}$ car $e^x e^{-x} = e^0 = 1$. De même $\left(e^x - e^{-x}\right)^2 = e^{2x} - 2 + e^{-2x}$. En soustrayant : $(e^{2x} + 2 + e^{-2x}) - (e^{2x} - 2 + e^{-2x}) = 4$. $\blacksquare$

\textbf{3.} $7 = e^{\ln 7}$, donc $7 \times e^{-\ln 7} = e^{\ln 7} \times e^{-\ln 7} = e^{0} = 1$. (On reconnaît que $e^{-\ln 7} = \dfrac{1}{e^{\ln 7}} = \dfrac{1}{7}$.)
\end{corrige}

\begin{savaistu}[Pourquoi l'exponentielle gouverne les phénomènes à taux constant]
La relation $e^{a+b} = e^a \times e^b$ est la raison profonde pour laquelle l'exponentielle modélise tous les phénomènes à \cle{taux d'accroissement constant}. Prenons un exemple concret : une coopérative de cacao de la région du Centre place un capital qui augmente de $5\%$ par an. Chaque année, le capital est \emph{multiplié} par le même facteur $1{,}05$ : après $a$ années puis $b$ années supplémentaires, on a multiplié par $1{,}05^a$ puis par $1{,}05^b$, soit au total par $1{,}05^{a+b} = 1{,}05^a \times 1{,}05^b$. C'est exactement la relation fonctionnelle ! Le même mécanisme « multiplicatif » décrit la décroissance radioactive des isotopes en physique, la propagation d'une rumeur sur les réseaux sociaux, ou la croissance d'une population de bactéries. Les mathématiciens démontrent que l'exponentielle est (à un facteur près) la \emph{seule} fonction continue vérifiant $f(a+b) = f(a) \times f(b)$ : cette relation fonctionnelle la caractérise entièrement. C'est d'ailleurs ainsi que Leonhard Euler, au XVIII\textsuperscript{e} siècle, a introduit le nombre $e$ qui porte aujourd'hui son initiale.
\end{savaistu}

\begin{aretenir}[L'essentiel de la section]
\begin{itemize}
\item \textbf{Relation mère} : $e^{a+b} = e^a \times e^b$ (démonstration exigible, via $\ln$ et l'injectivité).
\item \textbf{Conséquences} : $e^{-a} = \dfrac{1}{e^a}$ ; \; $e^{a-b} = \dfrac{e^a}{e^b}$ ; \; $(e^a)^n = e^{na}$ ($n \in \mathbb{Z}$) ; \; $e^{a/2} = \sqrt{e^a}$.
\item \textbf{Comparaison} : $e^a = e^b \iff a = b$ et $e^a < e^b \iff a < b$ (croissance stricte).
\item \textbf{Lien avec ln} : pour $a > 0$, $a = e^{\ln a}$ ; et pour tout réel $x$, $\ln(e^x) = x$.
\item \textbf{Réflexe de calcul} : regrouper les exposants, simplifier, et se souvenir que $e^0 = 1$ et que $e^x > 0$ pour tout réel $x$.
\end{itemize}
\end{aretenir}

\coursec{Équations, inéquations et systèmes avec $e^x$}

Dans la section précédente, nous avons établi les propriétés algébriques de l'exponentielle : $e^{a+b} = e^a\,e^b$, $e^{-a} = \dfrac{1}{e^a}$, $(e^a)^n = e^{na}$. Ces formules transforment les produits en sommes et les quotients en différences. Nous allons maintenant exploiter un fait encore plus puissant : la fonction exponentielle est \cle{strictement croissante} sur $\mathbb{R}$, donc deux nombres ont la même image par $\exp$ si et seulement s'ils sont égaux. C'est cette propriété qui va nous permettre de « descendre » de l'exponentielle à l'exposant, et de résoudre toutes sortes d'équations, d'inéquations et de systèmes où l'inconnue se cache dans un exposant.

\courssub{Le principe fondamental : injectivité et croissance de l'exponentielle}

\textbf{Intuition.} Imagine la courbe de la fonction exponentielle : elle monte sans cesse, sans jamais redescendre ni marquer de palier. Une droite horizontale ne peut donc la couper qu'en \emph{un seul point} (ou pas du tout). Autrement dit, si deux réels $A$ et $B$ ont la même exponentielle, c'est qu'il s'agit du même réel. De même, l'ordre est conservé : si $A < B$, alors $e^A < e^B$, car la courbe monte.

\begin{propriete}[Équivalences fondamentales]
Pour tous réels $A$ et $B$ :
\[ e^A = e^B \iff A = B \qquad \text{et} \qquad e^A < e^B \iff A < B. \]
De même : $e^A \leq e^B \iff A \leq B$ et $e^A > e^B \iff A > B$.
\end{propriete}

\begin{experience}[Pourquoi c'est vrai]
Rappelons que $\exp$ est la bijection réciproque de $\ln$, et que $\ln$ est strictement croissante sur $]0\,;\,+\infty[$.
\begin{enumerate}
\item \textbf{Injectivité.} Supposons $e^A = e^B$. Comme $e^A > 0$ et $e^B > 0$, on peut composer par $\ln$ : $\ln(e^A) = \ln(e^B)$, c'est-à-dire $A = B$. La réciproque est immédiate : si $A = B$, alors $e^A = e^B$.
\item \textbf{Conservation de l'ordre.} Supposons $e^A < e^B$. En composant par $\ln$, qui est strictement croissante et conserve donc le sens de l'inégalité : $\ln(e^A) < \ln(e^B)$, soit $A < B$. Réciproquement, si $A < B$, alors $e^A < e^B$ car $\exp$ est strictement croissante (sa dérivée $e^x$ est strictement positive partout).
\end{enumerate}
La stricte croissance est donc la clé de voûte de toute cette section : elle garantit que l'on peut « passer au logarithme » ou « passer à l'exponentielle » dans une équation ou une inéquation \emph{sans jamais changer le sens des inégalités}.
\end{experience}

\begin{aretenir}[La règle d'or]
L'exponentielle est strictement croissante : dans une équation ou une inéquation entre exponentielles, on peut supprimer les $\exp$ en \textbf{conservant le sens}. C'est exactement l'analogue de ce que fait le logarithme : $\ln a = \ln b \iff a = b$ (pour $a, b > 0$).
\end{aretenir}

\courssub{Équations du type $e^{u(x)} = k$ : discussion selon le signe de $k$}

L'exponentielle ne prend que des valeurs \textbf{strictement positives}. Cette remarque, toute simple, décide à elle seule du sort de l'équation $e^{u(x)} = k$.

\begin{propriete}[Résolution de $e^{u(x)} = k$]
Soit $u$ une fonction définie sur un ensemble $D$ et $k$ un réel.
\begin{itemize}
\item Si $k \leq 0$ : l'équation $e^{u(x)} = k$ \textbf{n'a aucune solution} (car $e^{u(x)} > 0$ pour tout $x$).
\item Si $k > 0$ : $e^{u(x)} = k \iff e^{u(x)} = e^{\ln k} \iff u(x) = \ln k$.
\end{itemize}
\end{propriete}

\begin{popfigure}
\begin{center}
\begin{tikzpicture}
\begin{axis}[
  width=0.85\linewidth, height=6cm,
  axis lines=middle, xlabel={$x$}, ylabel={$y$},
  xmin=-3.2, xmax=3.2, ymin=-1.2, ymax=6.5,
  xtick={-2,-1,1,2}, ytick={1,2,4},
  legend style={at={(0.03,0.97)}, anchor=north west, font=\small},
  clip=true]
\addplot[popcurve,domain=-3.1:1.85,samples=200]{exp(x)};
\addlegendentry{$y = e^x$}
\addplot[poporange,thick,domain=-3.1:3.1]{2};
\addlegendentry{$y = k = 2$ (une solution)}
\addplot[poppurple,thick,dashed,domain=-3.1:3.1]{-0.6};
\addlegendentry{$y = k = -0{,}6$ (aucune solution)}
\addplot[popgreen,mark=*,only marks,mark size=2pt] coordinates {(0.693,2)};
\node[popgreen,font=\small,anchor=south west] at (axis cs:0.75,2.1) {$(\ln 2\,;\,2)$};
\end{axis}
\end{tikzpicture}
\end{center}
\legende{Interprétation graphique de $e^x = k$ : la droite horizontale $y = k$ coupe la courbe de $\exp$ en un unique point si $k > 0$, et ne la coupe jamais si $k \leq 0$.}
\end{popfigure}

\begin{exemplebox}[Trois cas rapides]
\begin{enumerate}
\item $e^{2x-1} = 5$ : comme $5 > 0$, on a $2x - 1 = \ln 5$, soit $x = \dfrac{1 + \ln 5}{2}$. Solution unique : $S = \left\lbrace \dfrac{1+\ln 5}{2}\right\rbrace$.
\item $e^{x^2} = -3$ : comme $-3 < 0$, il n'y a \textbf{aucune solution} : $S = \varnothing$.
\item $e^{3-x} = 1$ : $3 - x = \ln 1 = 0$, donc $x = 3$.
\end{enumerate}
\end{exemplebox}

\begin{attention}[Le piège du signe de $k$]
L'erreur la plus fréquente au Baccalauréat est d'écrire mécaniquement $e^{u(x)} = k \iff u(x) = \ln k$ \emph{sans vérifier que $k > 0$}. Or $\ln k$ n'existe pas si $k \leq 0$ ! Avant tout calcul, observe le membre de droite : s'il est négatif ou nul, conclus immédiatement $S = \varnothing$ — c'est un point facile à gagner.
\end{attention}

\courssub{Équations et inéquations avec le logarithme : la condition d'existence d'abord}

Pour les équations faisant intervenir $\ln$, le principe est symétrique : $\ln A = \ln B \iff A = B$ (avec $A, B > 0$), et $\ln(u(x)) = k \iff u(x) = e^k$. Mais il y a une différence capitale avec l'exponentielle : le logarithme n'est défini que sur $]0\,;\,+\infty[$. Il faut donc \textbf{toujours} commencer par déterminer l'ensemble de résolution.

\begin{methode}[Résoudre une équation ou inéquation avec $\ln$]
\begin{enumerate}
\item Déterminer l'\textbf{ensemble de résolution} $D$ : on impose que chaque quantité sous un $\ln$ soit strictement positive.
\item Transformer l'équation à l'aide de $\ln A = \ln B \iff A = B$ ou de $\ln(u(x)) = k \iff u(x) = e^k$.
\item Résoudre l'équation obtenue (souvent polynomiale).
\item \textbf{Ne retenir que les solutions appartenant à $D$} : toute solution hors de $D$ est à rejeter, même si elle est correcte algébriquement.
\end{enumerate}
\end{methode}

\begin{exoresolu}[Équation avec logarithme]
Résoudre dans $\mathbb{R}$ l'équation $\ln(x+1) + \ln(x-2) = \ln 4$.
\tcblower
\textbf{Étape 1 — Ensemble de résolution.} Il faut $x + 1 > 0$ et $x - 2 > 0$, c'est-à-dire $x > -1$ et $x > 2$. Donc $D = \left]2\,;\,+\infty\right[$.

\textbf{Étape 2 — Transformation.} Pour $x \in D$ :
\[ \ln\big((x+1)(x-2)\big) = \ln 4 \iff (x+1)(x-2) = 4. \]

\textbf{Étape 3 — Résolution.} $x^2 - x - 2 = 4 \iff x^2 - x - 6 = 0$. Le discriminant est $\Delta = 1 + 24 = 25$, d'où $x = \dfrac{1+5}{2} = 3$ ou $x = \dfrac{1-5}{2} = -2$.

\textbf{Étape 4 — Filtrage.} $x = -2 \notin D$ : à rejeter. $x = 3 \in D$ : valide. Vérification : $\ln 4 + \ln 1 = \ln 4$. \checkmark

Conclusion : $S = \lbrace 3 \rbrace$.
\end{exoresolu}

\begin{attention}[Inéquations avec $\ln$ : ne jamais sauter l'étape du domaine]
Pour résoudre $\ln(u(x)) \geq k$, il faut \emph{d'abord} imposer $u(x) > 0$, \emph{ensuite} écrire $u(x) \geq e^k$ (le sens est conservé car $\exp$ est croissante), puis \textbf{intersecter} les deux conditions. Oublier la condition d'existence conduit à accepter des valeurs pour lesquelles le logarithme n'existe même pas : c'est une faute grave de logique, sanctionnée au Bac.
\end{attention}

\begin{exemplebox}[Une inéquation avec $\ln$]
Résoudre $\ln(2x - 3) \geq 1$.
\begin{itemize}
\item Condition d'existence : $2x - 3 > 0 \iff x > \dfrac{3}{2}$.
\item L'inéquation équivaut à $2x - 3 \geq e^1 = e$, soit $x \geq \dfrac{3 + e}{2}$.
\item Comme $\dfrac{3+e}{2} > \dfrac{3}{2}$, la condition d'existence est automatiquement satisfaite.
\end{itemize}
Conclusion : $S = \left[\dfrac{3+e}{2}\,;\,+\infty\right[$.
\end{exemplebox}

\courssub{Inéquations du type $e^{u(x)} \leq k$}

Le même raisonnement s'applique, avec la croissance de $\exp$ qui conserve le sens des inégalités.

\begin{propriete}[Résolution de $e^{u(x)} \leq k$]
\begin{itemize}
\item Si $k \leq 0$ : aucune solution (une exponentielle est strictement positive).
\item Si $k > 0$ : $e^{u(x)} \leq k \iff u(x) \leq \ln k$.
\end{itemize}
Variante utile : $e^{u(x)} \geq k$ avec $k \leq 0$ est \textbf{toujours vraie} (solution : tout l'ensemble de définition de $u$) ; avec $k > 0$, elle équivaut à $u(x) \geq \ln k$.
\end{propriete}

\begin{exemplebox}[Inéquation exponentielle]
Résoudre $e^{1-2x} \leq 3$.
\[ e^{1-2x} \leq 3 \iff 1 - 2x \leq \ln 3 \iff -2x \leq \ln 3 - 1 \iff x \geq \frac{1 - \ln 3}{2} \]
(attention au changement de sens lors de la division par $-2$ !). Donc $S = \left[\dfrac{1-\ln 3}{2}\,;\,+\infty\right[$.
\end{exemplebox}

\courssub{Le changement d'inconnue $X = e^x$ : équations du second degré déguisées}

\textbf{Intuition.} Observe l'équation $e^{2x} - 3e^x + 2 = 0$. Grâce à la propriété $(e^x)^2 = e^{2x}$, elle ressemble trait pour trait à une équation du second degré : il suffit de « photographier » $e^x$ sous le nom $X$. C'est le célèbre \cle{changement d'inconnue}, une technique que tu as déjà rencontrée avec les équations bicarrées.

\begin{methode}[Résoudre $a\,e^{2x} + b\,e^x + c = 0$]
\begin{enumerate}
\item Poser $X = e^x$, avec la contrainte fondamentale $X > 0$.
\item Résoudre l'équation du second degré $aX^2 + bX + c = 0$.
\item \textbf{Écarter toute racine négative ou nulle} : une telle racine ne correspond à aucun réel $x$, puisque $e^x > 0$.
\item Pour chaque racine strictement positive $X_i$, conclure $x_i = \ln X_i$.
\end{enumerate}
\end{methode}

\begin{popfigure}
\begin{center}
\begin{tikzpicture}[
  node distance=0.55cm,
  etape/.style={rectangle, rounded corners=3pt, draw=popblue, fill=popblueL, text width=5.6cm, align=center, font=\small, inner sep=5pt},
  decision/.style={diamond, draw=poporange, fill=poporangeL, aspect=2.2, align=center, font=\small, inner sep=2pt},
  fin/.style={rectangle, rounded corners=3pt, draw=popgreen, fill=popgreenL, text width=5.6cm, align=center, font=\small, inner sep=5pt},
  fleche/.style={-{Stealth[length=2.5mm]}, thick, popdark}]
\node[etape] (e1) {Équation en $e^x$ : $a\,e^{2x} + b\,e^x + c = 0$};
\node[etape, below=of e1] (e2) {Poser $X = e^x$ \quad (contrainte : $X > 0$)};
\node[etape, below=of e2] (e3) {Résoudre le trinôme $aX^2 + bX + c = 0$};
\node[decision, below=of e3] (d1) {Racine $X > 0$ ?};
\node[fin, below left=0.6cm and 1.6cm of d1] (oui) {\textbf{Oui} : retenir $x = \ln X$};
\node[fin, below right=0.6cm and 1.6cm of d1, draw=poppink, fill=poppinkL] (non) {\textbf{Non} ($X \leq 0$) : rejeter la racine};
\draw[fleche] (e1) -- (e2);
\draw[fleche] (e2) -- (e3);
\draw[fleche] (e3) -- (d1);
\draw[fleche] (d1) -- (oui);
\draw[fleche] (d1) -- (non);
\end{tikzpicture}
\end{center}
\legende{Organigramme de résolution d'une équation en $e^x$ par changement d'inconnue $X = e^x$.}
\end{popfigure}

\begin{exoresolu}[L'exemple fondateur]
Résoudre dans $\mathbb{R}$ : $e^{2x} - 3e^x + 2 = 0$.
\tcblower
\textbf{Étape 1.} Posons $X = e^x$ ($X > 0$). Comme $e^{2x} = (e^x)^2 = X^2$, l'équation devient :
\[ X^2 - 3X + 2 = 0. \]
\textbf{Étape 2.} Discriminant : $\Delta = 9 - 8 = 1 > 0$. Racines : $X_1 = \dfrac{3-1}{2} = 1$ et $X_2 = \dfrac{3+1}{2} = 2$.
\textbf{Étape 3.} Les deux racines sont strictement positives : on les retient toutes les deux.
\textbf{Étape 4.} $e^x = 1 \iff x = \ln 1 = 0$ ; \quad $e^x = 2 \iff x = \ln 2$.

Vérification pour $x = \ln 2$ : $e^{2\ln 2} - 3e^{\ln 2} + 2 = 4 - 6 + 2 = 0$. \checkmark

Conclusion : $S = \lbrace 0 \,;\, \ln 2 \rbrace$.
\end{exoresolu}

\begin{attention}[Une racine négative ne donne aucune solution]
Considère $e^{2x} + e^x - 2 = 0$. Avec $X = e^x$ : $X^2 + X - 2 = 0$, de racines $X = 1$ et $X = -2$. La racine $X = -2$ doit être \textbf{rejetée} : l'équation $e^x = -2$ n'a pas de solution. Seule subsiste $x = \ln 1 = 0$. Écrire « $x = \ln(-2)$ » est une absurdité qui coûte cher : $\ln$ d'un nombre négatif n'existe pas. Contrôle systématiquement le signe de chaque racine avant de conclure.
\end{attention}

\courssub{Systèmes d'inconnues $(e^x, e^y)$ ou $(e^x, \ln y)$}

Les propriétés algébriques de la section 2 transforment les systèmes multiplicatifs en systèmes additifs, que l'on résout ensuite par substitution ou combinaison.

\begin{exemplebox}[Système avec exponentielles]
Résoudre le système $\begin{cases} e^x \times e^y = e^3 \\ \dfrac{e^x}{e^y} = e \end{cases}$.

Grâce aux relations $e^x \times e^y = e^{x+y}$ et $\dfrac{e^x}{e^y} = e^{x-y}$, le système équivaut à :
\[ \begin{cases} e^{x+y} = e^3 \\ e^{x-y} = e^1 \end{cases} \iff \begin{cases} x + y = 3 \\ x - y = 1 \end{cases} \]
Par addition des deux lignes : $2x = 4$, donc $x = 2$ ; puis $y = 3 - x = 1$.
Conclusion : $S = \lbrace (2\,;\,1) \rbrace$. Vérification : $e^2 \times e^1 = e^3$ \checkmark\ et $\dfrac{e^2}{e^1} = e$ \checkmark.
\end{exemplebox}

\begin{methode}[Systèmes mixtes avec $e^x$ et $\ln y$]
Pour un système du type $\begin{cases} e^x \cdot y = a \\ x + \ln y = b \end{cases}$ :
\begin{enumerate}
\item Noter les conditions d'existence (ici $y > 0$).
\item Poser des inconnues auxiliaires : $X = e^x$ ($X > 0$) et $Y = \ln y$, ou transformer chaque équation pour faire apparaître un système linéaire en $x$ et $\ln y$.
\item Résoudre le système obtenu (substitution ou combinaison).
\item Revenir aux inconnues initiales : $x = \ln X$, $y = e^Y$, en vérifiant les conditions ($X > 0$, $y > 0$).
\end{enumerate}
\end{methode}

\begin{exoresolu}[Système mixte]
Résoudre $\begin{cases} x + \ln y = 2 \\ 2x - \ln y = 1 \end{cases}$ d'inconnues $(x, y) \in \mathbb{R} \times \left]0\,;\,+\infty\right[$.
\tcblower
\textbf{Condition d'existence :} $y > 0$.

Posons $Y = \ln y$. Le système devient linéaire : $\begin{cases} x + Y = 2 \\ 2x - Y = 1 \end{cases}$.

Par addition : $3x = 3$, donc $x = 1$ ; puis $Y = 2 - x = 1$.

Retour à $y$ : $\ln y = 1 \iff y = e^1 = e$, qui est bien strictement positif.

Conclusion : $S = \lbrace (1\,;\,e) \rbrace$. Vérification : $1 + \ln e = 1 + 1 = 2$ \checkmark\ et $2 - 1 = 1$ \checkmark.
\end{exoresolu}

\courssub{Bilan et entraînement}

\begin{aretenir}[L'essentiel de la section]
\begin{itemize}
\item $e^A = e^B \iff A = B$ et $e^A < e^B \iff A < B$ : la stricte croissance de $\exp$ fait tout le travail.
\item $e^{u(x)} = k$ : aucune solution si $k \leq 0$ ; $u(x) = \ln k$ si $k > 0$.
\item Avec $\ln$ : toujours déterminer l'ensemble de résolution \emph{avant} de lever le logarithme, et filtrer les solutions à la fin.
\item $a\,e^{2x} + b\,e^x + c = 0$ : poser $X = e^x > 0$, résoudre le trinôme, écarter les racines $\leq 0$, conclure par $x = \ln X$.
\item Systèmes : les propriétés algébriques ramènent à un système linéaire en $x$, $y$ ou $\ln y$.
\end{itemize}
\end{aretenir}

\begin{savaistu}[Pourquoi tant de soin avec les conditions ?]
Ces techniques de résolution sont le cœur de la modélisation exponentielle que nous verrons en section 6 : datation au carbone 14, décroissance radioactive, remboursement d'un crédit à la microfinance, propagation d'une rumeur sur les réseaux sociaux. Dans chacun de ces problèmes, une solution algébrique qui violerait une condition d'existence (un temps négatif, une quantité de matière négative) serait physiquement absurde. Les mathématiciens du XVIII\textsuperscript{e} siècle comme Euler, qui a popularisé la notation $e$, ont mis des décennies à discipliner l'usage des logarithmes : exige la même rigueur de toi-même, et le correcteur du Bac te le rendra.
\end{savaistu}

\begin{exercice}[À toi de jouer]
\begin{enumerate}
\item Résoudre dans $\mathbb{R}$ : $e^{2x} - 4e^x + 3 = 0$.
\item Résoudre dans $\mathbb{R}$ : $e^{2x} + 2e^x - 3 = 0$.
\item Résoudre : $\ln(x+3) \leq 2$.
\item Résoudre le système $\begin{cases} e^x \cdot e^y = e^5 \\ e^x / e^y = e^3 \end{cases}$.
\end{enumerate}
\end{exercice}

\begin{corrige}[Exercice précédent]
\textbf{1.} Posons $X = e^x > 0$ : $X^2 - 4X + 3 = 0$, $\Delta = 16 - 12 = 4$, $X = 1$ ou $X = 3$, toutes deux positives. Donc $x = 0$ ou $x = \ln 3$ : $S = \lbrace 0\,;\,\ln 3\rbrace$.

\textbf{2.} $X^2 + 2X - 3 = 0$, $\Delta = 4 + 12 = 16$, $X = 1$ ou $X = -3$. La racine $-3$ est rejetée. $S = \lbrace 0 \rbrace$.

\textbf{3.} Condition : $x > -3$. L'inéquation donne $x + 3 \leq e^2$, soit $x \leq e^2 - 3$. Comme $e^2 - 3 > -3$, on obtient $S = \left]-3\,;\,e^2 - 3\right]$.

\textbf{4.} Le système équivaut à $\begin{cases} x + y = 5 \\ x - y = 3 \end{cases}$, d'où $x = 4$ et $y = 1$ : $S = \lbrace (4\,;\,1)\rbrace$.
\end{corrige}

Tu maîtrises désormais la résolution des équations, inéquations et systèmes faisant intervenir l'exponentielle. La section suivante étudiera le comportement de cette fonction aux bornes de son ensemble de définition : les limites, étape indispensable avant l'étude complète des fonctions exponentielles.

\coursec{Limites de la fonction exponentielle}

Dans la section précédente, tu as appris à \emph{résoudre} des équations et des inéquations où l'inconnue se cache dans un exposant : tu as vu que l'exponentielle et le logarithme se « défont » mutuellement. Mais résoudre ne suffit pas. Pour étudier une fonction contenant $\mathrm{e}^x$ — et ce sera le cœur de nombreux problèmes du Baccalauréat — il faut savoir ce qui se passe \emph{aux extrémités} : que devient $\mathrm{e}^x$ lorsque $x$ devient très grand ? Lorsque $x$ devient très petit (négativement) ? C'est tout l'objet de cette section : les \cle{limites} de la fonction exponentielle, d'abord aux bornes de son ensemble de définition, puis pour toutes les fonctions composées de la forme $\mathrm{e}^{u(x)}$.

\courssub{Limites aux bornes de $\mathbb{R}$}

\textbf{Intuition.}\par
Rappelle-toi la croissance d'une population de bactéries qui double chaque heure : après 1 heure, 2 individus ; après 10 heures, plus de 1000 ; après 20 heures, plus d'un million. Une croissance \emph{exponentielle} est une croissance qui s'emballe : elle ne ralentit jamais, elle accélère sans cesse. À l'inverse, remonte le temps : la population est divisée par 2 chaque heure écoulée en arrière ; elle se rapproche de zéro sans jamais l'atteindre. C'est exactement le comportement de la fonction exponentielle.

\begin{propriete}[Théorème (admis) : limites de l'exponentielle aux bornes]
La fonction exponentielle vérifie :
\[ \lim_{x \to +\infty} \mathrm{e}^x = +\infty \qquad \text{et} \qquad \lim_{x \to -\infty} \mathrm{e}^x = 0. \]
Ce résultat est admis au programme de Terminale ; il découle de la définition de $\exp$ comme réciproque de $\ln$ : l'ensemble image de $\ln$ étant $\mathbb{R}$ tout entier, l'exponentielle prend toutes les valeurs de $]0\,;\,+\infty[$ et « suit » le comportement de $\ln$ en miroir.
\end{propriete}

\begin{aretenir}[Interprétation graphique : l'asymptote horizontale]
Puisque $\displaystyle\lim_{x \to -\infty} \mathrm{e}^x = 0$, la courbe représentative de la fonction exponentielle se rapproche indéfiniment de la droite d'équation $y = 0$ lorsque $x$ tend vers $-\infty$, sans jamais la toucher (car $\mathrm{e}^x > 0$ pour tout réel $x$). On dit que \cle{l'axe des abscisses est asymptote horizontale} à la courbe de $\exp$ en $-\infty$.
En revanche, en $+\infty$, la courbe « monte » de plus en plus vite : il n'y a \textbf{aucune} asymptote en $+\infty$.
\end{aretenir}

\begin{popfigure}
\begin{center}
\begin{tikzpicture}
\begin{axis}[
    width=0.85\linewidth, height=6.5cm,
    axis lines=middle,
    xlabel={$x$}, ylabel={$y$},
    xmin=-4.5, xmax=3.2, ymin=-0.6, ymax=8.5,
    xtick={-4,-3,-2,-1,1,2,3}, ytick={1,2,4,6,8},
    tick label style={font=\small},
    label style={font=\small},
    clip=false
]
\addplot[popcurve,domain=-4.5:2.1,samples=200]{exp(x)};
\addplot[poporange,dashed,thick,domain=-4.5:3.2]{0};
\node[popdark,font=\small] at (axis cs:1.4,5.4) {$y=\mathrm{e}^x$};
\node[poporange,font=\small,anchor=south west] at (axis cs:-4.4,0.15) {asymptote $y=0$ en $-\infty$};
\draw[->,popgreen,thick] (axis cs:-2.2,0.35) -- (axis cs:-3.6,0.12);
\node[popgreen,font=\small,anchor=east] at (axis cs:-2.2,0.5) {$\mathrm{e}^x \to 0$};
\draw[->,popblue,thick] (axis cs:1.7,4.5) -- (axis cs:2.05,7.3);
\node[popblue,font=\small,anchor=west] at (axis cs:1.55,6.6) {$\mathrm{e}^x \to +\infty$};
\end{axis}
\end{tikzpicture}
\end{center}
\legende{Courbe de la fonction exponentielle : asymptote horizontale $y = 0$ en $-\infty$, croissance explosive en $+\infty$.}
\end{popfigure}

\begin{attention}[Erreur fréquente : inverser les deux limites]
Beaucoup d'élèves écrivent, par automatisme mal placé, que « $\mathrm{e}^x$ tend vers 0 »... en $+\infty$ ! C'est \textbf{faux}. Retiens bien :
\begin{itemize}
\item c'est en $\boldsymbol{-\infty}$ que la courbe rejoint l'axe des abscisses ($\mathrm{e}^x \to 0$) ;
\item en $\boldsymbol{+\infty}$, l'exponentielle \emph{explose} : $\mathrm{e}^x \to +\infty$.
\end{itemize}
Une image mentale : l'exponentielle est une fusée qui décolle vers la droite ; vers la gauche, elle se couche doucement sur le sol.
\end{attention}

\courssub{Limites par composition : limite de $\mathrm{e}^{u(x)}$}

En pratique, l'exposant n'est presque jamais $x$ tout seul : c'est une expression $u(x)$, par exemple $-2x+1$, $x^2-3$ ou $\dfrac{1}{x}$. La règle de composition des limites s'applique alors naturellement : on regarde d'abord où va l'exposant, puis on applique l'exponentielle au résultat.

\begin{propriete}[Limite d'une composée avec l'exponentielle]
Soit $u$ une fonction définie au voisinage d'un point $a$ (réel, $+\infty$ ou $-\infty$).
\begin{itemize}
\item Si $\displaystyle\lim_{x \to a} u(x) = \ell$ (réel), alors $\displaystyle\lim_{x \to a} \mathrm{e}^{u(x)} = \mathrm{e}^{\ell}$ ;
\item Si $\displaystyle\lim_{x \to a} u(x) = +\infty$, alors $\displaystyle\lim_{x \to a} \mathrm{e}^{u(x)} = +\infty$ ;
\item Si $\displaystyle\lim_{x \to a} u(x) = -\infty$, alors $\displaystyle\lim_{x \to a} \mathrm{e}^{u(x)} = 0$.
\end{itemize}
\end{propriete}

\begin{methode}[Déterminer la limite de $\mathrm{e}^{ax+b}$ en $\pm\infty$]
\begin{enumerate}
\item \textbf{Étape 1 :} déterminer la limite de l'\cle{exposant} $u(x) = ax + b$ en la borne considérée (c'est une fonction affine : sa limite dépend du signe de $a$).
\item \textbf{Étape 2 :} appliquer la règle de composition : exposant qui tend vers $+\infty$ donne $+\infty$ ; exposant qui tend vers $-\infty$ donne $0$ ; exposant qui tend vers un réel $\ell$ donne $\mathrm{e}^{\ell}$.
\item \textbf{Étape 3 :} conclure en rédigeant proprement : « par composition des limites, ... ».
\end{enumerate}
\end{methode}

\begin{exemplebox}[Deux limites de référence à connaître par cœur]
\textbf{1.} Déterminons $\displaystyle\lim_{x \to +\infty} \mathrm{e}^{-2x+1}$.

\emph{Étape 1 :} l'exposant est $u(x) = -2x + 1$, fonction affine de coefficient $-2 < 0$, donc $\displaystyle\lim_{x \to +\infty} (-2x+1) = -\infty$.

\emph{Étape 2 :} par composition, $\displaystyle\lim_{x \to +\infty} \mathrm{e}^{-2x+1} = 0$.

\textbf{2.} Déterminons $\displaystyle\lim_{x \to +\infty} \mathrm{e}^{3x-5}$.

\emph{Étape 1 :} l'exposant $3x - 5$ a un coefficient $3 > 0$, donc $\displaystyle\lim_{x \to +\infty} (3x-5) = +\infty$.

\emph{Étape 2 :} par composition, $\displaystyle\lim_{x \to +\infty} \mathrm{e}^{3x-5} = +\infty$.

\textbf{Bilan à retenir :} en $+\infty$, $\mathrm{e}^{ax+b}$ tend vers $+\infty$ si $a > 0$, et vers $0$ si $a < 0$ (le signe de $a$ décide de tout ; $b$ ne joue aucun rôle dans la limite).
\end{exemplebox}

\begin{exoresolu}[Limite d'une composée avec un exposant non affine]
\textbf{Énoncé.} Déterminer les limites suivantes :
\[ \text{a) } \lim_{x \to +\infty} \mathrm{e}^{\,2 - \frac{3}{x}} \qquad ; \qquad \text{b) } \lim_{\substack{x \to 0 \\ x > 0}} \mathrm{e}^{-1/x}. \]
\tcblower
\textbf{Solution détaillée.}

\textbf{a)} Posons $u(x) = 2 - \dfrac{3}{x}$. Lorsque $x \to +\infty$, $\dfrac{3}{x} \to 0$, donc $\displaystyle\lim_{x \to +\infty} u(x) = 2$ (limite finie).

Par composition : $\displaystyle\lim_{x \to +\infty} \mathrm{e}^{\,2 - \frac{3}{x}} = \mathrm{e}^{2}$.

\textbf{b)} Posons $u(x) = -\dfrac{1}{x}$. Lorsque $x \to 0$ avec $x > 0$, $\dfrac{1}{x} \to +\infty$, donc $u(x) \to -\infty$.

Par composition : $\displaystyle\lim_{\substack{x \to 0 \\ x > 0}} \mathrm{e}^{-1/x} = 0$.

\emph{Remarque :} dans les deux cas, on n'a fait qu'\textbf{une} chose : regarder où va l'exposant, puis appliquer la règle des trois cas. C'est toujours ainsi qu'on procède, quel que soit l'exposant.
\end{exoresolu}

\courssub{Complément Bac : les croissances comparées (admis)}

\begin{savaistu}[Hors programme strict, mais classique au Baccalauréat]
Le résultat qui suit n'est pas exigible en tant que tel dans le programme MINESEC, mais il apparaît très fréquemment dans les sujets de Bac (souvent admis dans l'énoncé, ou à établir à partir d'une étude de fonction). Le connaître te fera gagner un temps précieux.
\end{savaistu}

\begin{propriete}[Croissances comparées (admises)]
\[ \lim_{x \to +\infty} \frac{\mathrm{e}^x}{x} = +\infty \qquad \text{et} \qquad \lim_{x \to -\infty} x\,\mathrm{e}^x = 0. \]
Plus généralement, pour tout entier $n \geq 1$ : $\displaystyle\lim_{x \to +\infty} \frac{\mathrm{e}^x}{x^n} = +\infty$ et $\displaystyle\lim_{x \to -\infty} x^n \mathrm{e}^x = 0$.
\end{propriete}

\begin{aretenir}[Comment lire ces résultats]
On dit que \cle{l'exponentielle l'emporte sur toute puissance de $x$} :
\begin{itemize}
\item en $+\infty$ : $\mathrm{e}^x$ tend vers l'infini « plus vite » que n'importe quel polynôme, donc dans un quotient $\dfrac{\mathrm{e}^x}{x^n}$, c'est le numérateur qui gagne ;
\item en $-\infty$ : $\mathrm{e}^x$ tend vers 0 « plus vite » que n'importe quel polynôme ne tend vers l'infini, donc dans un produit $x^n \mathrm{e}^x$, c'est l'exponentielle qui gagne et le produit tend vers $0$.
\end{itemize}
\end{aretenir}

\begin{popfigure}
\begin{center}
\begin{tikzpicture}
\begin{axis}[
    width=0.85\linewidth, height=7cm,
    axis lines=middle,
    xlabel={$x$}, ylabel={$y$},
    xmin=-0.3, xmax=6.2, ymin=-2, ymax=60,
    xtick={1,2,3,4,5,6}, ytick={10,20,30,40,50},
    tick label style={font=\small},
    label style={font=\small},
    legend style={font=\small, at={(0.03,0.97)}, anchor=north west}
]
\addplot[popcurve,domain=0:4.15,samples=200]{exp(x)};
\addlegendentry{$y=\mathrm{e}^x$}
\addplot[poporange,thick,domain=0:6,samples=100]{x^2};
\addlegendentry{$y=x^2$}
\addplot[popgreen,thick,domain=0:6,samples=2]{x};
\addlegendentry{$y=x$}
\end{axis}
\end{tikzpicture}
\end{center}
\legende{Croissances comparées : dès que $x$ grandit, $\mathrm{e}^x$ dépasse $x^2$ (et donc $x$) de façon définitive. L'écart ne cesse ensuite de se creuser.}
\end{popfigure}

La figure le montre sans appel : sur $[0\,;\,6]$, la parabole $y = x^2$ semble d'abord rivaliser avec l'exponentielle, mais à partir d'environ $x = 4$, la courbe de $\exp$ s'envole et laisse toutes les puissances de $x$ loin derrière. C'est pourquoi on parle de \emph{croissance exponentielle} dans le langage courant pour désigner une croissance qui devient irrattrapable.

\courssub{Application : lever les indéterminations}

Les croissances comparées sont précisément l'outil qui permet de lever les formes indéterminées « $+\infty - \infty$ », « $\dfrac{\infty}{\infty}$ » ou « $0 \times \infty$ » mêlant polynômes et exponentielles. La stratégie reine : \textbf{factoriser par le terme dominant}, qui est toujours l'exponentielle.

\begin{methode}[Lever une indétermination avec $\mathrm{e}^x$]
\begin{enumerate}
\item Identifier la forme indéterminée (par exemple $+\infty - \infty$).
\item \textbf{Factoriser par $\mathrm{e}^x$} (le terme dominant) dans l'expression.
\item Dans le facteur restant, faire apparaître des termes du type $x^n \mathrm{e}^{-x}$ ou $\dfrac{x^n}{\mathrm{e}^x}$, qui tendent tous vers $0$ en $+\infty$ d'après les croissances comparées.
\item Conclure par les opérations sur les limites.
\end{enumerate}
\end{methode}

\begin{exoresolu}[Limites de fonctions polynôme $\times$ exponentielle et quotient]
\textbf{Énoncé.} Déterminer les limites suivantes :
\[ \text{a) } \lim_{x \to +\infty} \left( \mathrm{e}^x - x^2 \right) \qquad ; \qquad \text{b) } \lim_{x \to -\infty} (x+1)\,\mathrm{e}^{x} \qquad ; \qquad \text{c) } \lim_{x \to +\infty} \frac{\mathrm{e}^x}{x + 3}. \]
\tcblower
\textbf{Solution détaillée.}

\textbf{a)} Forme indéterminée « $+\infty - \infty$ ». On factorise par le terme dominant $\mathrm{e}^x$ :
\[ \mathrm{e}^x - x^2 = \mathrm{e}^x \left( 1 - \frac{x^2}{\mathrm{e}^x} \right) = \mathrm{e}^x \left( 1 - x^2 \mathrm{e}^{-x} \right). \]
Or $\displaystyle\lim_{x \to +\infty} \frac{x^2}{\mathrm{e}^x} = 0$ (croissances comparées : l'exponentielle l'emporte sur $x^2$). Donc la parenthèse tend vers $1 - 0 = 1$, et par produit :
\[ \lim_{x \to +\infty} \left( \mathrm{e}^x - x^2 \right) = +\infty. \]

\textbf{b)} Forme indéterminée « $-\infty \times 0$ ». On développe : $(x+1)\mathrm{e}^x = x\,\mathrm{e}^x + \mathrm{e}^x$.

D'après les croissances comparées, $\displaystyle\lim_{x \to -\infty} x\,\mathrm{e}^x = 0$, et on sait que $\displaystyle\lim_{x \to -\infty} \mathrm{e}^x = 0$. Par somme :
\[ \lim_{x \to -\infty} (x+1)\,\mathrm{e}^x = 0 + 0 = 0. \]

\textbf{c)} Forme indéterminée « $\dfrac{+\infty}{+\infty}$ ». On écrit :
\[ \frac{\mathrm{e}^x}{x+3} = \frac{\mathrm{e}^x}{x} \times \frac{x}{x+3}. \]
D'une part $\displaystyle\lim_{x \to +\infty} \frac{\mathrm{e}^x}{x} = +\infty$ (croissances comparées). D'autre part $\dfrac{x}{x+3} = \dfrac{1}{1 + \frac{3}{x}} \xrightarrow[x \to +\infty]{} 1$. Par produit :
\[ \lim_{x \to +\infty} \frac{\mathrm{e}^x}{x+3} = +\infty. \]
\end{exoresolu}

\begin{attention}[Pièges de rédaction à éviter]
\begin{itemize}
\item \textbf{Ne jamais écrire} « $\displaystyle\lim_{x \to +\infty}(\mathrm{e}^x - x^2) = +\infty - \infty = 0$ » : une forme indéterminée ne se « calcule » pas directement, elle se \emph{lève} par factorisation. Écris toujours l'étape de factorisation.
\item \textbf{Ne pas confondre} $\mathrm{e}^{-x}$ et $-\mathrm{e}^x$ : $\mathrm{e}^{-x} \to 0$ en $+\infty$, tandis que $-\mathrm{e}^x \to -\infty$.
\item \textbf{Vérifier le sens de la borne} avant d'appliquer les croissances comparées : $x\,\mathrm{e}^x \to 0$ en $-\infty$, mais $x\,\mathrm{e}^x \to +\infty$ en $+\infty$ (produit de deux quantités qui tendent vers $+\infty$).
\item Dans un quotient $\dfrac{\mathrm{e}^x}{P(x)}$ avec $P$ polynôme, c'est \textbf{toujours} l'exponentielle qui gagne en $+\infty$ : la limite est $+\infty$ (au signe près du coefficient dominant de $P$).
\end{itemize}
\end{attention}

\begin{exercice}[À toi de jouer]
Déterminer les limites suivantes, en justifiant chaque étape :
\[ \text{a) } \lim_{x \to +\infty} \mathrm{e}^{-x^2 + 4} \qquad ; \qquad \text{b) } \lim_{x \to +\infty} \left( \mathrm{e}^x - 3x \right) \qquad ; \qquad \text{c) } \lim_{x \to -\infty} \left( x^2 + x \right) \mathrm{e}^x. \]
\end{exercice}

\begin{corrige}[Exercice : limites avec l'exponentielle]
\textbf{a)} L'exposant $u(x) = -x^2 + 4$ est un trinôme de coefficient dominant $-1 < 0$, donc $\displaystyle\lim_{x \to +\infty}(-x^2 + 4) = -\infty$. Par composition : $\displaystyle\lim_{x \to +\infty} \mathrm{e}^{-x^2+4} = 0$.

\textbf{b)} Forme « $+\infty - \infty$ ». Factorisation par $\mathrm{e}^x$ : $\mathrm{e}^x - 3x = \mathrm{e}^x\left(1 - \dfrac{3x}{\mathrm{e}^x}\right)$. Par croissances comparées, $\dfrac{3x}{\mathrm{e}^x} \to 0$ en $+\infty$, donc la parenthèse tend vers $1$ et $\displaystyle\lim_{x \to +\infty}(\mathrm{e}^x - 3x) = +\infty$.

\textbf{c)} On développe : $(x^2 + x)\mathrm{e}^x = x^2 \mathrm{e}^x + x\,\mathrm{e}^x$. Par croissances comparées en $-\infty$, $x^2 \mathrm{e}^x \to 0$ et $x\,\mathrm{e}^x \to 0$. Par somme : $\displaystyle\lim_{x \to -\infty}(x^2 + x)\mathrm{e}^x = 0$.
\end{corrige}

Tu maîtrises désormais le comportement de l'exponentielle aux bornes de $\mathbb{R}$, la composition des limites avec un exposant quelconque, et l'arme redoutable des croissances comparées pour lever les indéterminations. Il ne reste plus qu'à donner à cette fonction son visage complet : dans la section suivante, nous étudierons sa \emph{dérivée}, son \emph{sens de variation} et nous tracerons rigoureusement sa courbe représentative — tu verras que tout ce que nous venons d'établir sur les limites y trouvera immédiatement sa place.

\coursec{Dérivation, variations et courbe représentative}

Dans la section précédente, nous avons découvert le comportement de la fonction exponentielle aux bornes de son domaine : elle tend vers $0$ en $-\infty$ et explose vers $+\infty$ en $+\infty$, plus vite que n'importe quelle puissance de $x$. Mais entre ces deux extrêmes, comment se comporte-t-elle exactement ? Est-elle croissante ? Où passe-t-elle ? Quelle est l'allure de sa courbe ? Pour répondre à toutes ces questions, l'outil roi est la \cle{dérivation}. Et c'est ici que l'exponentielle va révéler son secret le plus étonnant : c'est une fonction qui est \emph{sa propre dérivée}. Cette propriété, unique en son genre, explique pourquoi l'exponentielle apparaît partout dans la nature : dès qu'une quantité croît proportionnellement à elle-même (une population de bactéries, un capital placé à intérêts composés continus), c'est elle qui se cache derrière.

\courssub{Dérivabilité de la fonction exponentielle}

\textbf{Intuition.} La fonction exponentielle a été définie comme la réciproque du logarithme népérien. Or nous savons dériver $\ln$ : $\ln'(x) = \dfrac{1}{x}$ sur $]0\,;+\infty[$. Il existe un théorème général qui permet de dériver une fonction réciproque à partir de la dérivée de la fonction de départ. Appliquons-le.

\begin{propriete}[Dérivabilité de exp — démonstration exigible au Bac]
La fonction exponentielle est \cle{dérivable} sur $\mathbb{R}$ et, pour tout réel $x$ :
\[ \exp'(x) = \mathrm{e}^x. \]
Autrement dit, l'exponentielle est \textbf{égale à sa propre dérivée}.
\end{propriete}

\begin{experience}[Démonstration guidée : dérivée d'une fonction réciproque]
Rappelons le théorème de dérivation d'une bijection réciproque : si $f$ est une bijection dérivable en $g(x)$ avec $f'(g(x)) \neq 0$, alors sa réciproque $g$ est dérivable en $x$ et
\[ g'(x) = \frac{1}{f'\big(g(x)\big)}. \]
\textbf{Étape 1.} Ici, $f = \ln$ (bijection de $]0\,;+\infty[$ sur $\mathbb{R}$) et $g = \exp$ (bijection de $\mathbb{R}$ sur $]0\,;+\infty[$). On a bien $\ln(\mathrm{e}^x) = x$ pour tout $x \in \mathbb{R}$ et $\mathrm{e}^{\ln y} = y$ pour tout $y > 0$.

\textbf{Étape 2.} On sait que $\ln'(y) = \dfrac{1}{y}$ pour tout $y > 0$. Donc, pour tout réel $x$ :
\[ \ln'\big(\mathrm{e}^x\big) = \frac{1}{\mathrm{e}^x}. \]
Ce nombre est strictement positif (donc non nul), le théorème s'applique : $\exp$ est dérivable en tout point $x$ de $\mathbb{R}$.

\textbf{Étape 3.} On calcule :
\[ \exp'(x) = \frac{1}{\ln'\big(\mathrm{e}^x\big)} = \frac{1}{\dfrac{1}{\mathrm{e}^x}} = \mathrm{e}^x. \]
Conclusion : $\exp' = \exp$ sur $\mathbb{R}$. $\blacksquare$
\end{experience}

\begin{savaistu}[L'unique fonction égale à sa dérivée]
On peut démontrer que les seules fonctions dérivables sur $\mathbb{R}$ vérifiant $f' = f$ sont les fonctions de la forme $f(x) = k\,\mathrm{e}^x$ avec $k \in \mathbb{R}$. Si l'on impose de plus $f(0) = 1$, alors $k = 1$ : l'exponentielle est donc \textbf{l'unique} fonction dérivable sur $\mathbb{R}$ telle que $f' = f$ et $f(0) = 1$. C'est cette propriété qui la rend indispensable en physique (radioactivité), en biologie (croissance des populations) et en économie (intérêts composés en continu).
\end{savaistu}

\begin{propriete}[Sens de variation de exp]
Pour tout réel $x$, $\mathrm{e}^x > 0$. Comme $\exp'(x) = \mathrm{e}^x > 0$, la fonction exponentielle est \cle{strictement croissante} sur $\mathbb{R}$.
\end{propriete}

Cela confirme rigoureusement ce que nous avions utilisé dans les sections précédentes : $\mathrm{e}^a < \mathrm{e}^b \iff a < b$.

\courssub{Dérivée d'une composée : $(\mathrm{e}^u)' = u'\,\mathrm{e}^u$}

En pratique, on dérive rarement $\mathrm{e}^x$ seul : on rencontre $\mathrm{e}^{2x}$, $\mathrm{e}^{-x^2}$, $\mathrm{e}^{u(x)}$. On applique alors la formule de dérivation des composées.

\begin{propriete}[Dérivée de $\mathrm{e}^u$]
Si $u$ est une fonction dérivable sur un intervalle $I$, alors la fonction $x \mapsto \mathrm{e}^{u(x)}$ est dérivable sur $I$ et :
\[ \big(\mathrm{e}^{u}\big)' = u' \times \mathrm{e}^{u}. \]
\end{propriete}

\begin{exemplebox}[Calculs de dérivées]
\begin{itemize}
\item $f(x) = \mathrm{e}^{2x}$ : avec $u(x) = 2x$, $u'(x) = 2$, donc $f'(x) = 2\mathrm{e}^{2x}$.
\item $g(x) = \mathrm{e}^{-x}$ : $g'(x) = -\mathrm{e}^{-x}$.
\item $h(x) = \mathrm{e}^{x^2 - 3x}$ : $h'(x) = (2x - 3)\mathrm{e}^{x^2 - 3x}$.
\item $k(x) = x\,\mathrm{e}^{x}$ (produit !) : $k'(x) = 1 \times \mathrm{e}^x + x \times \mathrm{e}^x = (x+1)\mathrm{e}^x$.
\end{itemize}
\end{exemplebox}

\begin{attention}[Le piège classique : oublier le facteur $u'$]
L'erreur la plus fréquente au Baccalauréat est d'écrire $\big(\mathrm{e}^{2x}\big)' = \mathrm{e}^{2x}$. C'est \textbf{faux} : il faut multiplier par la dérivée de l'exposant. Retiens le réflexe :
\[ \big(\mathrm{e}^{u}\big)' = \underbrace{u'}_{\text{ne pas oublier !}} \times \mathrm{e}^{u}. \]
Ainsi $\big(\mathrm{e}^{2x}\big)' = 2\mathrm{e}^{2x}$ et $\big(\mathrm{e}^{-3x+1}\big)' = -3\mathrm{e}^{-3x+1}$. Vérifie systématiquement ce facteur avant de poursuivre un calcul.
\end{attention}

\courssub{Tableau de variations et courbe représentative de exp}

\textbf{Valeurs remarquables.} On connaît déjà $\mathrm{e}^0 = 1$ (car $\ln 1 = 0$) et $\mathrm{e}^1 = \mathrm{e} \approx 2{,}718$. La courbe passe donc par les points $(0\,;1)$ et $(1\,;\mathrm{e})$.

\textbf{Tangente en 0.} L'équation de la tangente à la courbe de $\exp$ au point d'abscisse $0$ est :
\[ y = \exp'(0)(x - 0) + \exp(0) = \mathrm{e}^0 \cdot x + 1 = x + 1. \]
La tangente à la courbe de $\exp$ en $(0\,;1)$ a donc pour équation $\boxed{y = x + 1}$.

Réunissons tout dans un tableau de variations complet :

\begin{center}
\begin{tabular}{|c|ccccc|}
\hline
\rowcolor{popblueL}
$x$ & $-\infty$ & & $0$ & & $+\infty$ \\
\hline
$\exp'(x) = \mathrm{e}^x$ & & $+$ & $+$ & $+$ & \\
\hline
$\exp(x)$ & $0$ & $\nearrow$ & $1$ & $\nearrow$ & $+\infty$ \\
\hline
\end{tabular}
\end{center}

La droite d'équation $y = 0$ (l'axe des abscisses) est \cle{asymptote horizontale} à la courbe en $-\infty$, car $\lim\limits_{x \to -\infty} \mathrm{e}^x = 0$.

\begin{popfigure}
\begin{center}
\begin{tikzpicture}
\begin{axis}[
    width=0.85\linewidth, height=7cm,
    axis lines=middle,
    xlabel={$x$}, ylabel={$y$},
    xmin=-3.2, xmax=2.6, ymin=-0.8, ymax=5,
    xtick={-2,-1,1,2}, ytick={1,2,3,4},
    grid=both, grid style={popline!40},
    tick label style={font=\small},
]
\addplot[popcurve,domain=-3:1.6,samples=200]{exp(x)};
\addplot[poporange,thick,domain=-2.5:2.2,samples=2]{x+1};
\addplot[popmuted,dashed,domain=-3:2.4,samples=2]{0};
\addplot[only marks,mark=*,popdark] coordinates {(0,1) (1,2.718)};
\node[popblue,font=\small] at (axis cs:-1.4,2.6) {$y = \mathrm{e}^x$};
\node[poporange,font=\small] at (axis cs:1.9,3.4) {$y = x+1$};
\node[popmuted,font=\small] at (axis cs:2.2,0.25) {asymptote $y=0$};
\node[popdark,font=\small,anchor=west] at (axis cs:0.05,0.55) {$(0\,;1)$};
\node[popdark,font=\small,anchor=south] at (axis cs:1,2.75) {$(1\,;\mathrm{e})$};
\end{axis}
\end{tikzpicture}
\end{center}
\legende{Courbe de la fonction exponentielle, sa tangente $y = x+1$ au point $(0\,;1)$ et son asymptote horizontale $y = 0$ en $-\infty$.}
\end{popfigure}

\begin{propriete}[Complément Bac : inégalité fondamentale $\mathrm{e}^x \geq x + 1$]
Pour tout réel $x$ : $\mathrm{e}^x \geq x + 1$, avec égalité si et seulement si $x = 0$. Graphiquement : la courbe de $\exp$ est \textbf{toujours au-dessus} de sa tangente en $0$.
\end{propriete}

\begin{experience}[Démonstration guidée de l'inégalité]
Posons $\varphi(x) = \mathrm{e}^x - x - 1$ pour tout réel $x$, et étudions cette fonction.

\textbf{Étape 1.} $\varphi$ est dérivable sur $\mathbb{R}$ et $\varphi'(x) = \mathrm{e}^x - 1$.

\textbf{Étape 2.} Signe de $\varphi'$ : comme $\exp$ est strictement croissante, $\mathrm{e}^x - 1 > 0 \iff \mathrm{e}^x > \mathrm{e}^0 \iff x > 0$. Donc $\varphi' < 0$ sur $]-\infty\,;0[$ et $\varphi' > 0$ sur $]0\,;+\infty[$.

\textbf{Étape 3.} $\varphi$ est donc strictement décroissante sur $]-\infty\,;0]$ puis strictement croissante sur $[0\,;+\infty[$ : elle admet un \textbf{minimum} en $x = 0$, qui vaut $\varphi(0) = \mathrm{e}^0 - 0 - 1 = 0$.

\textbf{Étape 4.} Pour tout réel $x$, $\varphi(x) \geq \varphi(0) = 0$, c'est-à-dire $\mathrm{e}^x - x - 1 \geq 0$, donc $\mathrm{e}^x \geq x + 1$. L'égalité n'a lieu qu'en $x = 0$. $\blacksquare$
\end{experience}

\begin{savaistu}[Pourquoi cette inégalité est si utile]
L'inégalité $\mathrm{e}^x \geq 1 + x$ est l'un des outils préférés des correcteurs du Baccalauréat : elle permet d'encadrer des exponentielles par des expressions affines, de prouver des convergences de suites, et elle se généralise en $\ln(1+x) \leq x$ (en posant $x = \ln(1+t)$). Retiens-la comme un réflexe !
\end{savaistu}

\courssub{Étude complète de fonctions du type $f(x) = (ax+b)\mathrm{e}^x$ ou $f(x) = \mathrm{e}^{ax+b}$}

\begin{methode}[Étudier $f(x) = (ax+b)\mathrm{e}^x$]
\begin{enumerate}
\item \textbf{Domaine} : $f$ est définie et dérivable sur $\mathbb{R}$ (produit d'un polynôme et de $\exp$).
\item \textbf{Limites} :
      \begin{itemize}
      \item En $+\infty$ : forme directe. $\lim\limits_{x \to +\infty} f(x) = +\infty$ si $a > 0$, $-\infty$ si $a < 0$.
      \item En $-\infty$ : forme indéterminée que l'on lève par croissances comparées : $\lim\limits_{x \to -\infty} x\,\mathrm{e}^x = 0$ et $\lim\limits_{x \to -\infty} \mathrm{e}^x = 0$, donc $\lim\limits_{x \to -\infty} f(x) = 0$ (asymptote horizontale $y = 0$).
      \end{itemize}
\item \textbf{Dérivée} : on dérive un \emph{produit} :
\[ f'(x) = a\,\mathrm{e}^x + (ax+b)\mathrm{e}^x = (ax + a + b)\,\mathrm{e}^x. \]
\item \textbf{Signe de $f'$} : comme $\mathrm{e}^x > 0$ pour tout $x$, le signe de $f'(x)$ est celui du \textbf{facteur affine} $ax + a + b$.
\item \textbf{Tableau de variations} et extremum en $x = -\dfrac{a+b}{a}$ (si $a \neq 0$).
\item \textbf{Courbe} : placer l'extremum, l'intersection avec l'axe des ordonnées $f(0) = b$, l'asymptote.
\end{enumerate}
\end{methode}

\begin{exemplebox}[Exemple chiffré : $f(x) = (x-2)\mathrm{e}^x$]
\textbf{Dérivée} : $f'(x) = 1 \times \mathrm{e}^x + (x-2)\mathrm{e}^x = (x - 1)\mathrm{e}^x$.

\textbf{Signe} : $\mathrm{e}^x > 0$, donc $f'(x)$ a le signe de $x - 1$ : négatif sur $]-\infty\,;1[$, positif sur $]1\,;+\infty[$.

\textbf{Variations} : $f$ est décroissante sur $]-\infty\,;1]$, croissante sur $[1\,;+\infty[$, avec un \textbf{minimum} $f(1) = (1-2)\mathrm{e}^1 = -\mathrm{e} \approx -2{,}72$.

\textbf{Limites} : $\lim\limits_{x \to -\infty} f(x) = 0$ (asymptote $y = 0$) et $\lim\limits_{x \to +\infty} f(x) = +\infty$.
\end{exemplebox}

\begin{exoresolu}[Étude complète de $f(x) = (x-1)\mathrm{e}^x$]
Soit $f$ la fonction définie sur $\mathbb{R}$ par $f(x) = (x-1)\mathrm{e}^x$.
\begin{enumerate}
\item Calculer les limites de $f$ en $-\infty$ et $+\infty$. Interpréter graphiquement.
\item Calculer $f'(x)$ et dresser le tableau de variations de $f$.
\item Déterminer l'équation de la tangente à la courbe au point d'abscisse $0$.
\item Tracer l'allure de la courbe.
\end{enumerate}
\tcblower
\textbf{1. Limites.} En $+\infty$ : $\lim\limits_{x \to +\infty}(x-1) = +\infty$ et $\lim\limits_{x \to +\infty}\mathrm{e}^x = +\infty$, donc par produit $\lim\limits_{x \to +\infty} f(x) = +\infty$.

En $-\infty$ : on écrit $f(x) = x\mathrm{e}^x - \mathrm{e}^x$. Par croissances comparées, $\lim\limits_{x \to -\infty} x\mathrm{e}^x = 0$ et $\lim\limits_{x \to -\infty} \mathrm{e}^x = 0$, donc $\lim\limits_{x \to -\infty} f(x) = 0$. \textbf{Interprétation} : la droite $y = 0$ est asymptote horizontale à la courbe en $-\infty$.

\textbf{2. Dérivée et variations.} $f$ est un produit : $f'(x) = \mathrm{e}^x + (x-1)\mathrm{e}^x = x\,\mathrm{e}^x$. Comme $\mathrm{e}^x > 0$, $f'(x)$ a le signe de $x$ : négatif sur $]-\infty\,;0[$, positif sur $]0\,;+\infty[$. De plus $f(0) = (0-1)\mathrm{e}^0 = -1$.

\begin{center}
\begin{tabular}{|c|ccccc|}
\hline
\rowcolor{popblueL}
$x$ & $-\infty$ & & $0$ & & $+\infty$ \\
\hline
$f'(x) = x\mathrm{e}^x$ & & $-$ & $0$ & $+$ & \\
\hline
$f(x)$ & $0$ & $\searrow$ & $-1$ & $\nearrow$ & $+\infty$ \\
\hline
\end{tabular}
\end{center}

$f$ admet donc un \textbf{minimum} en $0$ valant $-1$.

\textbf{3. Tangente en 0.} $f'(0) = 0 \times \mathrm{e}^0 = 0$ : la tangente au point $(0\,;-1)$ est \textbf{horizontale}, d'équation $y = -1$.

\textbf{4. Courbe.} La courbe part de l'asymptote $y = 0$ en $-\infty$ (en restant négative, car $x - 1 < 0$ pour $x < 1$), descend jusqu'au minimum $(0\,;-1)$ où la tangente est horizontale, remonte en coupant l'axe des abscisses en $x = 1$ (car $f(1) = 0$), puis s'envole vers $+\infty$.
\end{exoresolu}

\begin{popfigure}
\begin{center}
\begin{tikzpicture}
\begin{axis}[
    width=0.85\linewidth, height=7.5cm,
    axis lines=middle,
    xlabel={$x$}, ylabel={$y$},
    xmin=-4.2, xmax=2.6, ymin=-1.8, ymax=4,
    xtick={-4,-3,-2,-1,1,2}, ytick={-1,1,2,3},
    grid=both, grid style={popline!40},
    tick label style={font=\small},
]
\addplot[popcurve,domain=-4.2:1.85,samples=200]{(x-1)*exp(x)};
\addplot[poporange,thick,dashed,domain=-4:2.4,samples=2]{-1};
\addplot[only marks,mark=*,popdark] coordinates {(0,-1) (1,0)};
\node[popblue,font=\small] at (axis cs:-2.6,1.8) {$y = (x-1)\mathrm{e}^x$};
\node[popdark,font=\small,anchor=north west] at (axis cs:0.1,-1.05) {minimum $(0\,;-1)$};
\node[popdark,font=\small,anchor=south west] at (axis cs:1,0.05) {$(1\,;0)$};
\node[poporange,font=\small] at (axis cs:2.1,-1.35) {$y=-1$};
\end{axis}
\end{tikzpicture}
\end{center}
\legende{Courbe de $f(x) = (x-1)\mathrm{e}^x$ : asymptote $y = 0$ en $-\infty$, minimum $(0\,;-1)$ à tangente horizontale, passage par $(1\,;0)$.}
\end{popfigure}

\begin{exoresolu}[Étude rapide de $g(x) = \mathrm{e}^{2x-3}$]
Soit $g(x) = \mathrm{e}^{2x-3}$ définie sur $\mathbb{R}$. Déterminer ses limites, sa dérivée et ses variations.
\tcblower
\textbf{Limites} : $\lim\limits_{x \to -\infty}(2x-3) = -\infty$ donc $\lim\limits_{x \to -\infty} g(x) = 0$ (asymptote $y = 0$) ; $\lim\limits_{x \to +\infty}(2x-3) = +\infty$ donc $\lim\limits_{x \to +\infty} g(x) = +\infty$.

\textbf{Dérivée} : avec $u(x) = 2x - 3$, $u'(x) = 2$, donc $g'(x) = 2\mathrm{e}^{2x-3}$.

\textbf{Variations} : $g'(x) > 0$ pour tout $x$ (produit de $2$ et d'une exponentielle) : $g$ est \textbf{strictement croissante} sur $\mathbb{R}$. Sa courbe a la même allure que celle de $\exp$, simplement « compressée » horizontalement et décalée : elle passe par $\left(\tfrac{3}{2}\,;1\right)$ car $g\left(\tfrac{3}{2}\right) = \mathrm{e}^0 = 1$.
\end{exoresolu}

\begin{attention}[Rédaction au Bac : justifier le signe de la dérivée]
Quand tu obtiens une dérivée du type $f'(x) = (ax + a + b)\mathrm{e}^x$, n'oublie \textbf{jamais} d'écrire explicitement la phrase : « Pour tout réel $x$, $\mathrm{e}^x > 0$, donc $f'(x)$ a le signe de $ax + a + b$. » C'est cette justification qui vaut des points. De même, pour les limites en $-\infty$ de $(ax+b)\mathrm{e}^x$, cite la croissance comparée : $\lim\limits_{x \to -\infty} x\mathrm{e}^x = 0$.
\end{attention}

\begin{aretenir}[L'essentiel de la section]
\begin{itemize}
\item $\exp$ est dérivable sur $\mathbb{R}$ et $\exp' = \exp$ : c'est (à coefficient près) la seule fonction égale à sa dérivée.
\item $\big(\mathrm{e}^{u}\big)' = u'\,\mathrm{e}^{u}$ — ne jamais oublier le facteur $u'$.
\item $\mathrm{e}^x > 0$ et $\exp' > 0$ : $\exp$ est strictement croissante sur $\mathbb{R}$.
\item Tangente en $0$ : $y = x + 1$ ; inégalité fondamentale : $\mathrm{e}^x \geq x + 1$ pour tout réel $x$.
\item Pour $f(x) = (ax+b)\mathrm{e}^x$ : $f'(x) = (ax + a + b)\mathrm{e}^x$, dont le signe est celui du facteur affine.
\item Asymptote horizontale $y = 0$ en $-\infty$ ; divergence vers $+\infty$ en $+\infty$ (si coefficient dominant positif).
\end{itemize}
\end{aretenir}

\begin{exercice}
Soit $h$ la fonction définie sur $\mathbb{R}$ par $h(x) = (2x + 1)\mathrm{e}^{-x}$.
\begin{enumerate}
\item Calculer les limites de $h$ en $-\infty$ et $+\infty$.
\item Montrer que $h'(x) = (1 - 2x)\mathrm{e}^{-x}$ et dresser le tableau de variations de $h$.
\item Déterminer l'équation de la tangente à la courbe de $h$ au point d'abscisse $0$.
\end{enumerate}
\end{exercice}

\begin{corrige}[Exercice 1]
\textbf{1.} En $+\infty$ : $h(x) = \dfrac{2x+1}{\mathrm{e}^{x}}$. Par croissances comparées, $\lim\limits_{x \to +\infty}\dfrac{x}{\mathrm{e}^x} = 0$, donc $\lim\limits_{x \to +\infty} h(x) = 0$ : asymptote $y = 0$. En $-\infty$ : $\mathrm{e}^{-x} \to +\infty$ et $2x+1 \to -\infty$, donc $\lim\limits_{x \to -\infty} h(x) = -\infty$.

\textbf{2.} $h$ est un produit : $h'(x) = 2\mathrm{e}^{-x} + (2x+1)(-\mathrm{e}^{-x}) = (2 - 2x - 1)\mathrm{e}^{-x} = (1 - 2x)\mathrm{e}^{-x}$. Comme $\mathrm{e}^{-x} > 0$, $h'(x)$ a le signe de $1 - 2x$ : positif sur $\left]-\infty\,;\tfrac{1}{2}\right[$, négatif sur $\left]\tfrac{1}{2}\,;+\infty\right[$. $h$ est croissante puis décroissante, avec un maximum $h\left(\tfrac{1}{2}\right) = 2\mathrm{e}^{-1/2} = \dfrac{2}{\sqrt{\mathrm{e}}} \approx 1{,}21$.

\textbf{3.} $h(0) = 1$ et $h'(0) = 1$ : la tangente a pour équation $y = x + 1$.
\end{corrige}

\coursec{Primitives et modélisation par l'exponentielle}

Dans la section précédente, nous avons établi que la fonction exponentielle est sa propre dérivée : $(e^x)' = e^x$, et plus généralement que $(e^u)' = u'\,e^u$. Cette propriété, remarquable entre toutes, va maintenant être « lue à l'envers » : si dériver $e^u$ donne $u'\,e^u$, alors \emph{primitiver} $u'\,e^u$ doit redonner $e^u$. C'est exactement le principe du calcul des primitives : remonter le cours de la dérivation. Nous allons exploiter cette idée pour calculer des primitives, des intégrales et des aires, puis nous verrons pourquoi l'exponentielle est \emph{la} fonction des phénomènes d'évolution : croissance d'une population, décroissance radioactive, refroidissement d'un corps. C'est là tout l'enjeu de cette dernière section, très prisée au Baccalauréat.

\courssub{Primitives de la fonction exponentielle}

\textbf{Le point de départ : primitiver $e^x$.}\par
Puisque $\exp' = \exp$, la fonction exponentielle est sa propre primitive. On obtient immédiatement :
\[ \int e^x \, dx = e^x + C, \quad C \in \mathbb{R}. \]
Autrement dit, sur $\mathbb{R}$, les primitives de $x \mapsto e^x$ sont les fonctions $x \mapsto e^x + C$. C'est aussi simple que la dérivation : l'exponentielle est « stable » dans les deux sens.

\textbf{Primitiver une composée : la forme $u'\,e^u$.}\par
Rappelons la formule de dérivation d'une composée vue à la section 5 : si $u$ est dérivable sur un intervalle $I$, alors
\[ \left(e^u\right)' = u'\,e^u. \]
En lisant cette égalité de droite à gauche, on obtient le résultat fondamental de cette section.

\begin{propriete}[Primitives et exponentielle]
Soit $u$ une fonction dérivable sur un intervalle $I$.
\begin{itemize}
\item Une primitive sur $I$ de la fonction $u'\,e^{u}$ est la fonction $e^{u}$.
\item En particulier, pour $a \neq 0$ et $b$ réels, une primitive sur $\mathbb{R}$ de $x \mapsto e^{ax+b}$ est
\[ x \longmapsto \frac{1}{a}\,e^{ax+b}. \]
\end{itemize}
\end{propriete}

\begin{experience}[Démonstration guidée (exigible)]
Démontrons qu'une primitive de $e^{ax+b}$ est $\dfrac{1}{a}e^{ax+b}$ pour $a \neq 0$.
\begin{enumerate}
\item Posons $F(x) = \dfrac{1}{a}\,e^{ax+b}$. Il s'agit d'une composée avec $u(x) = ax + b$.
\item On dérive en utilisant $(e^u)' = u'e^u$ : ici $u'(x) = a$, donc
\[ F'(x) = \frac{1}{a} \times a\,e^{ax+b} = e^{ax+b}. \]
\item Conclusion : $F$ a bien pour dérivée $x \mapsto e^{ax+b}$, donc $F$ est une primitive de $e^{ax+b}$. $\square$
\end{enumerate}
Le facteur $\dfrac{1}{a}$ sert exactement à « compenser » le facteur $a$ qui apparaît quand on dérive la composée.
\end{experience}

\begin{methode}[Primitiver $k \times u'\,e^{u}$ : reconnaître, ajuster, vérifier]
Pour trouver une primitive d'une expression contenant une exponentielle :
\begin{enumerate}
\item \textbf{Repérer} l'exposant : c'est la fonction $u(x)$. Calculer mentalement $u'(x)$.
\item \textbf{Comparer} : le facteur devant l'exponentielle est-il égal à $u'$, à une constante près ?
\item \textbf{Ajuster} le coefficient : si l'expression est $k \times u'\,e^{u}$, une primitive est $k\,e^{u}$. Si le facteur est $\lambda \times u'$, la primitive est $\lambda\,e^{u}$.
\item \textbf{Vérifier toujours} en dérivant le résultat obtenu : on doit retomber sur l'expression de départ.
\end{enumerate}
\end{methode}

\begin{exemplebox}[Exemples chiffrés fondamentaux]
\textbf{a)} Cherchons une primitive de $f(x) = 6x\,e^{x^2}$.
L'exposant est $u(x) = x^2$, donc $u'(x) = 2x$. Or $6x = 3 \times 2x = 3\,u'(x)$. Ainsi $f = 3\,u'\,e^{u}$, et une primitive est
\[ F(x) = 3\,e^{x^2}. \]
Vérification : $F'(x) = 3 \times 2x\,e^{x^2} = 6x\,e^{x^2} = f(x)$. \checkmark

\textbf{b)} Cherchons une primitive de $g(x) = e^{-2x+1}$.
C'est la forme $e^{ax+b}$ avec $a = -2$ et $b = 1$. Une primitive est
\[ G(x) = \frac{1}{-2}\,e^{-2x+1} = -\frac{1}{2}\,e^{-2x+1}. \]
Vérification : $G'(x) = -\dfrac{1}{2}\times(-2)e^{-2x+1} = e^{-2x+1} = g(x)$. \checkmark

\textbf{c)} Cherchons une primitive de $h(x) = x\,e^{x^2}$.
Ici $u(x) = x^2$, $u'(x) = 2x$, et $x = \dfrac{1}{2}\times 2x$. Donc $h = \dfrac{1}{2}u'e^u$ et une primitive est $H(x) = \dfrac{1}{2}e^{x^2}$.
Vérification : $H'(x) = \dfrac{1}{2} \times 2x\,e^{x^2} = x\,e^{x^2} = h(x)$. \checkmark
\end{exemplebox}

\begin{attention}[L'erreur classique : oublier de diviser par $a$]
L'erreur la plus fréquente au Baccalauréat est d'écrire qu'une primitive de $e^{ax+b}$ est $e^{ax+b}$, par analogie avec $e^x$. C'est \textbf{faux} dès que $a \neq 1$ : en dérivant $e^{ax+b}$, on obtient $a\,e^{ax+b}$ et non $e^{ax+b}$. Il faut \textbf{diviser par $a$}. De même, on ne peut primitiver $e^{u}$ en $e^{u}$ que si le facteur $u'$ est présent (à une constante près) devant l'exponentielle : la fonction $x \mapsto e^{x^2}$ seule, sans le facteur $x$, n'admet \emph{pas} de primitive exprimable avec les fonctions usuelles. Réflexe de survie : \textbf{dérive toujours ta réponse pour vérifier}.
\end{attention}

\courssub{Intégrales et aires avec l'exponentielle}

Le lien avec la leçon sur l'intégration est direct : pour calculer $\displaystyle\int_{\alpha}^{\beta} f(x)\,dx$, on cherche une primitive $F$ de $f$, puis on applique
\[ \int_{\alpha}^{\beta} f(x)\,dx = \big[F(x)\big]_{\alpha}^{\beta} = F(\beta) - F(\alpha). \]

\begin{exoresolu}[Calcul d'intégrale et aire sous la courbe exponentielle]
\textbf{Énoncé.} 1) Calculer $I = \displaystyle\int_{0}^{1} e^{3x}\,dx$.
2) On considère la fonction $f(x) = 4x\,e^{x^2}$ sur $[0\,;1]$. Calculer l'aire $\mathcal{A}$, en unités d'aire, du domaine délimité par la courbe de $f$, l'axe des abscisses et les droites $x = 0$ et $x = 1$.
\tcblower
\textbf{Solution.}
1) Une primitive de $e^{3x}$ est $\dfrac{1}{3}e^{3x}$ (forme $e^{ax+b}$ avec $a = 3$). Donc
\begin{align*}
I &= \left[\frac{1}{3}e^{3x}\right]_{0}^{1} = \frac{1}{3}e^{3} - \frac{1}{3}e^{0} = \frac{e^3 - 1}{3}.
\end{align*}
Numériquement : $I \approx \dfrac{20{,}086 - 1}{3} \approx 6{,}36$.

2) Sur $[0\,;1]$, $f(x) = 4x\,e^{x^2} \geq 0$, donc l'aire cherchée est $\mathcal{A} = \displaystyle\int_{0}^{1} 4x\,e^{x^2}\,dx$. On reconnaît la forme $u'e^u$ avec $u(x) = x^2$, $u'(x) = 2x$ : or $4x = 2\times 2x$, donc une primitive est $F(x) = 2e^{x^2}$. Ainsi
\[ \mathcal{A} = \left[2e^{x^2}\right]_{0}^{1} = 2e^{1} - 2e^{0} = 2e - 2 = 2(e-1) \text{ u.a.} \]
soit $\mathcal{A} \approx 3{,}44$ unités d'aire.
\end{exoresolu}

\begin{exercice}[À toi de jouer]
1) Déterminer une primitive de chacune des fonctions suivantes :
\[ f(x) = e^{5x-2} \quad ; \quad g(x) = (2x+1)e^{x^2+x} \quad ; \quad h(x) = -3e^{-x}. \]
2) Calculer $\displaystyle\int_{0}^{\ln 2} e^{x}\,dx$ et interpréter le résultat comme une aire.
3) Calculer $\displaystyle\int_{-1}^{1} x\,e^{-x^2}\,dx$ (on pourra remarquer une symétrie).
\end{exercice}

\begin{corrige}[Exercice]
1) $F(x) = \dfrac{1}{5}e^{5x-2}$ ; $G(x) = e^{x^2+x}$ car $(x^2+x)' = 2x+1$ ; $H(x) = 3e^{-x}$ car $(3e^{-x})' = -3e^{-x}$.
2) $\displaystyle\int_{0}^{\ln 2} e^x\,dx = e^{\ln 2} - e^0 = 2 - 1 = 1$. C'est l'aire, en unités d'aire, du domaine sous la courbe de $\exp$ entre $0$ et $\ln 2$ : elle vaut exactement $1$ u.a.
3) La fonction $x \mapsto x e^{-x^2}$ est impaire (produit d'une fonction impaire et d'une fonction paire), et l'intervalle $[-1\,;1]$ est centré en $0$, donc l'intégrale vaut $0$. Vérification par le calcul : une primitive est $-\dfrac{1}{2}e^{-x^2}$, et $-\dfrac{1}{2}e^{-1} - \left(-\dfrac{1}{2}e^{-1}\right) = 0$.
\end{corrige}

\courssub{Modélisation : les phénomènes régis par $f' = kf$}

\textbf{L'intuition : une croissance proportionnelle à la quantité présente.}\par
Imaginons une plantation de cacao dans la région du Centre. Plus il y a de plants, plus le nombre de nouveaux rejets chaque année est grand : la \emph{vitesse} de croissance est proportionnelle à la quantité déjà présente. De même, une masse de matière radioactive se désintègre d'autant plus vite qu'il reste d'atomes non désintégrés. Dans les deux cas, si $f(t)$ désigne la quantité à l'instant $t$, le phénomène obéit à la loi
\[ f'(t) = k\,f(t), \]
où $k$ est une constante : $k > 0$ pour une croissance, $k < 0$ pour une décroissance. Quelles fonctions vérifient cette relation ?

\begin{propriete}[Complément Bac (admis) : les solutions de $f' = kf$]
Soit $k$ un réel. Les fonctions dérivables sur $\mathbb{R}$ vérifiant $f' = kf$ sont exactement les fonctions de la forme
\[ f(t) = C\,e^{kt}, \quad \text{où } C \in \mathbb{R}. \]
De plus, $C = f(0)$ : la constante $C$ est la \textbf{valeur initiale}. Ce résultat, admis ici, sera démontré dans la leçon sur les équations différentielles. C'est le fondement mathématique de tous les modèles d'évolution exponentielle.
\end{propriete}

\begin{experience}[Vérification (formatrice)]
Vérifions que $f(t) = C e^{kt}$ satisfait bien $f' = kf$.
\begin{enumerate}
\item $f$ est une composée : $f(t) = C\,e^{u(t)}$ avec $u(t) = kt$, donc $u'(t) = k$.
\item $f'(t) = C \times k\,e^{kt} = k \times \left(C e^{kt}\right) = k\,f(t)$. $\square$
\end{enumerate}
La réciproque (ce sont les \emph{seules} solutions) est admise à ce niveau.
\end{experience}

\begin{popfigure}
\begin{center}
\begin{tikzpicture}[scale=1, every node/.style={font=\small}]
\node[draw=popblue, fill=popblueL, rounded corners, thick, minimum width=4.2cm, minimum height=1.1cm] (eq) at (0,0) {\textbf{Loi d'évolution : } $f' = kf$};
\node[draw=popgreen, fill=popgreenL, rounded corners, thick, minimum width=4.6cm, minimum height=1.1cm] (sol) at (9,0) {\textbf{Solution : } $f(t) = C\,e^{kt}$};
\draw[-{Stealth}, very thick, popdark] (eq) -- (sol) node[midway, above]{on résout};
\draw[-{Stealth}, very thick, popdark] (sol.south) .. controls (5,-1.6) .. (eq.south) node[midway, below]{on vérifie : $(Ce^{kt})' = kCe^{kt}$};
\node[draw=poporange, fill=poporangeL, rounded corners, align=center] (c) at (0,-3.2) {$k > 0$ : \textbf{croissance}\\ bactéries, population,\\ capital à intérêts composés};
\node[draw=poppink, fill=poppinkL, rounded corners, align=center] (d) at (9,-3.2) {$k < 0$ : \textbf{décroissance}\\ radioactivité, refroidissement,\\ extinction d'un feu de forêt};
\draw[-{Stealth}, thick, poporange] (c.north) -- (0,-0.7);
\draw[-{Stealth}, thick, poppink] (d.north) -- (9,-0.7);
\end{tikzpicture}
\end{center}
\legende{La boucle fondamentale de la modélisation exponentielle : $f' = kf \Leftrightarrow f(t) = Ce^{kt}$.}
\end{popfigure}

\textbf{Exploiter un modèle : déterminer $C$ et $k$, calculer des temps caractéristiques.}\par
En pratique, on dispose de données expérimentales : la valeur initiale donne $C = f(0)$, et une seconde mesure permet de trouver $k$ en résolvant une équation avec le logarithme népérien. Deux temps caractéristiques reviennent constamment :
\begin{itemize}
\item le \cle{temps de doublement} $T$ (si $k > 0$) : c'est la durée pour que la quantité double, soit $e^{kT} = 2$, d'où $T = \dfrac{\ln 2}{k}$ ;
\item la \cle{demi-vie} $T$ (si $k < 0$) : c'est la durée pour que la quantité soit divisée par deux, soit $e^{kT} = \dfrac{1}{2}$, d'où $T = \dfrac{\ln 2}{-k} = -\dfrac{\ln 2}{k}$.
\end{itemize}
Remarque essentielle : ces temps \emph{ne dépendent pas de la valeur initiale} $C$, uniquement de $k$.

\begin{exemplebox}[Temps de doublement d'une population]
Une ville du Cameroun compte \num{10000} habitants et croît selon le modèle $P(t) = \num{10000}\,e^{0{,}03t}$ ($t$ en années). Cherchons son temps de doublement :
\begin{align*}
P(t) = 2 \times \num{10000} &\iff \num{10000}\,e^{0{,}03t} = \num{20000} \\
&\iff e^{0{,}03t} = 2 \\
&\iff 0{,}03t = \ln 2 \quad (\text{on applique } \ln) \\
&\iff t = \frac{\ln 2}{0{,}03} \approx \frac{0{,}693}{0{,}03} \approx 23{,}1 \text{ ans}.
\end{align*}
La population double environ tous les $23$ ans, quelle que soit la population de départ.
\end{exemplebox}

\begin{exoresolu}[Décroissance radioactive et demi-vie]
\textbf{Énoncé.} Un échantillon contient $N_0 = 8 \times 10^{20}$ noyaux d'un isotope radioactif. Le nombre de noyaux restants à l'instant $t$ (en heures) est $N(t) = N_0\,e^{-0{,}12t}$.
1) Vérifier que $N' = -0{,}12\,N$.
2) Déterminer la demi-vie de cet isotope, arrondie à $0{,}1$ heure.
3) Combien reste-t-il de noyaux au bout de deux demi-vies ?
\tcblower
\textbf{Solution.}
1) $N'(t) = N_0 \times (-0{,}12)e^{-0{,}12t} = -0{,}12 \times N_0 e^{-0{,}12t} = -0{,}12\,N(t)$. \checkmark
2) La demi-vie $T$ vérifie $N(T) = \dfrac{N_0}{2}$ :
\begin{align*}
N_0\,e^{-0{,}12T} = \frac{N_0}{2} &\iff e^{-0{,}12T} = \frac{1}{2} \\
&\iff -0{,}12T = \ln\frac{1}{2} = -\ln 2 \\
&\iff T = \frac{\ln 2}{0{,}12} \approx \frac{0{,}693}{0{,}12} \approx 5{,}8 \text{ heures}.
\end{align*}
3) Après une demi-vie, il reste $\dfrac{N_0}{2}$ noyaux ; après une seconde demi-vie, la moitié de cette moitié, soit $\dfrac{N_0}{4} = 2 \times 10^{20}$ noyaux. On vérifie : $N(2T) = N_0 e^{-0{,}12 \times 2T} = N_0 \left(e^{-0{,}12T}\right)^2 = N_0 \times \left(\dfrac{1}{2}\right)^2 = \dfrac{N_0}{4}$.
\end{exoresolu}

\begin{popfigure}
\begin{center}
\begin{tikzpicture}
\begin{axis}[
  width=\linewidth, height=7.2cm,
  xlabel={$t$ (temps)}, ylabel={$N(t)$},
  xmin=0, xmax=40, ymin=0, ymax=1050,
  grid=both, grid style={popline!40},
  legend style={at={(0.97,0.95)}, anchor=north east, font=\small},
  axis lines=left, thick
]
\addplot[popcurve, popblue, domain=0:20, samples=100]{100*exp(0.035*x)};
\addlegendentry{$N(t) = 100\,e^{0{,}035t}$ (croissance)}
\addplot[popcurve, poppink, domain=0:40, samples=100]{1000*exp(-0.12*x)};
\addlegendentry{$N(t) = 1000\,e^{-0{,}12t}$ (décroissance)}
\addplot[popmuted, dashed, domain=0:5.8, samples=2]{500};
\draw[popmuted, dashed] (axis cs:5.8,0) -- (axis cs:5.8,500);
\node[popdark, font=\small, anchor=south] at (axis cs:5.8,510) {demi-vie $T \approx 5{,}8$};
\addplot[mark=*, popdark, only marks] coordinates {(5.8,500)};
\end{axis}
\end{tikzpicture}
\end{center}
\legende{Croissance exponentielle (en bleu) et décroissance radioactive (en rose) : la demi-vie est le temps au bout duquel la quantité initiale est divisée par deux.}
\end{popfigure}

\begin{savaistu}[La datation au carbone 14]
Le carbone 14 est un isotope radioactif de demi-vie environ \num{5730} ans, présent dans tout organisme vivant. À la mort de l'organisme, il n'est plus renouvelé et décroît selon $N(t) = N_0 e^{kt}$ avec $k = -\dfrac{\ln 2}{5730}$. En mesurant la proportion de carbone 14 restant dans un ossement ou un objet en bois — par exemple lors de fouilles archéologiques sur les sites anciens du Cameroun —, les scientifiques résolvent $\dfrac{N(t)}{N_0} = e^{kt}$ pour en déduire l'âge $t$ de l'échantillon : $t = \dfrac{1}{k}\ln\dfrac{N(t)}{N_0}$. C'est la méthode qui a valu le prix Nobel de chimie à Willard Libby en 1960.
\end{savaistu}

\begin{attention}[Pièges de rédaction dans les problèmes de modélisation]
\begin{itemize}
\item Ne pas confondre $C$ et $k$ : $C = f(0)$ est la \textbf{valeur initiale}, $k$ est le \textbf{taux} (relatif) d'évolution. On détermine d'abord $C$, puis $k$ grâce à une donnée supplémentaire.
\item Pour isoler $t$ dans $e^{kt} = A$, on applique $\ln$ : $kt = \ln A$, donc $t = \dfrac{\ln A}{k}$. Attention au signe de $k$ !
\item Le temps de doublement et la demi-vie sont indépendants de $C$ : inutile de traîner $N_0$ dans tous les calculs, il se simplifie.
\item Toujours préciser l'\textbf{unité} de $t$ (années, heures, jours...) cohérente avec celle de $k$.
\end{itemize}
\end{attention}

\begin{exercice}[Modèle de refroidissement]
Une tasse de café chaud, initialement à $90\,^{\circ}$C, est posée dans une salle à $25\,^{\circ}$C. On admet que l'écart de température $\theta(t) = T(t) - 25$ vérifie $\theta' = -0{,}08\,\theta$ ($t$ en minutes).
1) Exprimer $\theta(t)$ puis $T(t)$.
2) Au bout de combien de minutes le café atteint-il $50\,^{\circ}$C ? Arrondir à la minute.
\end{exercice}

\begin{corrige}[Exercice]
1) $\theta$ vérifie $\theta' = -0{,}08\,\theta$, donc $\theta(t) = C e^{-0{,}08t}$ avec $C = \theta(0) = 90 - 25 = 65$. Ainsi $\theta(t) = 65e^{-0{,}08t}$ et $T(t) = 25 + 65e^{-0{,}08t}$.
2) $T(t) = 50 \iff 65e^{-0{,}08t} = 25 \iff e^{-0{,}08t} = \dfrac{25}{65} = \dfrac{5}{13} \iff -0{,}08t = \ln\dfrac{5}{13} \iff t = \dfrac{\ln(13/5)}{0{,}08} \approx \dfrac{0{,}9555}{0{,}08} \approx 12$ minutes.
\end{corrige}

\begin{aretenir}[L'essentiel de la section]
\begin{itemize}
\item Une primitive de $e^x$ est $e^x$ ; une primitive de $u'\,e^{u}$ est $e^{u}$ ; une primitive de $e^{ax+b}$ ($a \neq 0$) est $\dfrac{1}{a}e^{ax+b}$.
\item Réflexe : repérer $u$, comparer avec $u'$, ajuster le coefficient, \textbf{vérifier en dérivant}.
\item Les phénomènes dont la vitesse d'évolution est proportionnelle à la quantité présente obéissent à $f' = kf$, dont les solutions sont $f(t) = C e^{kt}$ avec $C = f(0)$.
\item Temps de doublement ($k > 0$) : $T = \dfrac{\ln 2}{k}$ ; demi-vie ($k < 0$) : $T = \dfrac{\ln 2}{-k}$. On les obtient en résolvant une équation avec $\ln$.
\end{itemize}
\end{aretenir}

\newpage

\coursec{Exponentielles de base $a$ et fonctions puissances -- Généralisation}

Tu sais maintenant manipuler la fonction exponentielle népérienne $x \mapsto e^x$, réciproque du logarithme népérien. Mais dans la vie courante, on rencontre souvent des exponentielles « en base 10 » ou « en base 2 » : un capital qui double à chaque période s'écrit $2^n$, l'échelle de Richter ou le pH utilisent des puissances de $10$. De même, les fonctions $x \mapsto \sqrt{x}$, $x \mapsto x^{1{,}5}$ ou $x \mapsto x^{-2}$ ne sont ni des polynômes classiques ni des exponentielles. Ce complément montre que toutes ces fonctions se ramènent à un unique outil : l'écriture $e^{(\cdots)}$.

\courssub{Fonction exponentielle de base $a$}

\begin{definition}[Exponentielle de base $a$]
Soit $a$ un réel strictement positif. Pour tout réel $x$, on pose :
\[ a^x = e^{x \ln a}. \]
La fonction $x \mapsto a^x$, définie sur $\mathbb{R}$, est appelée \cle{fonction exponentielle de base $a$}.
\end{definition}

Cette définition prolonge naturellement les puissances entières : pour $n \in \mathbb{N}$, $e^{n \ln a} = \left(e^{\ln a}\right)^n = a^n$, on retrouve bien la puissance habituelle. Elle a maintenant un sens pour \emph{tout} exposant réel : $2^{\sqrt{2}} = e^{\sqrt{2}\,\ln 2} \approx 2{,}665$.

\begin{propriete}[Propriétés algébriques]
Pour tous réels $x$ et $y$, et tous réels $a, b > 0$ :
\[ a^{x+y} = a^x \times a^y \ ; \quad a^{-x} = \frac{1}{a^x} \ ; \quad (a^x)^y = a^{xy} \ ; \quad a^x b^x = (ab)^x \ ; \quad a^0 = 1. \]
Ces propriétés découlent immédiatement de celles de l'exponentielle népérienne.
\end{propriete}

\begin{propriete}[Dérivée, variations et limites]
La fonction $f : x \mapsto a^x$ est dérivable sur $\mathbb{R}$ et, comme $f(x) = e^{x\ln a}$ est de la forme $e^u$ avec $u(x) = x\ln a$ :
\[ f'(x) = (\ln a)\, a^x. \]
Comme $a^x > 0$, le signe de $f'$ est celui de $\ln a$ :
\begin{itemize}
\item si $a > 1$, alors $\ln a > 0$ : $f$ est \textbf{strictement croissante} sur $\mathbb{R}$ ;
\item si $0 < a < 1$, alors $\ln a < 0$ : $f$ est \textbf{strictement décroissante} sur $\mathbb{R}$ ;
\item si $a = 1$ : $f$ est constante, égale à $1$.
\end{itemize}
Limites aux bornes :
\begin{itemize}
\item Si $a > 1$ : $\lim_{x \to +\infty} a^x = +\infty$ et $\lim_{x \to -\infty} a^x = 0$ ;
\item Si $0 < a < 1$ : $\lim_{x \to +\infty} a^x = 0$ et $\lim_{x \to -\infty} a^x = +\infty$.
\end{itemize}
\textbf{Tableaux de variations (à savoir tracer) :}
\begin{center}
\begin{tabular}{|c|ccc|}\hline
$x$ & $-\infty$ & & $+\infty$ \\ \hline
$f'(x)$ & & $+$ & \\ \hline
$f(x)=a^x\ (a>1)$ & $0$ & $\nearrow$ & $+\infty$ \\ \hline
\end{tabular}
\qquad
\begin{tabular}{|c|ccc|}\hline
$x$ & $-\infty$ & & $+\infty$ \\ \hline
$f'(x)$ & & $-$ & \\ \hline
$f(x)=a^x\ (0<a<1)$ & $+\infty$ & $\searrow$ & $0$ \\ \hline
\end{tabular}
\end{center}
Toutes les courbes passent par le point $(0\,;\,1)$.
\end{propriete}

\begin{popfigure}
\begin{center}
\begin{tikzpicture}
\begin{axis}[width=0.85\linewidth, height=6cm, axis lines=middle, xlabel={$x$}, ylabel={$y$}, xmin=-3.2, xmax=3.2, ymin=-0.5, ymax=8, xtick={-3,-2,-1,1,2,3}, ytick={1,2,4,6}, legend style={at={(0.02,0.98)}, anchor=north west, font=\small}, clip=false]
\addplot[popcurve, domain=-3:3, samples=200] {exp(x*ln(2))};
\addlegendentry{$2^x$}
\addplot[popcurve, poporange, domain=-3:0.9, samples=200] {exp(x*ln(10))};
\addlegendentry{$10^x$}
\addplot[popcurve, popgreen, domain=-3:3, samples=200] {exp(x*ln(0.5))};
\addlegendentry{$(1/2)^x$}
\addplot[mark=*, popink] coordinates {(0,1)};
\end{axis}
\end{tikzpicture}
\end{center}
\legende{Exponentielles de base $a$ : croissance si $a>1$, décroissance si $0<a<1$. Toutes passent par $(0\,;\,1)$.}
\end{popfigure}

\begin{attention}[Piège fréquent : confusion $a^x$ vs $x^a$]
Ne confonds pas $a^x$ et $x^a$ !
\begin{itemize}
\item Dans $a^x$, la \emph{base} est fixe et l'\emph{exposant} varie : la dérivée est $(\ln a)\,a^x$.
\item Dans $x^a$, c'est l'inverse (voir ci-dessous) : la dérivée est $a\,x^{a-1}$.
\end{itemize}
Ainsi la dérivée de $2^x$ n'est \textbf{pas} $x\,2^{x-1}$ (qui est la dérivée de $x^2$), mais bien $(\ln 2)\,2^x$.
\end{attention}

\courssub{Fonction puissance $x \mapsto x^{\alpha}$}

\begin{definition}[Fonction puissance d'exposant réel]
Soit $\alpha$ un réel. Pour tout $x > 0$, on pose :
\[ x^{\alpha} = e^{\alpha \ln x}. \]
La fonction $x \mapsto x^{\alpha}$, définie sur $]0\,;\,+\infty[$, est appelée \cle{fonction puissance} d'exposant $\alpha$.
\end{definition}

On retrouve les cas connus : $x^{1/2} = \sqrt{x}$, $x^{-1} = \dfrac{1}{x}$, $x^2$ au sens usuel.

\begin{propriete}[Dérivée, variations et limites]
La fonction $g : x \mapsto x^{\alpha}$ est dérivable sur $]0\,;\,+\infty[$ et :
\[ g'(x) = \alpha\, x^{\alpha - 1}. \]
En effet, $g(x) = e^{\alpha \ln x}$ donne $g'(x) = \dfrac{\alpha}{x}\, e^{\alpha \ln x} = \alpha\, x^{\alpha - 1}$. Comme $x^{\alpha-1} > 0$ sur $]0\,;\,+\infty[$ :
\begin{itemize}
\item si $\alpha > 0$ : $g$ est strictement croissante sur $]0\,;\,+\infty[$ ;
\item si $\alpha < 0$ : $g$ est strictement décroissante sur $]0\,;\,+\infty[$.
\end{itemize}
Limites aux bornes (selon le signe de $\alpha$) :
\begin{itemize}
\item $\alpha > 0$ : $\lim_{x \to 0^+} x^{\alpha} = 0$, $\lim_{x \to +\infty} x^{\alpha} = +\infty$ ;
\item $\alpha < 0$ : $\lim_{x \to 0^+} x^{\alpha} = +\infty$, $\lim_{x \to +\infty} x^{\alpha} = 0$.
\end{itemize}
\textbf{Tableaux de variations (à savoir tracer) :}
\begin{center}
\begin{tabular}{|c|ccc|}\hline
$x$ & $0^+$ & & $+\infty$ \\ \hline
$g'(x)$ & & $+$ & \\ \hline
$g(x)=x^{\alpha}\ (\alpha>0)$ & $0$ & $\nearrow$ & $+\infty$ \\ \hline
\end{tabular}
\qquad
\begin{tabular}{|c|ccc|}\hline
$x$ & $0^+$ & & $+\infty$ \\ \hline
$g'(x)$ & & $-$ & \\ \hline
$g(x)=x^{\alpha}\ (\alpha<0)$ & $+\infty$ & $\searrow$ & $0$ \\ \hline
\end{tabular}
\end{center}
\end{propriete}

\begin{exemplebox}[Étude rapide de $x \mapsto x^{3/2}$]
Soit $g(x) = x^{3/2} = e^{\frac{3}{2}\ln x}$ sur $]0\,;\,+\infty[$. On a $g'(x) = \dfrac{3}{2} x^{1/2} > 0$ : $g$ est strictement croissante. De plus $\lim_{x \to 0^+} g(x) = 0$ et $\lim_{x \to +\infty} g(x) = +\infty$. Par exemple $g(4) = e^{\frac{3}{2}\ln 4} = 4^{3/2} = \left(\sqrt{4}\right)^3 = 8$.
\end{exemplebox}

\courssub{Croissances comparées}

\begin{aretenir}[Croissances comparées -- Résultats fondamentaux]
Pour tous réels $\alpha > 0$ et $a > 1$ :
\[ \lim_{x \to +\infty} \frac{\ln x}{x^{\alpha}} = 0 \ ; \qquad \lim_{x \to +\infty} \frac{x^{\alpha}}{a^x} = 0. \]
Autrement dit, à l'infini : \textbf{toute puissance l'emporte sur le logarithme, et toute exponentielle de base $a>1$ l'emporte sur toute puissance.}
\[ \ln x \ \ll\ x^{\alpha} \ \ll\ a^x \qquad (x \to +\infty). \]
\end{aretenir}

\begin{methode}[Lever une indétermination par croissance comparée]
Pour une limite en $+\infty$ du type $\dfrac{x^3}{e^x}$ ou $x^2 \ln x$ en $0^+$ :
\begin{enumerate}
\item Identifier les deux « familles » en présence (logarithme, puissance, exponentielle).
\item Factoriser ou réécrire pour faire apparaître un quotient du type $\dfrac{x^{\alpha}}{e^x}$ ou $\dfrac{\ln x}{x^{\alpha}}$.
\item Conclure par croissance comparée : l'exponentielle domine la puissance, la puissance domine le logarithme.
\end{enumerate}
Exemple : $\lim_{x \to +\infty} \dfrac{e^x}{x^2} = +\infty$ car $\dfrac{x^2}{e^x} \to 0$.
\end{methode}

\begin{exoresolu}[Limite avec croissances comparées]
Déterminer $\lim_{x \to +\infty} \left( e^x - x^5 \right)$.
\tcblower
\textbf{Solution.} On a une forme indéterminée « $+\infty - \infty$ ». Factorisons par le terme dominant, $e^x$ :
\[ e^x - x^5 = e^x \left( 1 - \frac{x^5}{e^x} \right). \]
Par croissance comparée avec $\alpha = 5 > 0$ : $\lim_{x \to +\infty} \dfrac{x^5}{e^x} = 0$, donc la parenthèse tend vers $1$. Comme $e^x \to +\infty$, on conclut :
\[ \lim_{x \to +\infty} \left( e^x - x^5 \right) = +\infty. \]
\end{exoresolu}

\begin{savaistu}[Pourquoi « base 2 » en informatique ?]
L'information se mesure en bits : $n$ bits permettent de coder $2^n$ états différents. C'est la fonction exponentielle de base $2$ ! Avec 8 bits (un octet), on code $2^8 = 256$ valeurs. La croissance comparée explique aussi pourquoi un algorithme « exponentiel » devient vite inutilisable : $2^x$ finit toujours par écraser $x^{10}$, même si $x^{10}$ semble d'abord plus grand.
\end{savaistu}

\newpage

\coursec{Activité d'intégration}

\courssub{Objectifs de l'activité}

À l'issue de cette activité d'intégration, tu seras capable de :
\begin{itemize}
\item \cle{modéliser} une situation réelle de croissance ou de décroissance par une fonction de la forme $t \mapsto A\,e^{kt}$ et justifier ce choix par la relation $f' = kf$ ;
\item utiliser les \cle{propriétés algébriques} de l'exponentielle pour transformer et simplifier des expressions ;
\item \cle{résoudre} des équations du type $e^{kt} = a$ (avec $a > 0$) en utilisant la bijectivité de l'exponentielle : $e^x = y \iff x = \ln y$ ;
\item \cle{étudier} le sens de variation d'une fonction $t \mapsto N_0 e^{kt}$ à l'aide de sa dérivée et \cle{tracer} sa courbe représentative ;
\item interpréter concrètement les résultats (prévision démographique, temps de doublement, datation d'un objet archéologique) et comparer les rôles du signe de $k$.
\end{itemize}

\courssub{Situation}

\textbf{Partie A : la croissance démographique de Douala.}

Douala, capitale économique du Cameroun, est l'une des villes à la croissance la plus rapide d'Afrique centrale. Les démographes de l'Institut national de la statistique estiment sa population à $3,9$ millions d'habitants en 2023, avec un taux de croissance annuel moyen de $3,5\,\%$. Ce taux étant supposé constant, on modélise la population (en millions) par la fonction
\[ P(t) = 3,9\,e^{0,035t} \]
où $t$ désigne le nombre d'années écoulées depuis 2023 (ainsi $t = 0$ correspond à l'année 2023).

Les urbanistes s'interrogent : quelle sera la population de Douala en 2030 ? En quelle année la population dépassera-t-elle $6$ millions d'habitants, seuil qui conditionne le dimensionnement du futur réseau d'eau potable de la CAMWATER ? Et au bout de combien d'années la population aura-t-elle doublé ?

\textbf{Partie B : la datation au carbone 14 d'un artefact de Bimbia.}

Sur le site historique de Bimbia (région du Sud-Ouest), lieu de mémoire de la traite négrière, une équipe d'archéologues découvre un fragment de pirogue en bois. Pour dater l'objet, le laboratoire mesure la quantité de carbone 14, isotope radioactif qui se désintègre après la mort de l'arbre. On admet que la quantité de carbone 14 présente dans l'échantillon suit la loi
\[ N(t) = N_0\,e^{-0,000121\,t} \]
où $N_0$ est la quantité initiale (à la mort de l'arbre) et $t$ le temps écoulé en années. L'analyse révèle qu'il ne reste que $60\,\%$ de la quantité initiale de carbone 14. Quel est l'âge de la pirogue ? Comment évolue la quantité de carbone 14 au fil des siècles ?

\courssub{Consignes}

Le travail se fait en groupes de trois ou quatre élèves. Chaque groupe rédige un rapport répondant aux questions suivantes.

\textbf{Partie A.}
\begin{enumerate}
\item Vérifier que $P(0) = 3,9$ et interpréter ce résultat.
\item Calculer $P'(t)$ et montrer que $P'(t) = 0,035\,P(t)$. Expliquer pourquoi cette relation traduit un taux de croissance constant de $3,5\,\%$ par an.
\item En déduire le sens de variation de $P$ sur $[0\,;+\infty[$.
\item Estimer la population de Douala en 2030 (arrondir à $0,1$ million près).
\item Résoudre l'équation $P(t) = 6$ et en déduire l'année au cours de laquelle la population dépassera $6$ millions.
\item Déterminer le temps de doublement $T$, c'est-à-dire la solution de $P(t) = 2 \times 3,9$. Que remarque-t-on ?
\end{enumerate}

\textbf{Partie B.}
\begin{enumerate}
\item[7.] Traduire l'information « il reste $60\,\%$ du carbone initial » par une équation d'inconnue $t$, puis la résoudre. Donner l'âge de la pirogue arrondi à la dizaine d'années près.
\item[8.] Calculer $N'(t)$ et dresser le tableau de variations de $N$ sur $[0\,;+\infty[$. Quelle est la limite de $N(t)$ quand $t \to +\infty$ ? Interpréter.
\item[9.] Tracer l'allure de la courbe de $t \mapsto e^{-0,000121\,t}$ (on pourra calculer quelques valeurs : $t = 0$, $t = 4000$, $t = 8000$).
\item[10.] Calculer la demi-vie du carbone 14, c'est-à-dire le temps $t_{1/2}$ tel que $N(t_{1/2}) = \dfrac{N_0}{2}$.
\end{enumerate}

\textbf{Synthèse.}
\begin{enumerate}
\item[11.] Les deux modèles sont de la forme $f(t) = A\,e^{kt}$. Comparer les signes de $k$ dans les deux parties et conclure sur le rôle du signe de $k$ dans le sens de variation et le comportement en $+\infty$.
\end{enumerate}

\courssub{Ressources à mobiliser}

\begin{aretenir}[Boîte à outils de la leçon]
\begin{itemize}
\item \cle{Bijectivité} : pour tout $y > 0$, $e^x = y \iff x = \ln y$. En particulier, pour $a > 0$ et $k \neq 0$ : $e^{kt} = a \iff t = \dfrac{\ln a}{k}$.
\item \cle{Propriétés algébriques} : $e^{a+b} = e^a e^b$, $e^{-a} = \dfrac{1}{e^a}$, $(e^a)^n = e^{na}$ ($n \in \mathbb{Z}$), $\ln(e^x) = x$ et $e^{\ln x} = x$ pour $x > 0$.
\item \cle{Dérivation} : $(e^u)' = u' e^u$ ; donc si $f(t) = A e^{kt}$, alors $f'(t) = kA e^{kt} = k f(t)$.
\item \cle{Signe et variations} : $e^x > 0$ pour tout réel $x$ ; si $A > 0$, le signe de $f'$ est celui de $k$.
\item \cle{Limites} : $\lim_{x \to +\infty} e^x = +\infty$ et $\lim_{x \to -\infty} e^x = 0$.
\item \cle{Primitive} : une primitive de $u' e^u$ est $e^u$ ; pour $k \neq 0$, une primitive de $e^{kt}$ est $\dfrac{1}{k}e^{kt}$.
\end{itemize}
\end{aretenir}

\begin{attention}[Conditions d'application à ne pas oublier]
\begin{itemize}
\item L'équivalence $e^x = y \iff x = \ln y$ n'est valable que pour $y > 0$ : l'équation $e^x = -3$ n'a \textbf{aucune} solution, car une exponentielle est toujours strictement positive.
\item Dans la formule $t = \dfrac{\ln a}{k}$, le coefficient $k$ doit être non nul, et il faut bien diviser par $k$ \emph{avec son signe} : en partie B, $k = -0,000121 < 0$, et $\ln 0,6 < 0$, donc le quotient est bien positif.
\item Ne pas confondre $\ln(e^x) = x$ (valable pour tout $x$ réel) et $e^{\ln x} = x$ (valable seulement pour $x > 0$).
\end{itemize}
\end{attention}

\courssub{Démarche et corrigé commenté}

\begin{exoresolu}[Corrigé complet de l'activité]
\textbf{Partie A.} 1. Vérifier $P(0)$ et interpréter. 2. Montrer que $P' = 0,035\,P$. 3. Variations de $P$. 4. Population en 2030. 5. Année où $P > 6$ millions. 6. Temps de doublement. \textbf{Partie B.} 7. Âge de la pirogue. 8. Variations et limite de $N$. 9. Courbe. 10. Demi-vie. 11. Synthèse.
\tcblower
\textbf{1.} $P(0) = 3,9\,e^{0} = 3,9 \times 1 = 3,9$. En 2023 ($t = 0$), le modèle redonne bien la population initiale de $3,9$ millions : le modèle est cohérent avec la donnée de départ.

\textbf{2.} $P$ est de la forme $A\,e^{u}$ avec $u(t) = 0,035t$, donc $u'(t) = 0,035$. D'après la formule $(e^u)' = u'e^u$ :
\[ P'(t) = 3,9 \times 0,035\,e^{0,035t} = 0,035 \times \underbrace{3,9\,e^{0,035t}}_{P(t)} = 0,035\,P(t). \]
La vitesse de croissance $P'(t)$ est à chaque instant proportionnelle à la population $P(t)$, avec un coefficient de proportionnalité constant $0,035$ : c'est exactement la signification d'un taux de croissance annuel constant de $3,5\,\%$. C'est ce qui justifie le choix du modèle exponentiel.

\textbf{3.} Pour tout $t \geq 0$, $e^{0,035t} > 0$ donc $P(t) > 0$ ; comme $0,035 > 0$, on a $P'(t) = 0,035\,P(t) > 0$ : $P$ est \cle{strictement croissante} sur $[0\,;+\infty[$. De plus, quand $t \to +\infty$, $0,035t \to +\infty$, donc $\lim_{t \to +\infty} e^{0,035t} = +\infty$ et $\lim_{t \to +\infty} P(t) = +\infty$ : la population croît sans borne dans ce modèle, ce qui en montre les limites à très long terme (les ressources d'une ville ne sont pas infinies).

\textbf{4.} L'année 2030 correspond à $t = 2030 - 2023 = 7$ :
\[ P(7) = 3,9\,e^{0,035 \times 7} = 3,9\,e^{0,245} \approx 3,9 \times 1,2779 \approx 4,98. \]
Douala compterait environ $5,0$ millions d'habitants en 2030.

\textbf{5.} On résout $P(t) = 6$. Comme $\dfrac{6}{3,9} > 0$, l'équivalence $e^x = y \iff x = \ln y$ s'applique :
\begin{align*}
3,9\,e^{0,035t} = 6 &\iff e^{0,035t} = \frac{6}{3,9} \\
&\iff 0,035t = \ln\left(\frac{6}{3,9}\right) \\
&\iff t = \frac{\ln(6/3,9)}{0,035} \approx \frac{0,4308}{0,035} \approx 12,3.
\end{align*}
Comme $P$ est strictement croissante, $P(t) > 6$ dès que $t > 12,3$. Or $2023 + 12,3 \approx 2035,3$ : la population dépasserait $6$ millions \textbf{au cours de l'année 2035}.

\textbf{6.} $P(t) = 7,8 \iff 3,9\,e^{0,035t} = 2 \times 3,9 \iff e^{0,035t} = 2 \iff t = \dfrac{\ln 2}{0,035} \approx \dfrac{0,6931}{0,035} \approx 19,8$ ans.
Le temps de doublement est d'environ $20$ ans et \textbf{ne dépend pas de la population initiale} $3,9$ (elle s'est simplifiée dans le calcul) : c'est une propriété caractéristique de la croissance exponentielle.

\textbf{7.} « Il reste $60\,\%$ du carbone initial » se traduit par $N(t) = 0,6\,N_0$. Comme $N_0 > 0$, on peut simplifier par $N_0$ :
\begin{align*}
N_0\,e^{-0,000121t} = 0,6\,N_0 &\iff e^{-0,000121t} = 0,6 \\
&\iff -0,000121\,t = \ln 0,6 \\
&\iff t = \frac{\ln 0,6}{-0,000121} = \frac{-0,5108}{-0,000121} \approx 4222.
\end{align*}
La pirogue a environ \textbf{4220 ans} : elle daterait d'environ 2200 avant notre ère, bien avant la période de la traite négrière — un résultat que les archéologues devront interpréter avec soin (l'arbre a pu être abattu longtemps avant l'utilisation du bois, ou l'échantillon peut être antérieur au site).

\textbf{8.} $N$ est de la forme $N_0 e^u$ avec $u(t) = -0,000121t$, donc :
\[ N'(t) = -0,000121\,N_0\,e^{-0,000121t}. \]
Comme $N_0 > 0$, $e^{-0,000121t} > 0$ et $-0,000121 < 0$, on a $N'(t) < 0$ pour tout $t \geq 0$ : $N$ est \cle{strictement décroissante} sur $[0\,;+\infty[$. Quand $t \to +\infty$, $-0,000121t \to -\infty$, donc $\lim_{t \to +\infty} N(t) = 0$ : le carbone 14 finit par disparaître presque totalement, ce qui limite la méthode de datation aux objets de moins d'environ $50\,000$ ans.

\begin{center}
\begin{tabular}{|c|ccc|}
\hline
\rowcolor{popblueL} $t$ & $0$ & \hspace{2.5cm} & $+\infty$ \\
\hline
Signe de $N'(t)$ & \multicolumn{3}{c|}{$-$} \\
\hline
$N(t)$ & $N_0$ & $\searrow$ & $0$ \\
\hline
\end{tabular}
\end{center}

\textbf{9.} Valeurs repères : $e^{0} = 1$ pour $t = 0$ ; $e^{-0,000121 \times 4000} = e^{-0,484} \approx 0,62$ pour $t = 4000$ ; $e^{-0,000121 \times 8000} = e^{-0,968} \approx 0,38$ pour $t = 8000$. La courbe part de $1$ et décroît en s'approchant de l'axe des abscisses, qui est asymptote horizontale.

\begin{popfigure}
\begin{center}
\begin{tikzpicture}[scale=1]
\draw[->,thick] (0,0) -- (9,0) node[below left]{\small $t$ (milliers d'années)};
\draw[->,thick] (0,0) -- (0,3) node[below left]{\small $N/N_0$};
\draw[domain=0:8.5,samples=100,very thick,popblue] plot (\x,{2.4*exp(-0.484*\x/4)});
\draw[dashed,popmuted] (0,2.4) node[left]{\small $1$} -- (0.1,2.4);
\foreach \x/\y in {4/{2.4*exp(-0.484)},8/{2.4*exp(-0.968)}} {\fill[poporange] (\x,{\y}) circle (1.6pt);}
\draw[dashed,popmuted] (4,0) node[below]{\small $4$} -- (4,{2.4*exp(-0.484)});
\draw[dashed,popmuted] (8,0) node[below]{\small $8$} -- (8,{2.4*exp(-0.968)});
\end{tikzpicture}
\end{center}
\legende{Décroissance radioactive du carbone 14 : $t \mapsto e^{-0,000121t}$.}
\end{popfigure}

\textbf{10.} $N(t) = \dfrac{N_0}{2} \iff e^{-0,000121t} = 0,5 \iff -0,000121\,t = \ln 0,5 = -\ln 2 \iff t = \dfrac{\ln 2}{0,000121} \approx \dfrac{0,6931}{0,000121} \approx 5728$ ans. On retrouve la célèbre demi-vie du carbone 14, environ $5730$ ans : tous les $5730$ ans, la quantité de carbone 14 est divisée par deux.

\textbf{11.} Dans les deux cas, $f(t) = A e^{kt}$ vérifie $f' = kf$. Si $k > 0$ (Douala, $k = 0,035$), la fonction est strictement croissante et tend vers $+\infty$ ; si $k < 0$ (carbone 14, $k = -0,000121$), elle est strictement décroissante et tend vers $0$. Le signe de $k$ gouverne donc entièrement le comportement du phénomène : croissance ou décroissance exponentielle.
\end{exoresolu}

\courssub{Bilan de l'activité}

Cette activité a permis de mettre en évidence les points essentiels de la leçon :
\begin{itemize}
\item une même famille de fonctions, $t \mapsto A e^{kt}$, modélise des phénomènes très différents (démographie urbaine, radioactivité) : c'est l'\cle{universalité} du modèle exponentiel ;
\item la relation $f' = kf$ est la \cle{signature} des phénomènes à taux d'évolution constant, et elle justifie le choix du modèle ;
\item la résolution d'équations en $e^x$ repose toujours sur la même clé : $e^x = y \iff x = \ln y$ (pour $y > 0$), qui a donné aussi bien l'année du seuil des $6$ millions que l'âge de la pirogue ;
\item le \cle{signe de $k$} détermine tout : croissance illimitée si $k > 0$, extinction progressive si $k < 0$ ;
\item les grandeurs caractéristiques (temps de doublement, demi-vie) sont indépendantes de la quantité initiale : elles ne dépendent que de $k$ via la formule $\dfrac{\ln 2}{|k|}$.
\end{itemize}
Tu es désormais prêt à modéliser par toi-même d'autres situations : propagation d'une épidémie, refroidissement d'un liquide, amortissement d'un prêt bancaire, charge d'un condensateur.

\newpage

\coursec{Méthodes et exercices résolus}

\coursec{Fonction exponentielle népérienne}

\courssub{Définition et premières propriétés}

\begin{definition}[Fonction exponentielle népérienne]
La fonction exponentielle népérienne, notée $\exp$ ou $x\mapsto e^{x}$, est la fonction réciproque de la fonction logarithme népérien $\ln : ]0\ ;\ +\infty[ \to \mathbb{R}$.
Elle est définie sur $\mathbb{R}$ à valeurs dans $]0\ ;\ +\infty[$ et vérifie :
\[ \forall x\in\mathbb{R},\ \forall y>0,\qquad e^{x}=y \;\Longleftrightarrow\; \ln y = x. \]
Le nombre $e$ est l'unique réel tel que $\ln e = 1$ ($e\approx 2,71828$).
\end{definition}

\begin{propriete}[Propriétés fondamentales issues de la réciproque]
Pour tout réel $x$ et tout réel strictement positif $y$ :
\begin{itemize}
\item $e^{\ln y} = y$ \quad (exponentielle après logarithme)
\item $\ln(e^{x}) = x$ \quad (logarithme après exponentielle)
\item $e^{0}=1$ \quad et \quad $e^{1}=e$.
\end{itemize}
La fonction $\exp$ est \textbf{strictement croissante}, \textbf{continue} et \textbf{dérivable} sur $\mathbb{R}$.
\end{propriete}

\begin{experience}[Découvrir la dérivée de l'exponentielle par la réciproque]
On sait que $\ln$ est dérivable sur $]0\ ;\ +\infty[$ avec $(\ln)'(x)=\dfrac{1}{x}$.
Soit $y=e^{x}$ (donc $x=\ln y$). Dérivons l'égalité $x=\ln y$ par rapport à $x$ :
\[ 1 = \frac{1}{y}\, \frac{\mathrm{d}y}{\mathrm{d}x} \quad\Longrightarrow\quad \frac{\mathrm{d}y}{\mathrm{d}x} = y = e^{x}. \]
On retrouve ainsi $(\exp)' = \exp$. Cette méthode (dérivée de la réciproque) est une démonstration classique exigible au Bac.
\end{experience}

\begin{savaistu}[Le nombre $e$ et le Cameroun]
Le nombre $e$ apparaît naturellement dans tout phénomène où la vitesse d'évolution est proportionnelle à la quantité présente : croissance d'une population (ex: démographie de Yaoundé ou Douala), décroissance radioactive (datation au carbone 14 des sites archéologiques de la vallée de la Bénoué), intérêts composés continus (épargne à la BCEAO), ou refroidissement d'un corps (loi de Newton). C'est la base « naturelle » des mathématiques du changement.
\end{savaistu}

\courssub{Propriétés algébriques de l'exponentielle}

\begin{propriete}[Formules algébriques de $\exp$]
Pour tous réels $a,b$ et tout entier relatif $n$ :
\begin{align*}
e^{a+b} &= e^{a}\,e^{b} & e^{a-b} &= \dfrac{e^{a}}{e^{b}} & e^{-a} &= \dfrac{1}{e^{a}} \\
\left(e^{a}\right)^{n} &= e^{na} & e^{\ln a} &= a\ (a>0) & \ln(e^{a}) &= a.
\end{align*}
\textbf{Réflexe clé :} une \textbf{somme} dans l'exposant devient un \textbf{produit} d'exponentielles ; un \textbf{produit} dans l'exposant devient une \textbf{puissance}.
\end{propriete}

\begin{methode}[Simplifier une expression contenant des exponentielles]
\begin{enumerate}
\item \cle{Repérer} les exposants : ce sont eux qui subissent les opérations.
\item Appliquer les \cle{formules fondamentales} (propriété ci-dessus).
\item Utiliser le \cle{lien avec le logarithme} : $e^{\ln x}=x\ (x>0)$ et $\ln(e^{x})=x$.
\item \cle{Factoriser} si possible par le facteur exponentiel commun (souvent $e^{x}$ ou $e^{2x}$).
\item \cle{Vérifier} le résultat sur une valeur numérique simple ($x=0$ ou $x=1$).
\end{enumerate}
\end{methode}

\begin{exoresolu}[Simplification d'expressions exponentielles]
\textbf{Énoncé.} Simplifier les expressions suivantes sous la forme la plus simple :
\begin{enumerate}
\item $A=\dfrac{e^{3x+1}\times e^{2-x}}{e^{2x}}$ ;
\item $B=\left(e^{x}\right)^{2}\times e^{-3x}+\dfrac{e^{x+1}}{e^{x}}$ ;
\item $C=e^{\ln 5}+e^{2\ln 3}-\ln\left(e^{4}\right)$.
\end{enumerate}
\tcblower
\textbf{Solution détaillée.}
\begin{enumerate}
\item \textbf{Calcul de $A$.} Numérateur : $e^{(3x+1)+(2-x)}=e^{2x+3}$.
\[ A = \dfrac{e^{2x+3}}{e^{2x}} = e^{(2x+3)-2x} = e^{3}. \]
$A$ est une constante ($e^{3}$), indépendante de $x$.

\item \textbf{Calcul de $B$.} $\left(e^{x}\right)^{2}=e^{2x}$, donc $e^{2x}\times e^{-3x}=e^{-x}$.
\[ \dfrac{e^{x+1}}{e^{x}}=e^{(x+1)-x}=e^{1}=e. \]
\textbf{Conclusion :} $B=e^{-x}+e$.

\item \textbf{Calcul de $C$.} $e^{\ln 5}=5$ ; $e^{2\ln 3}=e^{\ln(3^{2})}=9$ ; $\ln(e^{4})=4$.
\[ C = 5 + 9 - 4 = 10. \]

\item \textbf{Vérification (réflexe) :} Pour $x=0$, $A_{\text{init}}=\dfrac{e^{1}e^{2}}{1}=e^{3}$ (OK). $B_{\text{init}}=1\times 1+\dfrac{e}{1}=1+e$ et $e^{-0}+e=1+e$ (OK).
\end{enumerate}
\end{exoresolu}

\courssub{Équations, inéquations et systèmes d'inconnue $e^{x}$}

\begin{methode}[Résoudre avec l'exponentielle]
\textbf{Cas 1 : Équation ou inéquation $e^{u(x)}=e^{v(x)}$ ou $e^{u(x)}=a$.}
\begin{enumerate}
\item Se ramener à une égalité entre \cle{deux exponentielles}, ou écrire $a=e^{\ln a}$ si $a>0$.
\item Utiliser l'\cle{injectivité} : $e^{A}=e^{B}\Leftrightarrow A=B$.
\item Pour les inéquations, utiliser la \cle{stricte croissance} : $e^{A}<e^{B}\Leftrightarrow A<B$ (sens conservé).
\item Si $a\leq 0$ : $e^{x}=a$ pas de solution ; $e^{x}>a$ vrai $\forall x$ ; $e^{x}<a$ pas de solution.
\end{enumerate}

\textbf{Cas 2 : Équation bicarrée $\alpha e^{2x}+\beta e^{x}+\gamma=0$.}
\begin{enumerate}
\item Poser $X=e^{x}$ avec la condition \cle{$X>0$}.
\item Résoudre $\alpha X^{2}+\beta X+\gamma=0$.
\item \cle{Écarter} toute racine $\leq 0$.
\item Conclure par $x=\ln X$ pour chaque racine positive.
\end{enumerate}

\textbf{Cas 3 : Système.} Poser $X=e^{x}>0$, $Y=e^{y}>0$, résoudre en $(X,Y)$, vérifier positivité, conclure.
\end{methode}

\begin{exoresolu}[Équation bicarrée, inéquation et système]
\textbf{Énoncé (type Bac).}
\begin{enumerate}
\item Résoudre dans $\mathbb{R}$ : $(E)\ e^{2x}-3e^{x}+2=0$.
\item Résoudre dans $\mathbb{R}$ : $(I)\ e^{2x}-3e^{x}+2>0$.
\item Résoudre $\begin{cases} e^{x}e^{y}=e^{5} \\ e^{x}-e^{y}=e^{2}-e^{3}\end{cases}$
\end{enumerate}
\tcblower
\textbf{Solution détaillée.}
\begin{enumerate}
\item Posons $X=e^{x}>0$. $X^{2}-3X+2=0 \Rightarrow \Delta=1$, $X_{1}=1$, $X_{2}=2$ (toutes deux $>0$).
$e^{x}=1 \Leftrightarrow x=0$ ; $e^{x}=2 \Leftrightarrow x=\ln 2$.
$\mathcal{S}=\{0\ ;\ \ln 2\}$.

\item Même changement de variable : $(X-1)(X-2)>0$ avec $X>0$.
$X<1$ ou $X>2$.
$e^{x}<1 \Leftrightarrow x<0$ ; $e^{x}>2 \Leftrightarrow x>\ln 2$.
$\mathcal{S}=]-\infty\ ;\ 0[\ \cup\ ]\ln 2\ ;\ +\infty[$.

\item Posons $X=e^{x}>0$, $Y=e^{y}>0$.
$\begin{cases} XY=e^{5} \\ X-Y=e^{2}-e^{3} \end{cases} \Rightarrow X=Y+e^{2}-e^{3}$.
$(Y+e^{2}-e^{3})Y=e^{5} \Leftrightarrow Y^{2}+(e^{2}-e^{3})Y-e^{5}=0$.
$\Delta=(e^{2}-e^{3})^{2}+4e^{5}=e^{4}+2e^{5}+e^{6}=(e^{2}+e^{3})^{2}$.
$\sqrt{\Delta}=e^{2}+e^{3}$.
$Y=\dfrac{-(e^{2}-e^{3})\pm(e^{2}+e^{3})}{2}$.
$Y_{+}=e^{3}>0$ (gardée) ; $Y_{-}=-e^{2}<0$ (rejetée).
$X=e^{3}+e^{2}-e^{3}=e^{2}>0$.
Solution unique : $(x\ ;\ y)=(2\ ;\ 3)$.
\end{enumerate}
\end{exoresolu}

\courssub{Limites de la fonction exponentielle}

\begin{propriete}[Limites de référence de $\exp$]
\begin{align*}
\lim_{x\to +\infty}e^{x} &= +\infty & \lim_{x\to -\infty}e^{x} &= 0^{+} \\
\lim_{x\to +\infty}\dfrac{e^{x}}{x} &= +\infty & \lim_{x\to -\infty}x\,e^{x} &= 0 \\
\lim_{x\to 0}\dfrac{e^{x}-1}{x} &= 1 \quad\text{(taux d'accroissement en 0)}.
\end{align*}
L'exponentielle l'emporte sur toute puissance de $x$ en $+\infty$ (croissance comparée).
\end{propriete}

\begin{methode}[Calculer une limite contenant $e^{u(x)}$]
\begin{enumerate}
\item Calculer $\ell=\lim u(x)$.
\item Appliquer le théorème de composition : $\lim e^{u(x)} = \lim_{X\to\ell}e^{X}$.
\item Pour $e^{ax+b}$ : le signe de $a$ décide ($a>0 \Rightarrow +\infty$ en $+\infty$ ; $a<0 \Rightarrow 0$ en $+\infty$).
\item En cas de forme indéterminée (somme/produit), \cle{factoriser par le terme dominant} (l'exponentielle la plus forte).
\item Pour $\dfrac{e^{u}-1}{u}$ avec $u\to 0$, utiliser la limite $1$.
\end{enumerate}
\end{methode}

\begin{exoresolu}[Limites de fonctions exponentielles]
\textbf{Énoncé (type Bac).} Calculer les limites suivantes :
\begin{enumerate}
\item $\lim\limits_{x\to +\infty}\left(e^{x}-3x^{2}+1\right)$
\item $\lim\limits_{x\to +\infty}\dfrac{e^{2x+1}}{x}$
\item $\lim\limits_{x\to +\infty}e^{-3x+2}$
\item $\lim\limits_{x\to 0}\dfrac{e^{5x}-1}{x}$
\end{enumerate}
\tcblower
\textbf{Solution détaillée.}
\begin{enumerate}
\item Forme $+\infty-\infty$. Factorisation par $e^{x}$ (dominant) :
$e^{x}-3x^{2}+1 = e^{x}\left(1-\dfrac{3x^{2}}{e^{x}}+\dfrac{1}{e^{x}}\right)$.
$\lim\dfrac{x^{2}}{e^{x}}=0$, $\lim\dfrac{1}{e^{x}}=0$ $\Rightarrow$ parenthèse $\to 1$.
$e^{x}\to +\infty$ $\Rightarrow$ limite $= +\infty$.

\item Posons $X=2x+1 \to +\infty$, $x=\frac{X-1}{2}$.
$\dfrac{e^{2x+1}}{x} = \dfrac{2e^{X}}{X-1} = 2\dfrac{X}{X-1}\dfrac{e^{X}}{X}$.
$\dfrac{X}{X-1}\to 1$, $\dfrac{e^{X}}{X}\to +\infty$ $\Rightarrow$ limite $= +\infty$.

\item $\lim_{x\to+\infty}(-3x+2)=-\infty$. Composition : $\lim_{X\to-\infty}e^{X}=0$.
Limite $= 0$. Asymptote horizontale $y=0$ en $+\infty$.

\item $\dfrac{e^{5x}-1}{x} = 5\dfrac{e^{5x}-1}{5x}$. Posons $u=5x\to 0$.
$\lim_{u\to 0}\dfrac{e^{u}-1}{u}=1$ $\Rightarrow$ limite $= 5\times 1 = 5$.
\end{enumerate}
\end{exoresolu}

\courssub{Dérivation, variations et courbe représentative}

\begin{propriete}[Dérivée de l'exponentielle composée]
Soit $u$ dérivable sur $I$. La fonction $x\mapsto e^{u(x)}$ est dérivable sur $I$ et :
\[ \left(e^{u}\right)' = u'\,e^{u}. \]
Comme $e^{u(x)}>0$ pour tout $x$, le signe de $\left(e^{u}\right)'$ est \textbf{exactement le signe de $u'$}.
\end{propriete}

\begin{methode}[Étude complète d'une fonction $f(x)=e^{u(x)}$ ou combinaison]
\begin{enumerate}
\item \cle{Ensemble de définition} : $\mathbb{R}$ si $u$ polynôme ; attention aux dénominateurs non exponentiels.
\item \cle{Limites} aux bornes ; en déduire asymptotes (horizontale si limite finie).
\item \cle{Dériver} : $(e^{u})'=u'e^{u}$. Ne jamais oublier $u'$.
\item \cle{Signe de $f'$} : signe de $u'$ (ou du facteur polynomial accompagnant $e^{u}$). Tableau de variations.
\item \cle{Points remarquables} : $f(0)$, tangente en $a$ : $y=f'(a)(x-a)+f(a)$.
\item \cle{Tracé} : asymptotes, points, variations.
\end{enumerate}
\end{methode}

\begin{exoresolu}[Étude complète de $f(x)=(x-1)e^{x}+2$]
\textbf{Énoncé (type Bac C/D).} Soit $f(x)=(x-1)e^{x}+2$ sur $\mathbb{R}$, $(\mathcal{C})$ sa courbe.
\begin{enumerate}
\item Limites en $\pm\infty$ et interprétation graphique.
\item Dérivée, tableau de variations.
\item Équation de la tangente $(T)$ en $0$.
\item Tracer $(\mathcal{C})$ et $(T)$.
\end{enumerate}
\tcblower
\textbf{Solution détaillée.}
\begin{enumerate}
\item \textbf{En $-\infty$ :} $f(x)=xe^{x}-e^{x}+2$. $\lim xe^{x}=0$, $\lim e^{x}=0$ $\Rightarrow$ $\lim f(x)=2$.
Asymptote horizontale $y=2$ en $-\infty$.
\textbf{En $+\infty$ :} $(x-1)\to+\infty$, $e^{x}\to+\infty$ $\Rightarrow$ $f(x)\to+\infty$.

\item $f'(x)=1\cdot e^{x}+(x-1)e^{x}=xe^{x}$.
$e^{x}>0$ $\Rightarrow$ signe de $f'$ = signe de $x$.
$f'(x)<0$ sur $]-\infty;0[$, $f'(0)=0$, $f'(x)>0$ sur $]0;+\infty[$.
Minimum $f(0)=1$.

\begin{center}
\begin{tabular}{|c|ccccc|}
\hline
$x$ & $-\infty$ & & $0$ & & $+\infty$ \\
\hline
$f'(x)$ & & $-$ & $0$ & $+$ & \\
\hline
$f(x)$ & $2$ & $\searrow$ & $1$ & $\nearrow$ & $+\infty$ \\
\hline
\end{tabular}
\end{center}

\item Tangente en $0$ : $y=f'(0)(x-0)+f(0)=0\cdot x+1=1$. $(T): y=1$ (horizontale).

\item \textbf{Tracé} : voir figure ci-dessous.
\end{enumerate}
\textbf{Conclusion :} $f(x)\geq 1$ pour tout $x\in\mathbb{R}$.
\end{exoresolu}

\begin{popfigure}
\begin{tikzpicture}[scale=0.7, domain=-3:2]
\draw[->, thick] (-3.5,0) -- (2.5,0) node[right] {$x$};
\draw[->, thick] (0,-1) -- (0,8) node[above] {$y$};
% Asymptote y=2
\draw[dashed, popmuted] (-3.5,2) -- (2.5,2) node[right] {$y=2$};
% Tangent y=1
\draw[dashed, poporange] (-3.5,1) -- (2.5,1) node[right] {$(T): y=1$};
% Axes graduations
\foreach \x in {-3,-2,-1,1,2} \draw (\x,0.1) -- (\x,-0.1) node[below] {$\x$};
\foreach \y in {1,2,3,4,5,6,7} \draw (0.1,\y) -- (-0.1,\y) node[left] {$\y$};
% Function plot
\draw[popcurve, thick, samples=200] plot (\x, {(\x-1)*exp(\x)+2});
% Points
\fill[popblue] (0,1) circle (2pt) node[below left] {$(0;1)$};
\fill[popblue] (1,2) circle (2pt) node[above right] {$(1;2)$};
\fill[popblue] (2,{exp(2)+2}) circle (2pt) node[above right] {$(2;e^2+2)$};
\end{tikzpicture}
\legende{Courbe $(\mathcal{C})$ de $f(x)=(x-1)e^x+2$, asymptote $y=2$ et tangente $(T): y=1$ au minimum.}
\end{popfigure}

\courssub{Primitives et modélisation par l'exponentielle}

\begin{propriete}[Primitives de la forme $u'\,e^{u}$]
Soit $u$ dérivable sur $I$. Une primitive de $x\mapsto u'(x)\,e^{u(x)}$ sur $I$ est $x\mapsto e^{u(x)}$.
\[ \int u'(x)\,e^{u(x)}\,\mathrm{d}x = e^{u(x)} + C. \]
Cas particulier : une primitive de $e^{ax+b}\ (a\neq 0)$ est $\dfrac{1}{a}e^{ax+b}$.
\end{propriete}

\begin{methode}[Rechercher une primitive de $u'\,e^{u}$]
\begin{enumerate}
\item Identifier $u$ dans l'exposant, calculer $u'$.
\item Comparer le facteur devant $e^{u}$ avec $u'$.
\item S'ils coïncident (à une constante $k$ près) : primitive $= k\,e^{u}$.
\item \cle{Vérifier} en dérivant le résultat.
\end{enumerate}
\end{methode}

\begin{exoresolu}[Recherche de primitives]
\textbf{Énoncé.} Une primitive sur $\mathbb{R}$ de :
\begin{enumerate}
\item $f(x)=e^{4x-1}$
\item $g(x)=x\,e^{x^{2}}$
\item $h(x)=(2x+1)e^{x^{2}+x-3}$
\end{enumerate}
\tcblower
\textbf{Solution.}
\begin{enumerate}
\item $f(x)=e^{4x-1}$, $a=4$. $F(x)=\dfrac{1}{4}e^{4x-1}$. Vérif : $F'=\frac{1}{4}\cdot 4e^{4x-1}=f$.
\item $u(x)=x^{2}$, $u'=2x$. $g(x)=\frac{1}{2}\cdot 2x\,e^{x^{2}}=\frac{1}{2}u'e^{u}$. $G(x)=\frac{1}{2}e^{x^{2}}$. Vérif : $G'=\frac{1}{2}\cdot 2x\,e^{x^{2}}=g$.
\item $u(x)=x^{2}+x-3$, $u'=2x+1$. $h=u'e^{u}$. $H(x)=e^{x^{2}+x-3}$. Vérif : $H'=(2x+1)e^{x^{2}+x-3}=h$.
\end{enumerate}
\end{exoresolu}

\begin{methode}[Modélisation exponentielle : croissance, décroissance, refroidissement]
\begin{enumerate}
\item \cle{Modèle} : $N(t)=N_{0}\,e^{kt}$ ($N_{0}=N(0)$, $k$ constante). $k>0$ croissance, $k<0$ décroissance.
\item \cle{Paramètres} : $N_{0}$ lu à $t=0$ ; $k=\dfrac{1}{t_{1}}\ln\dfrac{N(t_{1})}{N_{0}}$.
\item \cle{Exploitation} : valeur à $t$ donné, ou $t=\dfrac{1}{k}\ln\dfrac{N_{\text{cible}}}{N_{0}}$.
\item \cle{Demi-vie} $T$ (décroissance) : $e^{kT}=\frac{1}{2} \Rightarrow T=\dfrac{\ln 2}{|k|}$.
\item \cle{Conclusion} : phrase complète, unité, arrondi.
\end{enumerate}
\end{methode}

\begin{exoresolu}[Refroidissement d'une boisson chaude à Yaoundé]
\textbf{Énoncé (type Bac).} Café servi à $\theta_{0}=90^{\circ}\text{C}$, salle à $\theta_{a}=25^{\circ}\text{C}$.
Écart : $\theta(t)-25=(\theta_{0}-25)e^{-kt}$ ($t$ en min, $k>0$). $\theta(5)=65^{\circ}\text{C}$.
\begin{enumerate}
\item $k$ exact et approché $10^{-3}$.
\item $\theta(15)$ au degré près.
\item Temps pour $\theta=30^{\circ}\text{C}$ à la minute près.
\end{enumerate}
\tcblower
\textbf{Solution.}
\begin{enumerate}
\item Modèle : $\theta(t)=25+65\,e^{-kt}$.
$\theta(5)=65 \Rightarrow 25+65e^{-5k}=65 \Rightarrow e^{-5k}=\frac{40}{65}=\frac{8}{13}$.
$-5k=\ln\frac{8}{13} \Rightarrow k=\frac{1}{5}\ln\frac{13}{8}$.
$k\approx \frac{0,4855}{5}\approx 0,097\ \text{min}^{-1}$.

\item $\theta(15)=25+65\,e^{-15k}=25+65\left(e^{-5k}\right)^{3}=25+65\left(\frac{8}{13}\right)^{3}$.
$\left(\frac{8}{13}\right)^{3}=\frac{512}{2197}\approx 0,2330$.
$\theta(15)\approx 25+65\times 0,2330\approx 40,1$.
\textbf{Conclusion :} $\approx 40^{\circ}\text{C}$.

\item $\theta(t)=30 \Rightarrow 25+65e^{-kt}=30 \Rightarrow e^{-kt}=\frac{5}{65}=\frac{1}{13}$.
$-kt=\ln\frac{1}{13}=-\ln 13 \Rightarrow t=\frac{\ln 13}{k}$.
$t\approx \frac{2,565}{0,097}\approx 26,4$ min.
\textbf{Conclusion :} environ $26$ minutes. La température tend vers $25^{\circ}\text{C}$ (asymptote), cohérent physiquement.
\end{enumerate}
\end{exoresolu}

\begin{popfigure}
\begin{tikzpicture}[scale=0.7, domain=0:30]
\draw[->, thick] (0,0) -- (32,0) node[right] {$t$ (min)};
\draw[->, thick] (0,20) -- (0,95) node[above] {$\theta(t)\ (^{\circ}\text{C})$};
% Asymptote
\draw[dashed, popmuted] (0,25) -- (32,25) node[right] {$\theta_a=25^{\circ}\text{C}$};
% Points
\fill[popblue] (0,90) circle (2pt) node[left] {$90$};
\fill[popblue] (5,65) circle (2pt) node[above right] {$65$};
\fill[popblue] (15,40) circle (2pt) node[above right] {$40$};
\fill[popblue] (26.4,30) circle (2pt) node[below right] {$30$};
% Curve
\draw[popcurve, thick, samples=200] plot (\x, {25+65*exp(-0.097*\x)});
% Vertical lines for times
\draw[dotted, popgray] (5,0) -- (5,65);
\draw[dotted, popgray] (15,0) -- (15,40);
\draw[dotted, popgray] (26.4,0) -- (26.4,30);
\end{tikzpicture}
\legende{Refroidissement du café : courbe $\theta(t)=25+65e^{-kt}$, points caractéristiques.}
\end{popfigure}

\begin{attention}[Les 5 erreurs qui coûtent des points au Bac]
\begin{enumerate}
\item \textbf{Oublier que $e^{x}>0$ pour tout réel $x$.} Conséquences : affirmer que $e^{x}=-3$ a une solution, ou accepter une racine $X\leq 0$ après $X=e^{x}$. \textbf{Toute racine $\leq 0$ doit être explicitement rejetée.}
\item \textbf{Dériver $e^{u}$ en oubliant $u'$.} Dérivée de $e^{3x}$ est $3e^{3x}$, pas $e^{3x}$. Primitive de $e^{3x}$ est $\frac{1}{3}e^{3x}$, pas $e^{3x}$. \textbf{Vérifie toujours en dérivant.}
\item \textbf{Changer le sens d'une inégalité à tort.} $\exp$ est \textbf{strictement croissante} : $e^{A}<e^{B}\Leftrightarrow A<B$ (sens conservé). $\ln$ aussi. On ne « retourne » l'inégalité qu'en multipliant/divisant par un \textbf{négatif}.
\item \textbf{Confondre $e^{a+b}$ et $e^{a}+e^{b}$, ou $(e^{a})^{n}$ et $e^{a^{n}}$.} Bonnes formules : $e^{a+b}=e^{a}e^{b}$, $(e^{a})^{n}=e^{na}$. Une erreur ici fausse tout l'exercice.
\item \textbf{Conclure une modélisation sans phrase, sans unité, sans arrondi.} Résultat brut $t\approx 26,4$ $\to$ écrire : « Le café sera à $30^{\circ}\text{C}$ au bout d'environ $26$ minutes. » Vérifier le signe de $k$ (négatif pour refroidissement/décroissance).
\end{enumerate}
\end{attention}

\newpage

\coursec{Exercices}

\courssub{Je vérifie mes connaissances}

\begin{exercice}[QCM : une seule bonne réponse]
Pour chaque question, choisir la bonne réponse (aucune justification n'est demandée).
\begin{enumerate}
\item Pour tout réel $x$, $\mathrm{e}^{x+2}$ est égal à :
\begin{enumerate}
\item[a)] $\mathrm{e}^x + \mathrm{e}^2$ \qquad b) $\mathrm{e}^x \times \mathrm{e}^2$ \qquad c) $2\mathrm{e}^x$ \qquad d) $(\mathrm{e}^x)^2$
\end{enumerate}
\item L'équation $\ln x = 3$ a pour solution :
\begin{enumerate}
\item[a)] $x = \ln 3$ \qquad b) $x = 3\mathrm{e}$ \qquad c) $x = \mathrm{e}^3$ \qquad d) $x = \dfrac{3}{\mathrm{e}}$
\end{enumerate}
\item La dérivée de la fonction $x \longmapsto \mathrm{e}^{3x}$ est :
\begin{enumerate}
\item[a)] $x \longmapsto \mathrm{e}^{3x}$ \qquad b) $x \longmapsto 3\mathrm{e}^{3x}$ \qquad c) $x \longmapsto 3x\,\mathrm{e}^{3x-1}$ \qquad d) $x \longmapsto \mathrm{e}^{3}$
\end{enumerate}
\item $\displaystyle\lim_{x \to -\infty} \mathrm{e}^{x}$ est égale à :
\begin{enumerate}
\item[a)] $-\infty$ \qquad b) $+\infty$ \qquad c) $0$ \qquad d) $1$
\end{enumerate}
\end{enumerate}
\end{exercice}

\begin{exercice}[Vrai ou faux, en justifiant]
Pour chacune des affirmations suivantes, dire si elle est vraie ou fausse, puis justifier la réponse par une propriété du cours ou un contre-exemple.
\begin{enumerate}
\item Pour tous réels $a$ et $b$ : $\mathrm{e}^{a}\times \mathrm{e}^{b} = \mathrm{e}^{ab}$.
\item Pour tout réel $x$ : $\mathrm{e}^{-x} = \dfrac{1}{\mathrm{e}^{x}}$.
\item Pour tout réel $x$ strictement positif : $\mathrm{e}^{\ln x} = x$.
\item La fonction exponentielle est strictement croissante sur $\mathbb{R}$.
\item Pour tout réel $x$ : $\ln(\mathrm{e}^{x}) = x$.
\end{enumerate}
\end{exercice}

\begin{exercice}[Calculs directs]
Simplifier chacune des expressions suivantes, où $x$ désigne un réel quelconque :
\[ A = \mathrm{e}^{x}\times \mathrm{e}^{-x} \;;\qquad B = \dfrac{\mathrm{e}^{5x}}{\mathrm{e}^{2x}} \;;\qquad C = \left(\mathrm{e}^{x}\right)^{3}\times \mathrm{e}^{-2x} \;;\qquad D = \mathrm{e}^{\ln 7} + \ln\left(\mathrm{e}^{-4}\right). \]
\end{exercice}

\begin{exercice}[Lecture de cours]
\begin{enumerate}
\item Donner la définition de la fonction exponentielle népérienne comme fonction réciproque du logarithme népérien, en précisant son ensemble de définition et son ensemble image.
\item Compléter : $\ln x = y \iff x = \ldots$ (avec $x > 0$).
\item Donner, sans démonstration, une primitive sur $\mathbb{R}$ de la fonction $x \longmapsto u'(x)\,\mathrm{e}^{u(x)}$, où $u$ est une fonction dérivable.
\end{enumerate}
\end{exercice}

\courssub{Je m'entraîne}

\begin{exercice}[Simplification d'expressions]
Écrire le plus simplement possible les expressions suivantes, où $x$ désigne un réel :
\[ A = \dfrac{\mathrm{e}^{2x}\times \mathrm{e}^{3-x}}{\left(\mathrm{e}^{x+1}\right)^{2}} \;;\qquad B = \mathrm{e}^{\ln 5} + \ln\left(\mathrm{e}^{-3}\right) - \mathrm{e}^{2\ln 3}. \]
\end{exercice}

\begin{methode}[Résolution d'équations avec $\mathrm{e}^x$]
\begin{enumerate}
\item Identifier la structure : équation du type $\mathrm{e}^{u(x)} = k$ ($k>0$), ou polynôme en $\mathrm{e}^x$ (positon $X=\mathrm{e}^x>0$), ou mélange $\ln$ et $\mathrm{e}^x$.
\item Si $\mathrm{e}^{u(x)} = k$ avec $k>0$ : passer au logarithme $\iff u(x) = \ln k$.
\item Si polynôme en $\mathrm{e}^x$ : poser $X = \mathrm{e}^x$ ($X>0$), résoudre l'équation polynomiale, ne garder que les racines $X>0$, puis revenir à $x = \ln X$.
\item Si logarithmes : déterminer d'abord l'ensemble de définition (arguments $>0$), utiliser les propriétés $\ln a + \ln b = \ln(ab)$, puis passer à l'exponentielle.
\item Vérifier que les solutions trouvées appartiennent à l'ensemble de définition.
\end{enumerate}
\end{methode}

\begin{exercice}[Équations]
Résoudre dans $\mathbb{R}$ les équations suivantes :
\begin{enumerate}
\item $\mathrm{e}^{x^{2}-3x} = 1$ ;
\item $\mathrm{e}^{2x} - 4\mathrm{e}^{x} + 3 = 0$ (on pourra poser $X = \mathrm{e}^{x}$) ;
\item $\ln(x+1) + \ln(x-1) = \ln 8$ (on précisera d'abord l'ensemble de résolution).
\end{enumerate}
\end{exercice}

\begin{methode}[Résolution d'inéquations avec $\mathrm{e}^x$ ou $\ln x$]
\begin{enumerate}
\item Pour $\mathrm{e}^{u(x)} \lesseqgtr k$ : 
      \begin{itemize}
      \item si $k \le 0$ : inéquation toujours vraie (si $>$) ou fausse (si $\le$) car $\mathrm{e}^{u(x)} > 0$.
      \item si $k > 0$ : équivalent à $u(x) \lesseqgtr \ln k$ (car $\ln$ est croissante).
      \end{itemize}
\item Pour $\ln(u(x)) \lesseqgtr k$ : 
      \begin{itemize}
      \item Déterminer la condition $u(x) > 0$ (ensemble de définition).
      \item L'inéquation est équivalente au système $\begin{cases} u(x) > 0 \\ u(x) \lesseqgtr \mathrm{e}^k \end{cases}$ (car $\exp$ est croissante).
      \end{itemize}
\item Pour les inéquations polynomiales en $\mathrm{e}^x$ : poser $X = \mathrm{e}^x > 0$, résoudre en $X$, puis revenir à $x$.
\end{enumerate}
\end{methode}

\begin{exercice}[Inéquations]
\begin{enumerate}
\item Résoudre dans $\mathbb{R}$ l'inéquation $\mathrm{e}^{2x} - 2\mathrm{e}^{x} > 0$ (on pourra factoriser par $\mathrm{e}^{x}$).
\item Résoudre dans $\mathbb{R}$ l'inéquation $\ln(2x-4) \leq 1$, en précisant au préalable l'ensemble de résolution.
\end{enumerate}
\end{exercice}

\begin{methode}[Recherche de primitives de la forme $u' \mathrm{e}^u$ ou $u' u^n$]
\begin{enumerate}
\item Identifier $u(x)$ et calculer $u'(x)$.
\item Vérifier que l'expression à intégrer s'écrit $u'(x) \times \mathrm{e}^{u(x)}$ (ou $u'(x) \times [u(x)]^n$).
\item Une primitive est alors $\mathrm{e}^{u(x)}$ (ou $\frac{[u(x)]^{n+1}}{n+1}$ pour $n \neq -1$).
\item Ne pas oublier la constante d'intégration $+C$ si on cherche l'ensemble des primitives.
\end{enumerate}
\end{methode}

\begin{exercice}[Primitives]
Déterminer une primitive sur $\mathbb{R}$ de chacune des fonctions suivantes :
\begin{enumerate}
\item $f(x) = (2x+1)\,\mathrm{e}^{x^{2}+x}$ ;
\item $g(x) = 3\mathrm{e}^{3x-1}$ ;
\item $h(x) = \mathrm{e}^{x}\left(1 + \mathrm{e}^{x}\right)^{2}$ (on pourra reconnaître la forme $u'u^{n}$).
\end{enumerate}
\end{exercice}

\courssub{Je me prépare au Bac}

\begin{exercice}[Étude complète d'une fonction — type Bac C/D]
On considère la fonction $f$ définie sur $\mathbb{R}$ par :
\[ f(x) = \left(x^{2} - 2x + 2\right)\mathrm{e}^{-x} \]
et on note $(\mathcal{C})$ sa courbe représentative dans un repère orthonormé $(O\,;\vec{\imath},\vec{\jmath})$ (unité : \SI{2}{cm}).
\begin{enumerate}
\item \begin{enumerate}
\item Calculer $\displaystyle\lim_{x \to -\infty} f(x)$.
\item Montrer que pour tout réel $x$, $f(x) = \dfrac{x^{2}}{\mathrm{e}^{x}} - \dfrac{2x}{\mathrm{e}^{x}} + \dfrac{2}{\mathrm{e}^{x}}$, puis en déduire $\displaystyle\lim_{x \to +\infty} f(x)$. Quelle interprétation graphique peut-on en donner ?
\end{enumerate}
\item \begin{enumerate}
\item Montrer que pour tout réel $x$ : $f'(x) = \left(-x^{2} + 4x - 4\right)\mathrm{e}^{-x}$.
\item Vérifier que $f'(x) = -(x-2)^{2}\mathrm{e}^{-x}$, puis étudier le signe de $f'(x)$ sur $\mathbb{R}$.
\item Dresser le tableau de variations complet de $f$.
\end{enumerate}
\item Déterminer une équation de la tangente $(T)$ à $(\mathcal{C})$ au point d'abscisse $0$.
\item Tracer $(T)$ et $(\mathcal{C})$ dans le repère $(O\,;\vec{\imath},\vec{\jmath})$.
\item Déterminer une primitive $F$ de $f$ de la forme $F(x) = (ax^{2}+bx+c)\mathrm{e}^{-x}$, où $a$, $b$ et $c$ sont des réels à déterminer.
\end{enumerate}
\end{exercice}

\begin{exercice}[Système, limites et aire — type Bac C/D/TI]
\textbf{Partie A.} Résoudre dans $\mathbb{R}^{2}$, avec $x > 0$ et $y > 0$, le système :
\[ \begin{cases} \mathrm{e}^{x}\times \mathrm{e}^{y} = \mathrm{e}^{5} \\ \ln x + \ln y = \ln 6 \end{cases} \]

\textbf{Partie B.} Déterminer les limites suivantes, en justifiant soigneusement :
\begin{enumerate}
\item $\displaystyle\lim_{x \to +\infty} \left(\mathrm{e}^{x} - x\right)$ ;
\item $\displaystyle\lim_{x \to +\infty} x\,\mathrm{e}^{\frac{1}{x}}$ ;
\item $\displaystyle\lim_{x \to 0} \dfrac{\mathrm{e}^{2x}-1}{x}$ (on pourra poser $X = 2x$ et utiliser $\displaystyle\lim_{X \to 0}\dfrac{\mathrm{e}^{X}-1}{X} = 1$).
\end{enumerate}

\textbf{Partie C.} Soit $g$ la fonction définie sur $\mathbb{R}$ par $g(x) = \mathrm{e}^{x}$, de courbe représentative $(\Gamma)$ dans un repère orthonormé d'unité \SI{2}{cm}.
\begin{enumerate}
\item Déterminer une primitive de $g$ sur $\mathbb{R}$.
\item Calculer, en unités d'aire puis en \si{cm^{2}}, l'aire du domaine délimité par $(\Gamma)$, l'axe des abscisses et les droites d'équations $x = 0$ et $x = 1$. On donnera la valeur exacte, puis une valeur approchée à $10^{-2}$ près (on rappelle que $\mathrm{e} \approx \num{2,718}$).
\end{enumerate}
\end{exercice}

\begin{exercice}[Problème de modélisation : refroidissement d'un moteur — type Bac C/D/E/TI]
Après une course à travers Yaoundé, un conducteur de moto-taxi gare sa moto à l'ombre. La température du moteur, exprimée en degrés Celsius, est modélisée en fonction du temps $t$ (en minutes) écoulé depuis l'arrêt par :
\[ T(t) = 25 + 75\,\mathrm{e}^{-0{,}08t}, \qquad t \geq 0. \]
La température ambiante ce jour-là est de $25\,^{\circ}\mathrm{C}$.
\begin{enumerate}
\item Quelle est la température du moteur à l'instant de l'arrêt ($t = 0$) ?
\item \begin{enumerate}
\item Calculer $T'(t)$ et étudier son signe sur $[0\,;+\infty[$.
\item Déterminer $\displaystyle\lim_{t \to +\infty} T(t)$. Interpréter ce résultat dans le contexte de la situation.
\item Dresser le tableau de variations de $T$ sur $[0\,;+\infty[$.
\end{enumerate}
\item Le mécanicien du quartier affirme qu'on peut toucher le moteur sans risque dès que sa température descend sous $40\,^{\circ}\mathrm{C}$.
\begin{enumerate}
\item Montrer que résoudre l'inéquation $T(t) < 40$ revient à résoudre $\mathrm{e}^{-0{,}08t} < \dfrac{1}{5}$.
\item En déduire, à l'aide de la fonction $\ln$, au bout de combien de minutes (arrondir à la minute supérieure) on peut toucher le moteur sans risque. On donne $\ln 5 \approx \num{1,609}$.
\end{enumerate}
\item Calculer la température du moteur au bout de $30$ minutes (arrondir au dixième de degré). On donne $\mathrm{e}^{-2{,}4} \approx \num{0,0907}$.
\end{enumerate}
\end{exercice}



\newpage

\coursec{Corrigés des exercices}

\begin{corrige}[Exercice 1 — QCM]
\begin{enumerate}
\item \textbf{Réponse b).} Propriété algébrique fondamentale : pour tous réels $a$ et $b$, $\mathrm{e}^{a+b} = \mathrm{e}^{a}\times\mathrm{e}^{b}$. Avec $a=x$ et $b=2$ : $\mathrm{e}^{x+2} = \mathrm{e}^{x}\times\mathrm{e}^{2}$.
\item \textbf{Réponse c).} Par définition de l'exponentielle comme réciproque du logarithme népérien : $\ln x = y \iff x = \mathrm{e}^{y}$. Donc $\ln x = 3 \iff x = \mathrm{e}^{3}$.
\item \textbf{Réponse b).} Dérivée d'une composée : $\left(\mathrm{e}^{u}\right)' = u'\,\mathrm{e}^{u}$. Ici $u(x)=3x$, $u'(x)=3$, d'où la dérivée $x \longmapsto 3\mathrm{e}^{3x}$.
\item \textbf{Réponse c).} Limite usuelle du cours : $\displaystyle\lim_{x \to -\infty} \mathrm{e}^{x} = 0$ (l'axe des abscisses est asymptote horizontale à la courbe de $\exp$ en $-\infty$).
\end{enumerate}
\textbf{Points de vigilance :} $\mathrm{e}^{x+2}$ n'est \emph{pas} une somme (l'exponentielle ne transforme pas l'addition en addition, mais en produit) ; la dérivée de $\mathrm{e}^{3x}$ n'est pas $\mathrm{e}^{3x}$ : ne jamais oublier le facteur $u'$.
\end{corrige}

\begin{corrige}[Exercice 2 — Vrai ou faux]
\begin{enumerate}
\item \textbf{Faux.} La propriété correcte est $\mathrm{e}^{a}\times\mathrm{e}^{b} = \mathrm{e}^{a+b}$. Contre-exemple : pour $a=b=1$, $\mathrm{e}^{1}\times\mathrm{e}^{1} = \mathrm{e}^{2}$, alors que $\mathrm{e}^{ab} = \mathrm{e}^{1} = \mathrm{e}$. Or $\mathrm{e}^{2} \neq \mathrm{e}$.
\item \textbf{Vrai.} C'est la propriété $\mathrm{e}^{-a} = \dfrac{1}{\mathrm{e}^{a}}$, qui découle de $\mathrm{e}^{a}\times\mathrm{e}^{-a} = \mathrm{e}^{0} = 1$.
\item \textbf{Vrai.} Les fonctions $\exp$ et $\ln$ sont bijections réciproques : pour tout $x > 0$, $\mathrm{e}^{\ln x} = x$.
\item \textbf{Vrai.} La fonction $\exp$ est dérivable sur $\mathbb{R}$ et $(\exp)'(x) = \mathrm{e}^{x} > 0$ pour tout réel $x$ : elle est donc strictement croissante sur $\mathbb{R}$.
\item \textbf{Vrai.} Par réciprocité : pour tout réel $x$, $\ln\left(\mathrm{e}^{x}\right) = x$.
\end{enumerate}
\textbf{Points de vigilance :} aux affirmations 3 et 5, les ensembles ne sont pas les mêmes : $\mathrm{e}^{\ln x} = x$ n'a de sens que pour $x > 0$, tandis que $\ln(\mathrm{e}^{x}) = x$ est valable pour tout réel $x$. Citer le domaine fait partie de la justification.
\end{corrige}

\begin{corrige}[Exercice 3 — Calculs directs]
\begin{align*}
A &= \mathrm{e}^{x}\times\mathrm{e}^{-x} = \mathrm{e}^{x+(-x)} = \mathrm{e}^{0} = 1. \\[0.3em]
B &= \dfrac{\mathrm{e}^{5x}}{\mathrm{e}^{2x}} = \mathrm{e}^{5x-2x} = \mathrm{e}^{3x}. \\[0.3em]
C &= \left(\mathrm{e}^{x}\right)^{3}\times\mathrm{e}^{-2x} = \mathrm{e}^{3x}\times\mathrm{e}^{-2x} = \mathrm{e}^{3x-2x} = \mathrm{e}^{x}. \\[0.3em]
D &= \mathrm{e}^{\ln 7} + \ln\left(\mathrm{e}^{-4}\right) = 7 + (-4) = 3.
\end{align*}
\textbf{Points de vigilance :} pour $C$, on utilise d'abord $\left(\mathrm{e}^{a}\right)^{n} = \mathrm{e}^{na}$ \emph{avant} la règle du produit ; pour $D$, on reconnaît les deux relations de réciproque $\mathrm{e}^{\ln 7} = 7$ (car $7>0$) et $\ln(\mathrm{e}^{-4}) = -4$.
\end{corrige}

\begin{corrige}[Exercice 4 — Lecture de cours]
\begin{enumerate}
\item La fonction logarithme népérien $\ln : \left]0\,;+\infty\right[ \to \mathbb{R}$ est une bijection. Sa bijection réciproque est appelée \cle{fonction exponentielle népérienne}, notée $\exp$ ou $x \mapsto \mathrm{e}^{x}$. Elle est définie sur $\mathbb{R}$ et à valeurs dans $\left]0\,;+\infty\right[$ :
\[ \exp : \mathbb{R} \longrightarrow \left]0\,;+\infty\right[. \]
\item $\ln x = y \iff x = \mathrm{e}^{y}$ (avec $x > 0$).
\item Une primitive sur $\mathbb{R}$ de $x \longmapsto u'(x)\,\mathrm{e}^{u(x)}$ est $x \longmapsto \mathrm{e}^{u(x)}$, puisque $\left(\mathrm{e}^{u}\right)' = u'\,\mathrm{e}^{u}$.
\end{enumerate}
\textbf{Points de vigilance :} dans la définition, il faut impérativement préciser les deux ensembles (départ $\mathbb{R}$, arrivée $]0\,;+\infty[$) : c'est ce qui traduit la bijectivité.
\end{corrige}

\begin{corrige}[Exercice 5 — Simplification d'expressions]
Pour $A$, on applique successivement $\mathrm{e}^{a}\mathrm{e}^{b} = \mathrm{e}^{a+b}$, $\left(\mathrm{e}^{a}\right)^{2} = \mathrm{e}^{2a}$ et $\dfrac{\mathrm{e}^{a}}{\mathrm{e}^{b}} = \mathrm{e}^{a-b}$ :
\begin{align*}
A &= \dfrac{\mathrm{e}^{2x}\times\mathrm{e}^{3-x}}{\left(\mathrm{e}^{x+1}\right)^{2}}
   = \dfrac{\mathrm{e}^{2x+3-x}}{\mathrm{e}^{2(x+1)}}
   = \dfrac{\mathrm{e}^{x+3}}{\mathrm{e}^{2x+2}}
   = \mathrm{e}^{(x+3)-(2x+2)}
   = \mathrm{e}^{-x+1} = \mathrm{e}^{1-x}.
\end{align*}
Pour $B$, on utilise les relations de réciproque et $2\ln 3 = \ln 9$ :
\begin{align*}
B &= \mathrm{e}^{\ln 5} + \ln\left(\mathrm{e}^{-3}\right) - \mathrm{e}^{2\ln 3}
   = 5 + (-3) - \mathrm{e}^{\ln 9}
   = 2 - 9 = -7.
\end{align*}
\textbf{Points de vigilance :} dans $A$, ne pas écrire $\left(\mathrm{e}^{x+1}\right)^{2} = \mathrm{e}^{x^{2}+1}$ : on \emph{multiplie} l'exposant par $2$, on ne l'élève pas au carré. Dans $B$, $\mathrm{e}^{2\ln 3} = \mathrm{e}^{\ln 9} = 9$ et non $6$.
\end{corrige}

\begin{corrige}[Exercice 6 — Équations]
\begin{enumerate}
\item Comme $1 = \mathrm{e}^{0}$ et que $\exp$ est injective :
\[ \mathrm{e}^{x^{2}-3x} = \mathrm{e}^{0} \iff x^{2} - 3x = 0 \iff x(x-3) = 0 \iff x = 0 \ \text{ou}\ x = 3. \]
\[ \mathcal{S} = \{0\,;\,3\}. \]
\item Posons $X = \mathrm{e}^{x}$, avec la condition $X > 0$. Alors $\mathrm{e}^{2x} = \left(\mathrm{e}^{x}\right)^{2} = X^{2}$ et l'équation devient :
\[ X^{2} - 4X + 3 = 0. \]
Discriminant : $\Delta = 16 - 12 = 4$, racines $X_{1} = \dfrac{4-2}{2} = 1$ et $X_{2} = \dfrac{4+2}{2} = 3$. Les deux racines sont strictement positives, donc toutes deux conviennent.
\[ X = 1 \iff \mathrm{e}^{x} = 1 \iff x = 0 \;;\qquad X = 3 \iff \mathrm{e}^{x} = 3 \iff x = \ln 3. \]
\[ \mathcal{S} = \{0\,;\,\ln 3\}. \]
\item \textbf{Ensemble de résolution :} il faut $x+1 > 0$ \emph{et} $x-1 > 0$, c'est-à-dire $x > 1$. On résout donc sur $\mathcal{D} = \left]1\,;+\infty\right[$. Sur $\mathcal{D}$ :
\[ \ln(x+1) + \ln(x-1) = \ln 8 \iff \ln\big((x+1)(x-1)\big) = \ln 8 \iff \ln(x^{2}-1) = \ln 8. \]
Par injectivité de $\ln$ : $x^{2} - 1 = 8 \iff x^{2} = 9 \iff x = 3$ ou $x = -3$. Or $-3 \notin \mathcal{D}$ : on rejette cette valeur.
\[ \mathcal{S} = \{3\}. \]
\end{enumerate}
\textbf{Points de vigilance :} à la question 2, la condition $X > 0$ doit être écrite dès le changement de variable ; à la question 3, la recherche de l'ensemble de résolution est \emph{indispensable} et la conclusion doit mentionner le rejet de $-3$. C'est une exigence de rédaction du Bac.
\end{corrige}

\begin{corrige}[Exercice 7 — Inéquations]
\begin{enumerate}
\item On factorise par $\mathrm{e}^{x}$ :
\[ \mathrm{e}^{2x} - 2\mathrm{e}^{x} > 0 \iff \mathrm{e}^{x}\left(\mathrm{e}^{x} - 2\right) > 0. \]
Or pour tout réel $x$, $\mathrm{e}^{x} > 0$ : le produit a donc le signe de $\mathrm{e}^{x} - 2$.
\[ \mathrm{e}^{x} - 2 > 0 \iff \mathrm{e}^{x} > 2 \iff x > \ln 2 \quad (\ln\ \text{est strictement croissante}). \]
\[ \mathcal{S} = \left]\ln 2\,;+\infty\right[. \]
\item \textbf{Ensemble de résolution :} $2x - 4 > 0 \iff x > 2$. On résout sur $\mathcal{D} = \left]2\,;+\infty\right[$. Sur $\mathcal{D}$, en écrivant $1 = \ln \mathrm{e}$ :
\[ \ln(2x-4) \leq \ln \mathrm{e} \iff 2x - 4 \leq \mathrm{e} \iff x \leq \dfrac{\mathrm{e}+4}{2}, \]
car $\ln$ est strictement croissante (le sens de l'inégalité est conservé). Comme $\dfrac{\mathrm{e}+4}{2} \approx \num{3,36} > 2$, la contrainte de domaine est compatible :
\[ \mathcal{S} = \left]2\,;\ \dfrac{\mathrm{e}+4}{2}\right]. \]
\end{enumerate}
\textbf{Points de vigilance :} la croissance de $\ln$ et de $\exp$ justifie la \emph{conservation} du sens des inégalités : il faut le dire. À la question 2, la borne $2$ est exclue (condition de définition stricte) alors que la borne $\frac{\mathrm{e}+4}{2}$ est incluse (inégalité large).
\end{corrige}

\begin{corrige}[Exercice 8 — Primitives]
\begin{enumerate}
\item Posons $u(x) = x^{2} + x$ ; alors $u'(x) = 2x + 1$. La fonction s'écrit exactement $f = u'\,\mathrm{e}^{u}$. Une primitive sur $\mathbb{R}$ est donc :
\[ F(x) = \mathrm{e}^{x^{2}+x}. \]
\emph{Vérification :} $F'(x) = (2x+1)\mathrm{e}^{x^{2}+x} = f(x)$. 
\item Posons $u(x) = 3x - 1$ ; alors $u'(x) = 3$. On a $g = u'\,\mathrm{e}^{u}$, donc une primitive sur $\mathbb{R}$ est :
\[ G(x) = \mathrm{e}^{3x-1}. \]
\emph{Vérification :} $G'(x) = 3\mathrm{e}^{3x-1} = g(x)$. 
\item Posons $u(x) = 1 + \mathrm{e}^{x}$ ; alors $u'(x) = \mathrm{e}^{x}$ et $h = u'\,u^{2}$ : c'est la forme $u'u^{n}$ avec $n = 2$, de primitive $\dfrac{u^{3}}{3}$. Une primitive sur $\mathbb{R}$ est :
\[ H(x) = \dfrac{1}{3}\left(1 + \mathrm{e}^{x}\right)^{3}. \]
\emph{Vérification :} $H'(x) = \dfrac{1}{3}\times 3\mathrm{e}^{x}\left(1+\mathrm{e}^{x}\right)^{2} = h(x)$. 
\end{enumerate}
\textbf{Points de vigilance :} la démarche exigible est : identifier $u$, calculer $u'$, \emph{constater} que l'expression est exactement $u'\mathrm{e}^{u}$ (à un facteur constant près qu'on ajuste), puis conclure. La vérification par dérivation est un excellent réflexe au Bac.
\end{corrige}

\begin{corrige}[Exercice 9 — Étude complète d'une fonction]
\textbf{Barème indicatif (sur 10 points) :} question 1 : 2 pts ; question 2 : 3 pts ; question 3 : 1 pt ; question 4 : 2 pts ; question 5 : 2 pts.

\begin{enumerate}
\item \begin{enumerate}
\item On a $\displaystyle\lim_{x \to -\infty}\left(x^{2} - 2x + 2\right) = +\infty$ (terme dominant $x^{2}$) et, par composition avec $X = -x \to +\infty$, $\displaystyle\lim_{x \to -\infty} \mathrm{e}^{-x} = +\infty$. Par produit :
\[ \lim_{x \to -\infty} f(x) = +\infty. \]
\item Pour tout réel $x$, $\mathrm{e}^{-x} = \dfrac{1}{\mathrm{e}^{x}}$, donc :
\[ f(x) = \dfrac{x^{2} - 2x + 2}{\mathrm{e}^{x}} = \dfrac{x^{2}}{\mathrm{e}^{x}} - \dfrac{2x}{\mathrm{e}^{x}} + \dfrac{2}{\mathrm{e}^{x}}. \]
Par croissances comparées, l'exponentielle l'emporte sur toute puissance de $x$ : $\displaystyle\lim_{x \to +\infty} \dfrac{x^{2}}{\mathrm{e}^{x}} = 0$, $\displaystyle\lim_{x \to +\infty} \dfrac{x}{\mathrm{e}^{x}} = 0$ et $\displaystyle\lim_{x \to +\infty} \dfrac{1}{\mathrm{e}^{x}} = 0$. Par somme :
\[ \lim_{x \to +\infty} f(x) = 0. \]
\textbf{Interprétation graphique :} la droite d'équation $y = 0$ (l'axe des abscisses) est \cle{asymptote horizontale} à $(\mathcal{C})$ au voisinage de $+\infty$.
\end{enumerate}
\item \begin{enumerate}
\item $f$ est dérivable sur $\mathbb{R}$ comme produit de fonctions dérivables. Avec $u(x) = x^{2}-2x+2$ et $v(x) = \mathrm{e}^{-x}$ : $u'(x) = 2x-2$ et $v'(x) = -\mathrm{e}^{-x}$. Alors :
\begin{align*}
f'(x) &= u'(x)v(x) + u(x)v'(x) = (2x-2)\mathrm{e}^{-x} - \left(x^{2}-2x+2\right)\mathrm{e}^{-x} \\
&= \mathrm{e}^{-x}\left(2x - 2 - x^{2} + 2x - 2\right) = \left(-x^{2} + 4x - 4\right)\mathrm{e}^{-x}.
\end{align*}
\item On factorise : $-x^{2} + 4x - 4 = -\left(x^{2} - 4x + 4\right) = -(x-2)^{2}$, d'où $f'(x) = -(x-2)^{2}\mathrm{e}^{-x}$.\\
Pour tout réel $x$ : $\mathrm{e}^{-x} > 0$ et $(x-2)^{2} \geq 0$, donc $f'(x) \leq 0$, avec $f'(x) = 0 \iff x = 2$. La dérivée est strictement négative sauf en un point isolé : $f$ est \textbf{strictement décroissante sur $\mathbb{R}$}.
\item On calcule $f(2) = (4 - 4 + 2)\mathrm{e}^{-2} = 2\mathrm{e}^{-2} \approx \num{0,27}$. Tableau de variations :
\begin{center}
\begin{tabular}{|c|ccccc|}
\hline
$x$ & $-\infty$ & & $2$ & & $+\infty$ \\
\hline
$f'(x)$ & & $-$ & $0$ & $-$ & \\
\hline
$f$ & $+\infty$ & $\searrow$ & $2\mathrm{e}^{-2}$ & $\searrow$ & $0$ \\
\hline
\end{tabular}
\end{center}
\end{enumerate}
\item $f(0) = 2\mathrm{e}^{0} = 2$ et $f'(0) = -(0-2)^{2}\mathrm{e}^{0} = -4$. La tangente $(T)$ au point $A(0\,;\,2)$ a pour équation :
\[ y = f'(0)(x - 0) + f(0) \iff \boxed{y = -4x + 2}. \]
\item Voir la figure ci-dessous : $(\mathcal{C})$ passe par $A(0\,;\,2)$, admet $(T)$ pour tangente en $A$, décroît et se rapproche de l'axe des abscisses en $+\infty$.
\begin{popfigure}
\begin{tikzpicture}
\begin{axis}[width=0.85\linewidth, height=6cm, axis lines=middle, xlabel={$x$}, ylabel={$y$}, xmin=-1.2, xmax=5, ymin=-1, ymax=6.5, xtick={1,2,3,4}, ytick={1,2,3,4,5,6}, legend style={at={(0.97,0.97)},anchor=north east,font=\small}]
\addplot[popcurve,domain=-1:5,samples=200]{(x^2-2*x+2)*exp(-x)};
\addlegendentry{$(\mathcal{C})$}
\addplot[poporange,thick,dashed,domain=-0.9:1.6,samples=2]{-4*x+2};
\addlegendentry{$(T)$}
\addplot[popmuted,domain=-1:5,samples=2]{0};
\node[circle,fill=popblue,inner sep=1.5pt,label={above right:{\small $A(0;2)$}}] at (axis cs:0,2) {};
\end{axis}
\end{tikzpicture}
\legende{Courbe $(\mathcal{C})$, tangente $(T)$ en $A$ et asymptote $y=0$}
\end{popfigure}
\item Soit $F(x) = (ax^{2}+bx+c)\mathrm{e}^{-x}$. $F$ est dérivable sur $\mathbb{R}$ et :
\begin{align*}
F'(x) &= (2ax+b)\mathrm{e}^{-x} - \left(ax^{2}+bx+c\right)\mathrm{e}^{-x} \\
&= \left(-ax^{2} + (2a-b)x + (b-c)\right)\mathrm{e}^{-x}.
\end{align*}
$F$ est une primitive de $f$ si et seulement si, pour tout réel $x$, $F'(x) = f(x)$, c'est-à-dire par identification des coefficients (car $\mathrm{e}^{-x} \neq 0$) :
\[ \begin{cases} -a = 1 \\ 2a - b = -2 \\ b - c = 2 \end{cases} \iff \begin{cases} a = -1 \\ b = 2a + 2 = 0 \\ c = b - 2 = -2. \end{cases} \]
Une primitive de $f$ sur $\mathbb{R}$ est donc :
\[ F(x) = \left(-x^{2} - 2\right)\mathrm{e}^{-x}. \]
\emph{Vérification :} $F'(x) = -2x\,\mathrm{e}^{-x} + \left(x^{2}+2\right)\mathrm{e}^{-x} = \left(x^{2}-2x+2\right)\mathrm{e}^{-x} = f(x)$. 
\end{enumerate}
\textbf{Points de vigilance :} en 1.b, citer explicitement le théorème des croissances comparées ; en 2.b, le signe de $f'$ repose sur les deux arguments « carré positif » et « exponentielle strictement positive » ; en 2.c, l'annulation de $f'$ en $2$ sans changement de signe ne crée \emph{pas} d'extremum (point d'inflexion à tangente horizontale) ; en 5, l'identification est licite parce que deux polynômes égaux pour tout $x$ ont les mêmes coefficients.
\end{corrige}

\begin{corrige}[Exercice 10 — Système, limites et aire]
\textbf{Barème indicatif (sur 10 points) :} Partie A : 3 pts ; Partie B : 3 pts (1 pt par limite) ; Partie C : 4 pts.

\textbf{Partie A.} Avec $x > 0$ et $y > 0$, on utilise $\mathrm{e}^{x}\mathrm{e}^{y} = \mathrm{e}^{x+y}$ et $\ln x + \ln y = \ln(xy)$ :
\[ \begin{cases} \mathrm{e}^{x+y} = \mathrm{e}^{5} \\ \ln(xy) = \ln 6 \end{cases} \iff \begin{cases} x + y = 5 \\ xy = 6 \end{cases} \quad (\text{injectivité de } \exp \text{ et de } \ln). \]
$x$ et $y$ sont donc les racines de l'équation $X^{2} - 5X + 6 = 0$. Discriminant : $\Delta = 25 - 24 = 1$ ; racines : $X_{1} = \dfrac{5-1}{2} = 2$ et $X_{2} = \dfrac{5+1}{2} = 3$, toutes deux strictement positives (conditions respectées).
\[ \mathcal{S} = \{(2\,;\,3)\,;\,(3\,;\,2)\}. \]

\textbf{Partie B.}
\begin{enumerate}
\item Pour tout $x > 0$, $\mathrm{e}^{x} - x = x\left(\dfrac{\mathrm{e}^{x}}{x} - 1\right)$. Par croissances comparées, $\displaystyle\lim_{x \to +\infty} \dfrac{\mathrm{e}^{x}}{x} = +\infty$, donc la parenthèse tend vers $+\infty$ et, par produit :
\[ \lim_{x \to +\infty} \left(\mathrm{e}^{x} - x\right) = +\infty. \]
\item Par composition, $\displaystyle\lim_{x \to +\infty} \dfrac{1}{x} = 0$ et $\exp$ est continue en $0$, donc $\displaystyle\lim_{x \to +\infty} \mathrm{e}^{\frac{1}{x}} = \mathrm{e}^{0} = 1$. Par produit avec $x \to +\infty$ :
\[ \lim_{x \to +\infty} x\,\mathrm{e}^{\frac{1}{x}} = +\infty. \]
\item Posons $X = 2x$ ; quand $x \to 0$, $X \to 0$, et $x = \dfrac{X}{2}$ :
\[ \dfrac{\mathrm{e}^{2x} - 1}{x} = \dfrac{\mathrm{e}^{X} - 1}{X/2} = 2\,\dfrac{\mathrm{e}^{X} - 1}{X} \xrightarrow[X \to 0]{} 2 \times 1 = 2. \]
\[ \lim_{x \to 0} \dfrac{\mathrm{e}^{2x} - 1}{x} = 2. \]
\end{enumerate}

\textbf{Partie C.}
\begin{enumerate}
\item Une primitive de $g(x) = \mathrm{e}^{x}$ sur $\mathbb{R}$ est $G(x) = \mathrm{e}^{x}$, car $G'(x) = \mathrm{e}^{x} = g(x)$.
\item La fonction $g$ est continue et positive sur $[0\,;\,1]$, donc l'aire du domaine est :
\[ \mathcal{A} = \int_{0}^{1} \mathrm{e}^{x}\,\mathrm{d}x = \Big[\mathrm{e}^{x}\Big]_{0}^{1} = \mathrm{e}^{1} - \mathrm{e}^{0} = \mathrm{e} - 1 \ \text{u.a.} \]
Valeur approchée : $\mathcal{A} \approx \num{2,718} - 1 = \num{1,718} \approx \num{1,72}$ u.a.\\
L'unité graphique étant \SI{2}{cm} sur chaque axe, une unité d'aire vaut $2 \times 2 = \SI{4}{cm^{2}}$. Donc :
\[ \mathcal{A} = 4(\mathrm{e} - 1)\ \si{cm^{2}} \approx 4 \times \num{1,718} \approx \num{6,87}\ \si{cm^{2}}. \]
\end{enumerate}
\begin{popfigure}
\begin{tikzpicture}
\begin{axis}[width=0.7\linewidth, height=5.5cm, axis lines=middle, xlabel={$x$}, ylabel={$y$}, xmin=-0.3, xmax=1.6, ymin=0, ymax=3.4, xtick={1}, ytick={1,2,3}]
\addplot[popblue!20, draw=none, domain=0:1, samples=50] {exp(x)} \closedcycle;
\addplot[popcurve, domain=-0.3:1.4, samples=100] {exp(x)};
\node[popblue] at (axis cs:1.35,2.6) {\small $(\Gamma)$};
\node[popdark] at (axis cs:0.5,0.8) {\small $\mathcal{A}$};
\draw[popmuted, dashed] (axis cs:1,0) -- (axis cs:1,2.718);
\end{axis}
\end{tikzpicture}
\legende{Aire sous la courbe de $\exp$ entre $x=0$ et $x=1$}
\end{popfigure}
\textbf{Points de vigilance :} en Partie A, ne pas oublier que le système en $(x,y)$ produit \emph{deux} couples solutions (symétrie) et vérifier les conditions $x>0$, $y>0$ ; en B.1, la forme « $+\infty - \infty$ » est indéterminée : il faut factoriser par $x$ et invoquer les croissances comparées ; en C.2, l'intégrale donne des \emph{unités d'aire} : la conversion en $\si{cm^{2}}$ (facteur $4$, pas $2$) est souvent oubliée au Bac.
\end{corrige}

\begin{corrige}[Exercice 11 — Refroidissement d'un moteur]
\textbf{Barème indicatif (sur 10 points) :} question 1 : 1 pt ; question 2 : 4 pts ; question 3 : 3 pts ; question 4 : 2 pts.

\begin{enumerate}
\item $T(0) = 25 + 75\,\mathrm{e}^{0} = 25 + 75 = 100$. Le moteur est à $\mathbf{100\,^{\circ}C}$ à l'instant de l'arrêt.
\item \begin{enumerate}
\item $T$ est dérivable sur $[0\,;+\infty[$ et, avec $u(t) = -0{,}08t$, $u'(t) = -0{,}08$ :
\[ T'(t) = 75 \times (-0{,}08)\,\mathrm{e}^{-0{,}08t} = -6\,\mathrm{e}^{-0{,}08t}. \]
Comme $\mathrm{e}^{-0{,}08t} > 0$ pour tout $t$, on a $T'(t) < 0$ sur $[0\,;+\infty[$ : $T$ est \textbf{strictement décroissante} sur $[0\,;+\infty[$.
\item Par composition, $\displaystyle\lim_{t \to +\infty}(-0{,}08t) = -\infty$ et $\displaystyle\lim_{X \to -\infty} \mathrm{e}^{X} = 0$, donc $\displaystyle\lim_{t \to +\infty} \mathrm{e}^{-0{,}08t} = 0$, puis :
\[ \lim_{t \to +\infty} T(t) = 25. \]
\textbf{Interprétation :} à long terme, la température du moteur se rapproche de la température ambiante $25\,^{\circ}\mathrm{C}$, sans jamais descendre en dessous (car $75\,\mathrm{e}^{-0{,}08t} > 0$). La droite $y = 25$ est asymptote horizontale à la courbe de $T$.
\item Tableau de variations :
\begin{center}
\begin{tabular}{|c|ccc|}
\hline
$t$ & $0$ & & $+\infty$ \\
\hline
$T'(t)$ & & $-$ & \\
\hline
$T$ & $100$ & $\searrow$ & $25$ \\
\hline
\end{tabular}
\end{center}
\end{enumerate}
\item \begin{enumerate}
\item Pour $t \geq 0$ :
\[ T(t) < 40 \iff 25 + 75\,\mathrm{e}^{-0{,}08t} < 40 \iff 75\,\mathrm{e}^{-0{,}08t} < 15 \iff \mathrm{e}^{-0{,}08t} < \dfrac{15}{75} = \dfrac{1}{5}. \]
\item La fonction $\ln$ étant strictement croissante sur $]0\,;+\infty[$ :
\[ \mathrm{e}^{-0{,}08t} < \dfrac{1}{5} \iff -0{,}08t < \ln\dfrac{1}{5} = -\ln 5 \iff t > \dfrac{\ln 5}{0{,}08}, \]
le sens de l'inégalité changeant lors de la division par $-0{,}08 < 0$. Numériquement :
\[ t > \dfrac{\num{1,609}}{0{,}08} \approx \num{20,11}. \]
On peut toucher le moteur sans risque à partir de la $\mathbf{21^{e}}$ \textbf{minute} (arrondi à la minute supérieure).
\end{enumerate}
\item $T(30) = 25 + 75\,\mathrm{e}^{-0{,}08 \times 30} = 25 + 75\,\mathrm{e}^{-2{,}4} \approx 25 + 75 \times \num{0,0907} \approx 25 + \num{6,80}$ :
\[ T(30) \approx \mathbf{31{,}8\,^{\circ}C}. \]
Cohérent avec la question 3 : au bout de 30 minutes, le moteur est déjà passé sous les $40\,^{\circ}\mathrm{C}$.
\end{enumerate}
\textbf{Points de vigilance :} en 2.a, ne pas oublier le facteur $-0{,}08$ issu de la dérivation de la composée ; en 3.b, le \emph{changement de sens} de l'inégalité lors de la division par le nombre négatif $-0{,}08$ est le piège classique ; l'arrondi demandé est « à la minute supérieure » (sécurité), donc $21$ et non $20$ ; en 4, conserver la valeur exacte $\mathrm{e}^{-2{,}4}$ jusqu'au calcul final avant d'arrondir.
\end{corrige}

\newpage

\coursec{Fiche bilan}

\coursec{Fonction exponentielle népérienne}

\begin{popfigure}
\begin{center}
\begin{tikzpicture}[mindmap, grow cyclic, every node/.style={concept, text=white, font=\small\bfseries},
  root concept/.append style={concept color=popdark, font=\large\bfseries},
  level 1/.append style={level distance=3.4cm, sibling angle=60},
  level 2/.append style={level distance=2.2cm, sibling angle=45, font=\scriptsize},
  concept connection/.append style={color=popmuted}]
\node[root concept]{Fonction exponentielle $e^x$}
  child[concept color=popblue]{ node {Définition}
    child { node[concept color=popblue!70] {réciproque de $\ln$} }
    child { node[concept color=popblue!70] {$\ln x = y \Leftrightarrow x = e^y$} }
  }
  child[concept color=poporange]{ node {Propriétés algébriques}
    child { node[concept color=poporange!70] {$e^{a+b}=e^a e^b$} }
    child { node[concept color=poporange!70] {$e^{-a}=\dfrac{1}{e^a}$} }
    child { node[concept color=poporange!70] {$(e^a)^n = e^{na}$} }
  }
  child[concept color=popgreen]{ node {Équations, inéquations}
    child { node[concept color=popgreen!70] {$e^u = e^v \Leftrightarrow u=v$} }
    child { node[concept color=popgreen!70] {$e^u \leq e^v \Leftrightarrow u \leq v$} }
    child { node[concept color=popgreen!70] {$e^x > 0$ toujours} }
  }
  child[concept color=poppurple]{ node {Limites}
    child { node[concept color=poppurple!70] {$e^x \to +\infty$ en $+\infty$} }
    child { node[concept color=poppurple!70] {$e^x \to 0$ en $-\infty$} }
    child { node[concept color=poppurple!70] {$\dfrac{e^x}{x} \to +\infty$} }
  }
  child[concept color=poppink]{ node {Dérivation, variations}
    child { node[concept color=poppink!70] {$(e^u)' = u' e^u$} }
    child { node[concept color=poppink!70] {exp strictement croissante} }
    child { node[concept color=poppink!70] {tangente $y = x+1$ en $0$} }
  }
  child[concept color=popgold]{ node {Primitives, modélisation}
    child { node[concept color=popgold!70] {$\int u' e^u = e^u$} }
    child { node[concept color=popgold!70] {croissance, décroissance} }
  };
\end{tikzpicture}
\end{center}
\legende{Carte mentale de la leçon : la fonction exponentielle népérienne.}
\end{popfigure}

\coursec{Définition et premières propriétés de la fonction exponentielle}

\courssub{Rappel : la fonction logarithme népérien}
La fonction $\ln$ est une bijection strictement croissante de $]0\,;\,+\infty[$ sur $\mathbb{R}$. Elle est continue, dérivable et $\ln'(x) = \frac{1}{x}$. Son ensemble de définition est $]0\,;\,+\infty[$ et son ensemble image est $\mathbb{R}$.

\courssub{Définition de la fonction exponentielle}
\begin{definition}[Fonction exponentielle népérienne]
La fonction \cle{exponentielle népérienne}, notée $\exp$ ou $x \mapsto e^x$, est la \textbf{bijection réciproque} de la fonction logarithme népérien $\ln$.
\end{definition}

\begin{propriete}[Conséquences immédiates de la définition]
\begin{itemize}
  \item $\exp$ est définie sur $\mathbb{R}$ (ensemble de départ de $\ln^{-1}$) et à valeurs dans $]0\,;\,+\infty[$ (ensemble de définition de $\ln$).
  \item $\exp$ est strictement croissante, continue et dérivable sur $\mathbb{R}$.
  \item \textbf{Équivalences fondamentales} : pour tout $x > 0$ et tout réel $y$,
        \[ \ln x = y \iff x = e^y. \]
  \item \textbf{Relations de composition} :
        \begin{align*}
          \forall x \in \mathbb{R},\quad & \ln(e^x) = x, \\
          \forall x > 0,\quad & e^{\ln x} = x.
        \end{align*}
  \item \textbf{Valeurs remarquables} : $e^0 = 1$ (car $\ln 1 = 0$) et $e^1 = e \approx 2,71828\ldots$
  \item \textbf{Positivité} : $\forall x \in \mathbb{R},\ e^x > 0$.
\end{itemize}
\end{propriete}

\begin{exemplebox}[Calculs directs]
\begin{enumerate}
  \item $e^{\ln 5} = 5$.
  \item $\ln(e^{-2}) = -2$.
  \item $e^{2\ln 3} = e^{\ln(3^2)} = e^{\ln 9} = 9$.
  \item $\ln(e^{x^2}) = x^2$ pour tout $x \in \mathbb{R}$.
\end{enumerate}
\end{exemplebox}

\begin{attention}[Pièges classiques]
\begin{itemize}
  \item Ne jamais écrire $e^{\ln x} = x$ sans préciser $x > 0$ (condition d'existence de $\ln x$).
  \item Ne pas confondre $e^{-x} = \frac{1}{e^x}$ (toujours positif) et $-e^x$ (toujours négatif).
  \item $e^x$ ne s'annule \textbf{jamais} : l'équation $e^x = 0$ n'a \textbf{aucune solution}.
\end{itemize}
\end{attention}

\coursec{Propriétés algébriques de l'exponentielle}

\courssub{Formules fondamentales}
\begin{propriete}[Propriétés algébriques]
Pour tous réels $a, b$ et tout entier relatif $n$ :
\begin{align*}
  e^{a+b} &= e^a \times e^b, \\
  e^{a-b} &= \frac{e^a}{e^b}, \\
  e^{-a} &= \frac{1}{e^a}, \\
  (e^a)^n &= e^{na}, \\
  e^{\frac{a}{2}} &= \sqrt{e^a} \quad (a \in \mathbb{R}).
\end{align*}
\end{propriete}

\begin{experience}[Démonstration guidée de $e^{a+b}=e^a e^b$]
On utilise la bijection $\ln$ et ses propriétés.
\begin{enumerate}
  \item Soient $a, b \in \mathbb{R}$. Posons $A = e^a$ et $B = e^b$. Alors $A>0, B>0$ et $\ln A = a$, $\ln B = b$.
  \item Calculer $\ln(AB)$ en fonction de $a$ et $b$.
  \item En déduire $AB = e^{a+b}$.
  \item Conclure.
\end{enumerate}
\tcblower
\textbf{Solution :}
\begin{enumerate}
  \item $\ln(AB) = \ln A + \ln B = a + b$.
  \item Comme $\ln$ est bijective, $AB$ est l'unique réel dont le logarithme est $a+b$, c'est-à-dire $e^{a+b}$.
  \item Donc $e^a e^b = e^{a+b}$.
\end{enumerate}
\end{experience}

\courssub{Simplification d'expressions mêlant $\ln$ et $\exp$}
\begin{methode}[Simplifier une expression contenant $e^x$ et $\ln x$]
\begin{enumerate}
  \item Repérer les compositions $e^{\ln(\cdot)}$ ou $\ln(e^{\cdot})$ et les simplifier ($e^{\ln u}=u$ si $u>0$, $\ln(e^v)=v$).
  \item Utiliser $e^{a+b}=e^a e^b$ et $e^{a-b}=e^a/e^b$ pour séparer ou regrouper les termes.
  \item Utiliser $e^{n\ln k} = k^n$ (pour $k>0$) et $e^{\ln k} = k$.
  \item Vérifier les conditions de validité (arguments de $\ln$ strictement positifs).
\end{enumerate}
\end{methode}

\begin{exoresolu}[Simplification]
Simplifier $A = \dfrac{e^{2\ln 3} \times e^{-\ln 2}}{e^{\ln 6}}$.
\tcblower
$e^{2\ln 3} = e^{\ln 9} = 9$ ; $e^{-\ln 2} = \frac{1}{e^{\ln 2}} = \frac{1}{2}$ ; $e^{\ln 6} = 6$.
Donc $A = \dfrac{9 \times \frac{1}{2}}{6} = \dfrac{9}{12} = \dfrac{3}{4}$.
\end{exoresolu}

\begin{exercice}[Entraînement]
Simplifier les expressions suivantes :
\begin{enumerate}
  \item $B = e^{3\ln 2 - \ln 4}$
  \item $C = \ln\left(e^{x^2} \times e^{2x+1}\right)$
  \item $D = \dfrac{e^{\ln(x^2+1)}}{e^{\ln(x+1)}}$ (préciser l'ensemble de définition)
\end{enumerate}
\end{exercice}

\begin{corrige}[Exercice Entraînement]
\begin{enumerate}
  \item $B = e^{\ln 8 - \ln 4} = e^{\ln 2} = 2$.
  \item $C = \ln(e^{x^2+2x+1}) = x^2+2x+1 = (x+1)^2$.
  \item $D = \dfrac{x^2+1}{x+1}$ pour $x \neq -1$ (condition $x+1>0 \Rightarrow x>-1$ pour le dénominateur initial $e^{\ln(x+1)}$, mais $e^{\ln u}=u$ impose $u>0$). Donc $D = \frac{x^2+1}{x+1}$ sur $]-1;+\infty[$.
\end{enumerate}
\end{corrige}

\coursec{Équations, inéquations et systèmes d'inconnue $e^x$}

\courssub{Principe de résolution}
\begin{propriete}[Monotonie et équivalences]
La fonction $\exp$ est \textbf{strictement croissante} sur $\mathbb{R}$. Par conséquent, pour tous réels $u, v$ :
\begin{align*}
  e^u = e^v &\iff u = v, \\
  e^u < e^v &\iff u < v, \\
  e^u \leq e^v &\iff u \leq v.
\end{align*}
De plus, pour tout $k > 0$ : $e^x = k \iff x = \ln k$.
\end{propriete}

\courssub{Équations faisant intervenir $e^x$}
\begin{methode}[Résolution d'une équation en $e^x$]
\begin{enumerate}
  \item Mettre l'équation sous la forme $P(e^x) = 0$ où $P$ est un polynôme.
  \item Effectuer le changement de variable $X = e^x$ (avec la condition \textbf{obligatoire} $X > 0$).
  \item Résoudre l'équation polynomiale $P(X) = 0$ dans $\mathbb{R}$.
  \item \textbf{Rejeter} les solutions $X \leq 0$.
  \item Pour chaque solution $X_k > 0$ retenue, en déduire $x = \ln X_k$.
\end{enumerate}
\end{methode}

\begin{exoresolu}[Équation du second degré en $e^x$]
Résoudre dans $\mathbb{R}$ : $e^{2x} - 5e^x + 6 = 0$.
\tcblower
Posons $X = e^x$ ($X > 0$). L'équation devient $X^2 - 5X + 6 = 0$.
$\Delta = 25 - 24 = 1$. Racines : $X_1 = \frac{5-1}{2} = 2$, $X_2 = \frac{5+1}{2} = 3$.
Les deux sont strictement positives, on les garde.
$x = \ln 2$ ou $x = \ln 3$.
$S = \{\ln 2\ ;\ \ln 3\}$.
\end{exoresolu}

\begin{exoresolu}[Équation avec solution rejetée]
Résoudre : $e^{2x} + e^x - 2 = 0$.
\tcblower
$X = e^x > 0$. $X^2 + X - 2 = 0 \iff (X+2)(X-1)=0$.
Solutions : $X = -2$ (rejetée car $<0$) ou $X = 1$ (acceptée).
$x = \ln 1 = 0$.
$S = \{0\}$.
\end{exoresolu}

\courssub{Inéquations faisant intervenir $e^x$}
\begin{methode}[Résolution d'une inéquation en $e^x$]
\begin{enumerate}
  \item Isoler le terme exponentiel (ou poser $X=e^x$ si forme polynomiale).
  \item Utiliser la monotonie : $e^u \leq e^v \iff u \leq v$ (le sens de l'inégalité est \textbf{conservé}).
  \item Si utilisation de $\ln$ : vérifier la condition d'existence (argument $>0$).
\end{enumerate}
\end{methode}

\begin{exoresolu}[Inéquation]
Résoudre : $e^{2x-1} \geq e^{x+2}$.
\tcblower
$\exp$ est croissante, donc l'inéquation est équivalente à :
$2x - 1 \geq x + 2 \iff x \geq 3$.
$S = [3\,;\,+\infty[$.
\end{exoresolu}

\begin{exoresolu}[Inéquation quadratique en $e^x$]
Résoudre : $e^{2x} - 3e^x - 4 < 0$.
\tcblower
$X = e^x > 0$. $X^2 - 3X - 4 < 0 \iff (X-4)(X+1) < 0$.
Tableau de signes : solution en $X \in ]-1\,;\,4[$.
Avec la condition $X > 0$ : $X \in ]0\,;\,4[$.
Donc $0 < e^x < 4 \iff e^x < 4 \iff x < \ln 4 = 2\ln 2$.
$S = ]-\infty\,;\,2\ln 2[$.
\end{exoresolu}

\courssub{Systèmes d'équations}
\begin{exemplebox}[Système symétrique]
Résoudre $\begin{cases} e^x + e^y = 5 \\ e^x e^y = 6 \end{cases}$.
$X=e^x>0, Y=e^y>0$. $\begin{cases} X+Y=5 \\ XY=6 \end{cases}$.
$X,Y$ sont les racines de $T^2 - 5T + 6 = 0 \Rightarrow \{2, 3\}$.
Cas 1 : $e^x=2, e^y=3 \Rightarrow (x,y)=(\ln 2, \ln 3)$.
Cas 2 : $e^x=3, e^y=2 \Rightarrow (x,y)=(\ln 3, \ln 2)$.
$S = \{(\ln 2, \ln 3)\ ;\ (\ln 3, \ln 2)\}$.
\end{exemplebox}

\begin{attention}[Erreurs fatales au Bac]
\begin{itemize}
  \item Oublier la condition $X = e^x > 0$ et garder des racines négatives ou nulles.
  \item Écrire $e^x = -2 \Rightarrow x = \ln(-2)$ (impossible).
  \item Inverser le sens de l'inégalité en appliquant $\ln$ (alors que $\ln$ est croissante, elle conserve l'ordre).
  \item Oublier de vérifier la condition d'existence ($>0$) avant d'appliquer $\ln$ sur une inéquation du type $\ln(u(x)) < \ln(v(x))$.
\end{itemize}
\end{attention}

\coursec{Limites de la fonction exponentielle}

\courssub{Limites usuelles}
\begin{propriete}[Limites de $\exp$ aux bornes de $\mathbb{R}$]
\begin{align*}
  \lim_{x \to +\infty} e^x &= +\infty, \\
  \lim_{x \to -\infty} e^x &= 0, \\
  \lim_{x \to +\infty} e^{-x} &= 0, \\
  \lim_{x \to -\infty} e^{-x} &= +\infty.
\end{align*}
La droite d'équation $y = 0$ (axe des abscisses) est une \textbf{asymptote horizontale} à la courbe de $\exp$ en $-\infty$.
\end{propriete}

\courssub{Croissances comparées (limites indispensables)}
\begin{propriete}[Limites comparées]
\begin{align*}
  \lim_{x \to +\infty} \frac{e^x}{x} &= +\infty \quad \text{(l'exponentielle l'emporte sur toute puissance de $x$)}, \\
  \lim_{x \to +\infty} \frac{x^n}{e^x} &= 0 \quad (n \in \mathbb{N}), \\
  \lim_{x \to -\infty} x e^x &= 0, \\
  \lim_{x \to 0} \frac{e^x - 1}{x} &= 1 \quad \text{(taux d'accroissement en 0, dérivée de $\exp$ en 0)}.
\end{align*}
\end{propriete}

\begin{experience}[Démonstration de $\lim_{x \to 0} \frac{e^x-1}{x} = 1$]
C'est la définition même du nombre dérivé de $\exp$ en $0$ : $(\exp)'(0) = \exp(0) = 1$.
On peut aussi utiliser l'équivalence $e^x - 1 \sim_0 x$ issue du développement limité d'ordre 1.
\end{experience}

\courssub{Limites de fonctions composées $e^{u(x)}$}
\begin{methode}[Limite de $e^{u(x)}$]
\begin{enumerate}
  \item Calculer $\ell = \lim u(x)$ (aux bornes de l'intervalle d'étude).
  \item Conclure par composition :
        \begin{itemize}
          \item Si $\ell = +\infty$, alors $e^{u(x)} \to +\infty$.
          \item Si $\ell = -\infty$, alors $e^{u(x)} \to 0$.
          \item Si $\ell \in \mathbb{R}$, alors $e^{u(x)} \to e^\ell$.
        \end{itemize}
  \item Attention aux formes indéterminées (ex: $e^x - e^x$, $x e^x$ en $-\infty$) : factoriser ou utiliser les croissances comparées.
\end{enumerate}
\end{methode}

\begin{exoresolu}[Limites composées]
\begin{enumerate}
  \item $\lim_{x \to +\infty} e^{x^2-3x} = +\infty$ car $x^2-3x \to +\infty$.
  \item $\lim_{x \to -\infty} e^{x^2-3x} = +\infty$ car $x^2-3x \to +\infty$.
  \item $\lim_{x \to 0} \frac{e^{2x}-1}{x} = \lim_{x \to 0} 2\frac{e^{2x}-1}{2x} = 2 \times 1 = 2$.
  \item $\lim_{x \to -\infty} (x^3 e^x) = 0$ (croissance comparée : $x^3$ polynomial vs $e^{-x}$ exponentiel en $+\infty$).
\end{enumerate}
\end{exoresolu}

\coursec{Dérivation, variations et courbe représentative}

\courssub{Dérivée de la fonction exponentielle}
\begin{propriete}[Dérivée de $\exp$ et de $e^{u(x)}$]
\begin{itemize}
  \item La fonction $\exp$ est dérivable sur $\mathbb{R}$ et $\forall x \in \mathbb{R},\ (\exp)'(x) = \exp(x) = e^x$.
  \item Si $u$ est dérivable sur $I$, alors $x \mapsto e^{u(x)}$ est dérivable sur $I$ et
        \[ (e^{u(x)})' = u'(x) e^{u(x)}. \]
\end{itemize}
\end{propriete}

\begin{experience}[Démonstration de $(\exp)' = \exp$ par la bijection réciproque]
On sait que $\ln$ est dérivable sur $]0;+\infty[$ avec $\ln'(x) = 1/x$.
$\exp$ est la bijection réciproque de $\ln$.
Pour tout $y \in \mathbb{R}$, posons $x = e^y > 0$. Alors $\ln x = y$.
D'après le théorème de dérivation de la bijection réciproque :
$(\exp)'(y) = \frac{1}{\ln'(x)} = \frac{1}{1/x} = x = e^y$.
Donc $(\exp)'(y) = \exp(y)$.
\end{experience}

\courssub{Variations de $x \mapsto e^{u(x)}$}
\begin{propriete}
Soit $u$ dérivable sur un intervalle $I$. La fonction $f(x) = e^{u(x)}$ a pour dérivée $f'(x) = u'(x) e^{u(x)}$.
Comme $e^{u(x)} > 0$ sur $I$, \textbf{le signe de $f'$ est exactement le signe de $u'$}.
\begin{itemize}
  \item $f$ est croissante sur les intervalles où $u$ est croissante.
  \item $f$ est décroissante sur les intervalles où $u$ est décroissante.
  \item Les extrema de $f$ correspondent aux extrema de $u$.
\end{itemize}
\end{propriete}

\begin{attention}[Cas particulier : $f(x) = u(x)e^{v(x)}$]
Si $f(x) = u(x)e^{v(x)}$, alors $f'(x) = \bigl(u'(x) + u(x)v'(x)\bigr)e^{v(x)}$.
Le signe de $f'$ dépend alors de $u'(x) + u(x)v'(x)$, \textbf{pas seulement de $u'$ ou $v'$}.
Ne pas confondre avec $f(x)=e^{u(x)}$.
\end{attention}

\begin{exoresolu}[Étude complète de $f(x) = (x-2)e^x$ sur $\mathbb{R}$]
\begin{enumerate}
  \item \textbf{Domaine :} $D_f = \mathbb{R}$.
  \item \textbf{Limites :}
        $\lim_{x \to -\infty} (x-2)e^x = 0$ (croissance comparée).
        $\lim_{x \to +\infty} (x-2)e^x = +\infty$.
  \item \textbf{Dérivée :} $f'(x) = 1\cdot e^x + (x-2)e^x = (x-1)e^x$.
        $e^x > 0$, donc signe de $f'$ = signe de $x-1$.
  \item \textbf{Tableau de variations :}
        \begin{center}
        \begin{tabular}{|c|ccccc|}
        \hline
        $x$ & $-\infty$ & & $1$ & & $+\infty$ \\ \hline
        $x-1$ & & $-$ & $0$ & $+$ & \\ \hline
        $f'(x)$ & & $-$ & $0$ & $+$ & \\ \hline
        $f(x)$ & $0$ & $\searrow$ & $-e$ & $\nearrow$ & $+\infty$ \\ \hline
        \end{tabular}
        \end{center}
  \item \textbf{Extremum :} Minimum global $f(1) = -e$ en $x=1$.
\end{enumerate}
\end{exoresolu}

\courssub{Courbe représentative de $y = e^x$}
\begin{popfigure}
\begin{center}
\begin{tikzpicture}[scale=0.8, domain=-3:1.5]
\begin{axis}[
    axis lines=middle,
    xlabel=$x$, ylabel=$y$,
    xmin=-3.5, xmax=2, ymin=-0.5, ymax=5,
    xtick={-3,-2,-1,0,1}, ytick={1,2,3,4},
    tick label style={font=\small},
    samples=200,
    width=\linewidth, height=0.6\linewidth,
    clip=false
]
\addplot[popcurve, thick, domain=-3:1.6] {exp(x)};
% Asymptote
\addplot[dashed, popmuted, domain=-3.5:2] {0};
% Tangente en 0 : y = x+1
\addplot[poporange, dashed, domain=-1.5:1.5] {x+1};
% Points remarquables
\addplot[only marks, mark=*, popblue] coordinates {(0,1) (1,2.718)};
\node[above left, popblue, font=\small] at (axis cs:0,1) {$(0;1)$};
\node[above right, popblue, font=\small] at (axis cs:1,2.718) {$(1;e)$};
% Annotation tangente
\node[poporange, font=\small, anchor=south west] at (axis cs:0.5, 1.5) {$y=x+1$};
% Annotation asymptote
\node[popmuted, font=\small, anchor=north east] at (axis cs:2, 0.2) {Asymptote $y=0$};
\end{axis}
\end{tikzpicture}
\end{center}
\legende{Courbe représentative $C_{\exp}$ : passage par $(0;1)$ et $(1;e)$, asymptote $y=0$ en $-\infty$, tangente $y=x+1$ en $0$.}
\end{popfigure}

\begin{aretenir}[Caractéristiques de la courbe $y=e^x$]
\begin{itemize}
  \item Passe par $(0;1)$ et $(1;e)$.
  \item Toujours au-dessus de l'axe des abscisses ($e^x > 0$).
  \item Asymptote horizontale $y=0$ en $-\infty$.
  \item Croissance "explosive" en $+\infty$.
  \item Tangente en $0$ : $y = x + 1$ (car $f'(0)=1$).
  \item Inégalité majeure : $\forall x \in \mathbb{R},\ e^x \geq x + 1$ (égalité seulement en $x=0$).
  \item Courbe convexe ($f''(x)=e^x > 0$).
\end{itemize}
\end{aretenir}

\coursec{Primitives et modélisation par l'exponentielle}

\courssub{Primitives usuelles}
\begin{propriete}[Primitives de la fonction exponentielle]
Soit $I$ un intervalle de $\mathbb{R}$.
\begin{itemize}
  \item Une primitive de $x \mapsto e^x$ sur $\mathbb{R}$ est $x \mapsto e^x + C$.
  \item Si $u$ est dérivable sur $I$, une primitive de $x \mapsto u'(x) e^{u(x)}$ sur $I$ est $x \mapsto e^{u(x)} + C$.
  \item Pour $a \neq 0$, une primitive de $x \mapsto e^{ax+b}$ sur $\mathbb{R}$ est $x \mapsto \frac{1}{a}e^{ax+b} + C$.
\end{itemize}
\end{propriete}

\begin{exemplebox}[Calcul d'intégrales]
\begin{enumerate}
  \item $\displaystyle \int_0^1 e^{2x}\,dx = \left[ \frac{1}{2}e^{2x} \right]_0^1 = \frac{1}{2}(e^2 - 1)$.
  \item $\displaystyle \int xe^{x^2}\,dx$ : on reconnaît $u(x)=x^2$, $u'(x)=2x$. $\int xe^{x^2}dx = \frac{1}{2}\int 2x e^{x^2}dx = \frac{1}{2}e^{x^2} + C$.
  \item $\displaystyle \int \frac{e^x}{1+e^x}\,dx = \ln(1+e^x) + C$ (car dérivée de $1+e^x$ est $e^x$).
\end{enumerate}
\end{exemplebox}

\courssub{Modélisation : croissance et décroissance exponentielles}
\begin{savaistu}[Modélisation réelle : la loi de l'évolution proportionnelle]
De nombreux phénomènes naturels ou économiques suivent une loi où \textbf{la vitesse d'évolution d'une quantité est proportionnelle à la quantité elle-même}.
Mathématiquement : $y'(t) = k \cdot y(t)$ où $k \in \mathbb{R}^*$.
La solution générale est $y(t) = y_0 e^{kt}$ avec $y_0 = y(0)$.
\begin{itemize}
  \item $k > 0$ : \textbf{Croissance exponentielle} (population, capital avec intérêts composés, épidémie au début).
  \item $k < 0$ : \textbf{Décroissance exponentielle} (radioactivité, refroidissement, décharge d'un condensateur, amortissement).
\end{itemize}
\end{savaistu}

\begin{exemplebox}[Contexte camerounais : Démographie et Finance]
\begin{itemize}
  \item \textbf{Croissance démographique :} La population de Yaoundé croît d'environ $3,5\%$ par an. Si $P_0 = 4\,000\,000$ habitants en 2020, $P(t) = 4 \times 10^6 \times e^{0,035 t}$ (modèle continu). En 2030 ($t=10$) : $P(10) \approx 5,7$ millions.
  \item \textbf{Épargne Mobile Money (MTN MoMo / Orange Money) :} Un taux d'intérêt annuel de $5\%$ capitalisé en continu donne $C(t) = C_0 e^{0,05 t}$.
  \item \textbf{Refroidissement (Loi de Newton) :} Un café à $90^\circ\text{C}$ dans une pièce à $25^\circ\text{C}$. L'écart de température $\Delta T = T - 25$ décroît exponentiellement : $\Delta T(t) = 65 e^{-kt}$.
  \item \textbf{Radioactivité :} Demi-vie du Carbone 14 $\approx 5730$ ans. $k = -\frac{\ln 2}{5730}$. Datation archéologique.
\end{itemize}
\end{exemplebox}

\begin{methode}[Modéliser par une exponentielle]
\begin{enumerate}
  \item Identifier la quantité $y(t)$ et l'instant initial $t=0$ (valeur $y_0$).
  \item Exprimer la relation "vitesse proportionnelle à la quantité" : $y' = ky$.
  \item Écrire la solution : $y(t) = y_0 e^{kt}$.
  \item Déterminer $k$ à partir d'une donnée (demi-vie, pourcentage d'évolution, seconde valeur connue).
  \item Répondre à la question (valeur à un instant, temps pour atteindre un seuil, etc.).
\end{enumerate}
\end{methode}

\begin{exoresolu}[Décroissance radioactive : Carbone 14]
Un fragment de bois trouvé sur un site archéologique au Cameroun contient $12,5\%$ de son Carbone 14 initial. La demi-vie du C14 est $5730$ ans. Dater le fragment.
\tcblower
$m(t) = m_0 e^{kt}$. Demi-vie : $m(5730) = \frac{m_0}{2} \Rightarrow e^{5730k} = \frac{1}{2} \Rightarrow 5730k = \ln(1/2) = -\ln 2 \Rightarrow k = -\frac{\ln 2}{5730}$.
On a $m(t) = 0,125 m_0 = \frac{m_0}{8}$.
$e^{kt} = \frac{1}{8} = 2^{-3} \Rightarrow kt = -3\ln 2$.
$t = \frac{-3\ln 2}{k} = \frac{-3\ln 2}{-\frac{\ln 2}{5730}} = 3 \times 5730 = 17190$ ans.
Le fragment a environ $17\,190$ ans.
\end{exoresolu}

\begin{exoresolu}[Refroidissement : Thé glacé]
Un thé à $80^\circ\text{C}$ est mis au réfrigérateur à $4^\circ\text{C}$. Au bout de $10$ min, il est à $50^\circ\text{C}$. Quelle température au bout de $30$ min ?
\tcblower
$T(t)$ température. Écart $E(t) = T(t) - 4$. $E(0) = 76$. $E(10) = 46$.
$E(t) = 76 e^{kt}$. $76 e^{10k} = 46 \Rightarrow e^{10k} = \frac{46}{76} = \frac{23}{38}$.
$10k = \ln(23/38) \Rightarrow k = \frac{1}{10}\ln(23/38) \approx -0,05$.
$E(30) = 76 e^{30k} = 76 (e^{10k})^3 = 76 \times \left(\frac{23}{38}\right)^3 = 76 \times \frac{12167}{54872} \approx 16,86$.
$T(30) = 4 + 16,86 \approx 20,9^\circ\text{C}$.
\end{exoresolu}

\begin{exercice}[Entraînement Bac - Modélisation]
La population d'une localité rurale électrifiée par un projet CAMTEL suit le modèle $P(t) = 5000 e^{0,04 t}$ ($t$ en années, $t=0$ en 2025).
\begin{enumerate}
  \item Quelle était la population en 2025 ?
  \item Quelle sera la population en 2035 ?
  \item Au bout de combien d'années la population aura-t-elle doublé ? (Donner une valeur approchée au dixième d'année).
  \item Calculer le taux d'évolution annuel moyen (en $\%$) sur la période 2025-2035.
\end{enumerate}
\end{exercice}

\begin{corrige}[Exercice Modélisation]
\begin{enumerate}
  \item $P(0) = 5000$ habitants.
  \item $t=10$ : $P(10) = 5000 e^{0,4} \approx 5000 \times 1,4918 \approx 7459$ habitants.
  \item Doublement : $10000 = 5000 e^{0,04 t} \iff e^{0,04 t} = 2 \iff 0,04 t = \ln 2 \iff t = \frac{\ln 2}{0,04} \approx \frac{0,693}{0,04} \approx 17,3$ ans.
  \item Taux annuel moyen : $\frac{P(10)-P(0)}{10 \times P(0)} \times 100 = \frac{e^{0,4}-1}{10} \times 100 \approx \frac{0,4918}{10} \times 100 \approx 4,9\%$. (Proche du $4\%$ instantané car continu).
\end{enumerate}
\end{corrige}

\coursec{Bilan de synthèse et Checklist Bac}

\begin{aretenir}[L'essentiel à maîtriser pour le Baccalauréat]
\begin{center}
\begin{tabularx}{\linewidth}{|l|X|}
\hline
\rowcolor{popblueL} \textbf{Thème} & \textbf{Points clés / Formules / Méthodes} \\
\hline
\textbf{Définition} & $\exp = \ln^{-1}$ bijection $\mathbb{R} \to ]0;+\infty[$. $\ln x = y \iff x = e^y$. $\ln(e^x)=x$, $e^{\ln x}=x\ (x>0)$. $e^0=1$, $e^1=e$. \\
\hline
\textbf{Algèbre} & $e^{a+b}=e^a e^b$, $e^{a-b}=e^a/e^b$, $e^{-a}=1/e^a$, $(e^a)^n=e^{na}$, $e^{\ln k}=k\ (k>0)$. \\
\hline
\textbf{Équations} & Poser $X=e^x>0$. Résoudre $P(X)=0$. \textbf{Rejeter $X \leq 0$}. $x=\ln X$. \\
\hline
\textbf{Inéquations} & $\exp$ croissante : $e^u \leq e^v \iff u \leq v$ (sens conservé). Si $\ln$ : vérifier $>0$. \\
\hline
\textbf{Limites} & $+\infty \to +\infty$, $-\infty \to 0$ (asymptote $y=0$). $\frac{e^x}{x}\to+\infty$, $x e^x\to 0\ (-\infty)$, $\frac{e^x-1}{x}\to 1\ (0)$. \\
\hline
\textbf{Dérivée} & $(e^x)'=e^x$, $(e^u)'=u'e^u$. Variations de $e^u$ : signe de $u'$. \\
\hline
\textbf{Courbe} & Points $(0;1)$, $(1;e)$. Asymptote $y=0$ ($-\infty$). Tangente $0$ : $y=x+1$. Convexe. $e^x \geq x+1$. \\
\hline
\textbf{Primitives} & $\int e^x = e^x+C$. $\int u'e^u = e^u+C$. $\int e^{ax+b} = \frac{1}{a}e^{ax+b}+C\ (a\neq 0)$. \\
\hline
\textbf{Modélisation} & $y'=ky \iff y=y_0 e^{kt}$. $k>0$ croissance, $k<0$ décroissance. Demi-vie $T_{1/2} = \frac{\ln 2}{|k|}$. \\
\hline
\end{tabularx}
\end{center}
\end{aretenir}

\begin{aretenir}[Checklist finale : Es-tu prêt ?]
Avant l'épreuve, vérifie que tu sais cocher toutes les cases :
\begin{itemize}
  \item[$\square$] Définir $\exp$ comme réciproque de $\ln$ et en déduire ensemble de définition, image, monotonie, parité (aucune).
  \item[$\square$] Utiliser $\ln x = y \iff x = e^y$ sans oublier $x>0$.
  \item[$\square$] Appliquer $e^{a+b}=e^a e^b$, $e^{-a}=1/e^a$, $(e^a)^n=e^{na}$ pour simplifier.
  \item[$\square$] Simplifier $e^{\ln u}=u$ ($u>0$) et $\ln(e^v)=v$.
  \item[$\square$] Résoudre $e^{2x}+ae^x+b=0$ par $X=e^x>0$ et \textbf{rejeter $X \leq 0$}.
  \item[$\square$] Résoudre inéquations en conservant le sens (croissance de $\exp$ et de $\ln$).
  \item[$\square$] Citer et utiliser les 5 limites usuelles ($+\infty, -\infty, e^x/x, xe^x, (e^x-1)/x$).
  \item[$\square$] Dériver $e^{u(x)}$ : $(e^u)' = u'e^u$ (ne pas oublier $u'$ !).
  \item[$\square$] Étudier variations de $f(x)=e^{u(x)}$ : signe de $f'$ = signe de $u'$.
  \item[$\square$] Tracer $y=e^x$ : points $(0;1)$, $(1;e)$, asymptote $y=0$, tangente $y=x+1$, convexité.
  \item[$\square$] Calculer primitives : $\int u'e^u = e^u$, $\int e^{ax+b} = \frac{1}{a}e^{ax+b}$.
  \item[$\square$] Modéliser : $y'=ky \to y=y_0 e^{kt}$. Interpréter $k$, calculer demi-vie ou temps de doublement.
\end{itemize}
\end{aretenir}

\findecours

\end{document}
